\documentclass[12pt]{article}

% Page layout
\usepackage[margin=1in]{geometry}
\usepackage{setspace}
\doublespacing  % or \onehalfspacing

% Math
\usepackage{amsmath}
\usepackage{amssymb}
\usepackage{amsthm}
\usepackage{mathtools}

% Tables
\usepackage{booktabs}
\usepackage{multirow}
\usepackage{threeparttable}
\usepackage{tabularx}

% General
\usepackage{graphicx}
\usepackage{hyperref}
\usepackage{natbib}

\usepackage{fancyhdr}
\fancyhead{}
\renewcommand{\headrulewidth}{0pt}
\fancyfoot{}

\usepackage{geometry}
\geometry{margin=1in}

\usepackage{enumerate}

\usepackage{rotating}

% Theorem environments
\newtheorem{theorem}{Theorem}
\newtheorem{lemma}{Lemma}
\newtheorem{corollary}{Corollary}
\newtheorem{conjecture}{Conjecture}
\newtheorem{remark}{Remark}
\newtheorem{definition}{Definition}
\newtheorem{assumption}{Assumption}
\newtheorem{proposition}{Proposition}


% Custom commands for asymptotic theory section
\newcommand{\norm}[1]{\left\| #1 \right\|}
\newcommand{\R}{\mathbb{R}}
\newcommand{\E}{\mathbb{E}}
\newcommand{\Lp}{L^p(\Omega)}
\newcommand{\Lone}{L^1(\Omega)}
\newcommand{\Wonep}{W^{1,p}(\Omega)}
\newcommand{\Wtwop}{W^{2,p}(\Omega)}
\newcommand{\BV}{BV(\Omega)}
\newcommand{\cM}{\mathcal{M}(\Omega)^d}
\newcommand{\fn}{\hat{f}_n}
\newcommand{\xrightarrowp}{\xrightarrow{p}}
\newcommand{\wstar}{\xrightarrow{w*}}
\newcommand{\Prob}{\mathbb{P}}
\newcommand{\abs}[1]{\left| #1 \right|}
\newcommand{\Dh}{\Delta_{j,h}}
\newcommand{\Ptwo}[1]{\norm{#1}_{P,2}}

\begin{document}
	
	\title{Model-Agnostic Interpretation of Machine Learning\\via Average Marginal Effects}
	\author{Jason Parker\thanks{Email: \texttt{jparker588@gmail.com}.}}
	\date{August 2026}
	
	\maketitle
	\thispagestyle{fancy}
	
	\begin{abstract}
		\noindent We propose a framework for interpreting any supervised machine learning
		model through average marginal effects, the standard inferential tool in the
		econometrics of nonlinear models. The estimand is the window average marginal
		effect --- the average centered finite difference of the regression function
		at a fixed, interpretable step --- which is well defined for learners with no
		differentiability anywhere and recovers the classical average derivative as
		the window shrinks. Because the functional is linear in the regression
		function, it admits an exact Riesz representation obtained by a change of
		variables, with no boundary terms and no density derivatives, and the
		resulting orthogonal score has estimation bias exactly equal to a product of
		two $L^2$ errors. A cross-fitted debiased estimator is therefore root-$n$
		consistent and asymptotically normal whenever the learner and a
		Riesz-regression estimate of the representer converge faster than $n^{-1/4}$
		in $L^2$ --- rather than at the parametric rate in a function-space norm ---
		attains the semiparametric efficiency bound, and returns a coefficient table
		directly analogous to parametric regression output without resampling; the
		classical generalized linear model result is nested as a special case, and
		binary contrasts arise from the same representation with inverse propensity
		weights. The score is doubly robust: when the representer is known or
		parametrically estimable --- designed experiments, randomized treatments,
		simulation designs --- inference remains valid with no consistency
		requirement on the learner, covering the fixed-complexity settings in which
		off-the-shelf learners converge to pseudo-true limits. A pairing identity
		characterizes, apparently for the first time,
		what the marginal effect of a tree-based model measures: the jump measure of
		the fitted surface integrated against a locally averaged covariate density.
		Gradient consistency results for neural networks in Sobolev spaces and for
		boosted trees and Mondrian forests in the space of functions of bounded
		variation delineate when uncorrected plug-in effects are already consistent.
		Experiments on three datasets and simulations illustrate the framework and
		display the regularization-bias mechanism that the orthogonal correction
		removes.
	\end{abstract}
	
	\section{Introduction}
	
	The rise of flexible machine learning methods has created a persistent tension in
	empirical research. Models such as random forests \citep{breiman2001rf}, gradient
	boosted trees \citep{chen2016xgboost}, and deep neural networks \citep{hastie2009esl}
	routinely achieve predictive accuracy that parametric models cannot match. Yet the
	standard tools of statistical inference --- coefficient estimates, standard errors,
	hypothesis tests --- apply naturally only to parametric models whose functional form
	is specified in advance. \citet{breiman2001cultures} framed this as a conflict between
	two cultures of statistical modeling: one oriented toward interpretable
	data-generating models, the other toward accurate algorithmic prediction. The
	conventional resolution has been to treat accuracy and interpretability as opposing
	ends of a spectrum, accepting reduced interpretability as the price of improved
	prediction. This paper argues that the tradeoff is not fundamental. It is an artifact
	of requiring parametric assumptions to make interpretation rigorous, and it disappears
	once that requirement is dropped.
	
	The machine learning literature has produced a rich set of post-hoc interpretability
	methods that attempt to explain black-box model predictions without parametric
	assumptions. Partial dependence plots \citep{friedman2001pdp} visualize the marginal
	relationship between a feature and the predicted outcome by averaging over the
	empirical distribution of other features. LIME \citep{ribeiro2016lime} fits a local
	linear approximation to the prediction function in the neighborhood of individual
	observations. SHAP \citep{lundberg2017shap} assigns each feature a contribution to
	each prediction based on Shapley values from cooperative game theory, providing a
	unified framework that nests several earlier methods. These tools are widely used in
	practice and are surveyed comprehensively in \citet{molnar2025} and
	\citet{doshivelez2017interpretability}. However, they share a fundamental limitation:
	they produce point estimates with no associated sampling distributions. A SHAP value,
	a PDP slope, and a LIME coefficient are descriptive quantities --- there is no
	standard error, no confidence interval, and no basis for a hypothesis test. A
	researcher cannot determine whether a SHAP value is statistically distinguishable
	from zero, or whether two features have significantly different effects, or whether
	the effect of a variable differs across model specifications. For applied researchers
	who need to draw defensible conclusions from their models, this is a serious
	constraint.
	
	Two prior contributions are closest in spirit to the present paper.
	\citet{scholbeck2024fme} propose forward marginal effects for nonlinear prediction
	functions as a model-agnostic interpretability method, establishing a general
	framework for computing feature effects for any supervised learner. Their framework
	is descriptive rather than inferential: it characterizes how feature effects vary
	across the covariate space but does not provide standard errors, confidence intervals,
	or hypothesis tests, and it explicitly argues against summarizing feature effects in
	a single metric such as the average marginal effect. The present paper takes the
	opposite position: we argue that the average marginal effect is precisely the right
	summary quantity, that it has a rigorous population-level interpretation inherited
	from the econometric literature, and that valid, semiparametrically efficient
	inference is achievable via a debiased, cross-fitted estimator with an explicit
	influence function. \cite{lowe2024fmeffects}
	provide the \texttt{fmeffects} R package implementing this descriptive framework,
	demonstrating the practical demand for marginal-effect-style interpretation of
	machine learning models, but without the theoretical justification --- no population estimand, no
	efficiency theory, no characterization of when the marginal effects are
	well-defined --- that the present paper provides through the Riesz representation
	of the effect functional and semiparametric efficiency theory.
	
	The econometric literature has solved exactly this inferential problem for parametric
	models through average marginal effects \citep{greene2012, wooldridge2010}. For a
	generalized linear model with inverse link function $\mu$ and linear predictor
	$\eta = \mathbf{x}^\top \boldsymbol{\beta}$, the average marginal effect of variable
	$j$ is $\mathrm{AME}_j = \beta_j \cdot \mathbb{E}[\mu'(\eta)]$ --- the regression
	coefficient scaled by the average derivative of the inverse link, averaged over the
	empirical distribution of the linear predictor. For the identity link this reduces
	to the OLS coefficient. For the logistic link it averages the logistic density over
	the sample. The result is a quantity measured in the natural units of the outcome,
	interpretable as the average change in the predicted outcome associated with a
	one-unit change in $x_j$, with standard errors obtained by the delta method or
	bootstrap. This is the output of \texttt{margins} in R and Stata \citep{leeper2021},
	and it is the standard form in which marginal effects are reported in applied
	economics, political science, sociology, and public health. The limitation is that
	it has been restricted to parametric models: to compute $\partial f / \partial x_j$
	analytically, one must know the functional form of $f$.
	
	A separate literature, largely disjoint from both the interpretability and the
	applied marginal effects traditions, has studied the population version of this
	quantity directly. The average derivative $\E[\partial_j f(X)]$ was introduced as
	an estimand in its own right by \citet{hardle1989} and \citet{powell1989}, its
	semiparametric efficiency bounds were derived by \citet{newey1993}, and the modern
	debiased machine learning program \citep{chernozhukov2018dml,
	chernozhukov2022riesz} shows how to estimate linear functionals of a regression
	function at the parametric rate using flexible learners as nuisance estimators.
	The present paper is deliberately positioned at the junction of these
	literatures, and its delta relative to each is stated in one sentence: the
	debiasing literature tells one how to estimate $\E[\partial_j f(X)]$ efficiently
	when the derivative exists and the learner is treated as a black-box nuisance;
	this paper establishes what the estimand \emph{is} when the learner is not
	differentiable --- the window effect, with an exact finite-difference Riesz
	representation requiring no boundary conditions and no density derivatives --- and
	shows that a single orthogonal score then covers smooth and piecewise constant
	learners, continuous and discrete regressors, within one construction. Conversely,
	relative to the interpretability literature, the contribution is to replace
	descriptive scores with an estimand--estimator--efficiency-bound triple of the
	classical kind.
	
	This paper removes that restriction. We propose a framework for computing average
	marginal effects for any supervised machine learning model, requiring only that the
	prediction function $\hat{f}$ can be evaluated at a point. The estimand is the
	\emph{window average marginal effect}: the expectation of a centered finite
	difference of the regression function at a fixed step $h$ scaled to the spread of
	each feature, trimmed away from the boundary of the support. Treating the finite
	difference as the estimand rather than as a numerical approximation to a
	derivative is what makes genuine model-agnosticism possible --- the window effect
	is well defined for tree-based models, for which classical derivatives do not
	exist anywhere --- while the classical average derivative is recovered as the
	$h \to 0$ limit whenever the regression function has one weak derivative, so
	nothing is given up relative to the econometric tradition. Estimation combines
	the plug-in average of finite differences with an orthogonal correction term
	built from an explicit Riesz representer, computed by cross-fitting; standard
	errors come from the estimated influence function in a single pass, with no
	resampling required, and a multiplier bootstrap is available at no additional
	model fits when a resampling distribution is wanted. The result is a coefficient
	table --- estimate, standard error, and significance level --- directly analogous
	to the output of a parametric regression but with no restrictions on the
	functional form of $\hat{f}$. The framework nests the classical GLM result as a
	special case: when $\hat{f}$ is a logistic regression, the estimator recovers the
	classical AME automatically. When $\hat{f}$ is XGBoost or a random forest or a
	neural network, the same formula applies and produces output on the same scale
	and with the same interpretation, and for binary regressors the construction
	specializes to augmented inverse probability weighting with the regressor in the
	role of the treatment.
	
	The theoretical contribution proceeds in three parts, which answer in turn the
	questions of inference, optimality, and interpretation. First, because the window
	functional is linear in the regression function, a change of variables yields an
	exact Riesz representation with a bounded, explicitly computable representer ---
	no integration by parts, no boundary terms, and no derivatives of the covariate
	density --- and the bias of the resulting orthogonal moment is exactly a product
	of two $L^2$ estimation errors. A cross-fitted estimator built on this score is
	root-$n$ consistent and asymptotically normal whenever the learner and a
	Riesz-regression estimate of the representer each converge faster than
	$n^{-1/4}$ in $L^2$ --- a condition of the kind modern learning theory actually
	verifies, in place of the parametric-rate function-space convergence that a
	delta-method treatment would demand and that no growing-complexity learner
	satisfies. Second, the score is the efficient influence function: the estimator
	attains the semiparametric efficiency bound for the window effect, whose variance
	decomposes into the heterogeneity of the effect and a noise term amplified by the
	representer, making the cost of narrow windows explicit. The score is moreover
	doubly robust: with a known or parametrically estimated representer --- simulation
	designs, designed experiments, randomized treatments --- inference remains valid
	with no consistency requirement on the learner at all, covering the
	fixed-complexity defaults under which off-the-shelf learners converge to
	pseudo-true limits rather than to the truth. Third, a pairing
	identity characterizes what the marginal effect of a non-differentiable model
	measures: for any estimator of bounded variation --- every tree-based model ---
	the window effect equals the jump measure of the fitted surface integrated
	against a locally averaged covariate density, an exact geometric description of
	the population quantity that a forest's coefficient tracks. Supporting this
	interpretation theory, gradient consistency is established for neural networks
	\citep{hornik1990universal} in the Sobolev space $W^{1,p}$, via the
	Rellich--Kondrachov compact embedding theorem \citep{evans2010pde}, and for
	gradient boosted models with tree base learners and Mondrian forests
	\citep{lakshminarayanan2014mondrian, mourtada2017mondrian} in the space of
	functions of bounded variation \citep{ambrosio2000bv}, via the BV compactness
	theorem. For Breiman forests \citep{breiman2001rf, scornet2015forests},
	gradient consistency is conjectured rather than proved, with the obstruction
	identified precisely; because the inference theory requires gradient
	consistency for no class, this leaves two Breiman questions of unequal
	difficulty --- valid inference needs only an $L^2$ rate from the forest, while
	the gradient question, which bears on interpretation rather than on coverage,
	remains a function-space problem and remains open.
	
	We validate the framework empirically on three datasets spanning binary classification
	in social science, binary classification in finance, and continuous regression: the
	UCI Adult Income dataset \citep{kohavi1996, UCI}, the UCI Credit Card Default dataset
	\citep{yeh2009, UCI}, and the Ames Housing dataset \citep{decock2011}. Across all
	three applications we compare AMEs from logistic regression, random forests
	\citep{breiman2001rf}, XGBoost \citep{chen2016xgboost}, and neural networks against
	SHAP values \citep{lundberg2017shap} and partial dependence plot slopes
	\citep{friedman2001pdp}. The debiased AMEs agree closely across the four model classes,
	which is what one expects once every learner is used to estimate the same window
	effect rather than to report its own slope, while SHAP values are on a different
	scale and in several cases have the wrong sign --- a systematic divergence that is
	most visible in the regression setting where units are dollars and the magnitudes are
	large enough to make the discrepancy concrete. PDP slopes, being themselves plug-in
	averages over a single full-sample fit, track the debiased AME where the plug-in is
	nearly unbiased and fall short of it, usually toward zero, wherever regularization
	flattens the fitted surface --- the regularization bias made directly visible.
	Because AMEs come with standard errors and SHAP values and PDP slopes do not, the AME
	framework supports conclusions that the alternatives cannot. An additional finding
	concerns The credit default application is the sharpest case, and it is worth stating at
	the outset. Under the plug-in the four learners put the effect of the payment ratio at
	$-$0.431 and $-$0.440 for the two tree ensembles, $-$0.045 for logistic regression, and
	0.0009 --- effectively nothing --- for the neural network: a set of estimates that
	cannot all be describing the same world, and among which a reader had no principled way
	to adjudicate. Under the correction they become $-$0.265, $-$0.489, $-$0.584 and
	$-$0.194, every one of them negative and substantial. The disagreement was not four
	models detecting different structure. It was the regularization bias
	$\E[\alpha_h(\hat{f} - f)]$, which differs across learners because their regularization
	differs, and which the orthogonal score removes.
	Breiman random forests, which are the learner the framework is most often asked
	to carry and the one for which the interpretation theory is least complete. Under
	the uncorrected plug-in they behave worst of the four classes: the simulation
	study shows a bias that does not vanish with the sample size but grows with it,
	from $-0.097$ at $n=250$ to $-0.091$ at $n=5{,}000$ in the linear regression
	design, so that intervals shrink at the usual rate around a center that is not
	moving and coverage collapses to $0.130$. The orthogonal correction removes that
	term. Estimating the same quantity on the same draws with the debiased estimator
	leaves a bias in the fourth decimal at every sample size and returns coverage to a
	mean of $0.948$ against the nominal $0.95$, within the band $0.922$ to $0.968$
	across all one hundred and twenty forest cells of the design --- every data
	generating process, sample size, feature and outcome type.
	This is more than a bias correction, and the reason matters for which learners
	the framework can be used with. A depth-constrained forest whose plug-in error
	grows with $n$ is not converging to $f$, and Proposition~\ref{prop:robust} says
	that with the representer known this costs nothing: the estimator is root-$n$
	normal about $\theta_h(f)$ even when $\hat{f}$ converges to some other function
	entirely. The simulation is that proposition observed rather than merely a
	demonstration that the correction is well calibrated. Because the inference
	theory consumes gradient consistency for no learner class, the Breiman forest is
	the natural choice for applied work rather than a concession, and what remains
	open for that class is the separate interpretive question of whether its window
	effect converges to a classical derivative as $h \to 0$.
	
	The remainder of the paper is organized as follows. Section~2 defines the estimand, presents the estimator, and establishes its relationship to classical average marginal effects for the GLM family. Section~3 develops the asymptotic theory: the window estimand and its Riesz representation, the orthogonal score with cross-fitting, root-$n$ asymptotic normality and semiparametric efficiency, the pairing identity for estimators of bounded variation, gradient consistency for three learner classes, and the nuisance rates that the conditions require. Section~4 presents the three empirical applications. Section~5 presents simulation evidence on 	finite-sample performance. Section~6 concludes.
	
	\section{Method}
	
	We propose a framework for interpreting any supervised machine learning model through the lens of average marginal effects. The framework is agnostic to the choice of learner --- the same estimation procedure applies whether $\hat{f}$ is a neural network, a gradient boosted model, a random forest, or any other supervised learner. Crucially, the framework produces standard errors from the influence function of a Neyman-orthogonal score, evaluated by cross-fitting, enabling hypothesis tests and confidence intervals that existing model-agnostic interpretability methods do not support. A refitting bootstrap is retained, but as a diagnostic rather than as the basis for inference: Section~\ref{sec:inference} shows that it is blind to the learner's regularization bias, which is precisely the term the orthogonal correction removes. SHAP values and permutation-based feature importance measures provide point estimates of variable importance but no basis for inference. The average marginal effect estimator proposed here yields a coefficient table --- estimate, standard error, and significance --- directly analogous to the output of a parametric regression, with no restrictions on the functional form of $\hat{f}$.
	
	\subsection{Estimand}
	
	Let $\mathbf{x}_i = (x_{i1}, \dots, x_{ip})^\top \in \mathbb{R}^p$ denote the feature vector for observation $i$, and let $y_i \in \mathbb{R}$ denote the outcome. We observe $n$ independent draws $\{(\mathbf{x}_i, y_i)\}_{i=1}^n$ from some population distribution $P$. Let $f: \mathbb{R}^p \to \mathbb{R}$ denote the conditional expectation function $f(\mathbf{x}) = E[Y \mid \mathbf{x}]$. For binary outcomes, $f(\mathbf{x}) = P(Y = 1 \mid \mathbf{x})$ is the conditional probability of the positive class; for continuous outcomes, $f(\mathbf{x}) = E[Y \mid \mathbf{x}]$ directly. In both cases $f$ is the same object --- the conditional expectation of $Y$ given $\mathbf{x}$ --- and the method requires no modification across outcome types.
	
	The marginal effect of variable $j$ at observation $i$ is the partial derivative of $f$ evaluated at $\mathbf{x}_i$:
	\begin{equation}
		m_{ij} = \frac{\partial f}{\partial x_j}\bigg|_{\mathbf{x}_i}
	\end{equation}
	The population average marginal effect of variable $j$ is:
	\begin{equation}
		\mathrm{AME}_j = \mathbb{E}\left[\frac{\partial f}{\partial x_j}(\mathbf{x})\right]
	\end{equation}
	where the expectation is taken over the marginal distribution of $\mathbf{x}$ under $P$. This is the classical target. The operational estimand of the paper is its finite-window counterpart, the \emph{window average marginal effect} $\theta_{j,h} = \mathbb{E}[w(\mathbf{x})\,\{f(\mathbf{x} + h e_j) - f(\mathbf{x} - h e_j)\}/2h]$, where $h > 0$ is a fixed step scaled to the spread of $x_j$ and $w$ is a trimming weight that excludes points within $h$ of the boundary of the support. The window effect is well defined for any bounded prediction function --- no derivative need exist anywhere, which is what a genuinely model-agnostic estimand requires --- and it converges to $\mathrm{AME}_j$ as $h \to 0$ whenever $f$ has one weak derivative. Both facts, and the conditions under which the estimators below converge, are established in Section~3.
	
	\subsection{Estimation}
	
	Given an estimator $\hat{f}$ of $f$, the sample analog is:
	\begin{equation}
		\widehat{\mathrm{AME}}_j = \frac{1}{n} \sum_{i=1}^n \frac{\partial \hat{f}}{\partial x_j}\bigg|_{\mathbf{x}_i}
	\end{equation}
	For most supervised learners the partial derivative $\partial \hat{f} / \partial x_j$ is not available in closed form, and for tree-based models it does not exist everywhere in the classical sense. We therefore estimate it via a centered finite difference:
	\begin{equation}
		\frac{\partial \hat{f}}{\partial x_j}\bigg|_{\mathbf{x}_i} \approx \frac{\hat{f}\!\left(\mathbf{x}_i^{(j,\, x_{ij} + h_j)}\right) - \hat{f}\!\left(\mathbf{x}_i^{(j,\, x_{ij} - h_j)}\right)}{2h_j}
	\end{equation}
	where $\mathbf{x}_i^{(j,v)}$ denotes the vector $\mathbf{x}_i$ with the $j$-th component replaced by $v$. The centered finite difference is preferred over a one-sided difference because its approximation error relative to the pointwise derivative is $O(h_j^2)$ rather than $O(h_j)$, which becomes material when $h_j$ is not negligibly small. Relative to the window estimand $\theta_{j,h}$ the centered difference involves no approximation at all: it is the defining operation of the estimand, computed exactly, and $h_j$ is a reported feature of the target --- the perturbation scale at which effects are measured --- rather than a numerical tolerance.
	
	The step size $h_j$ is chosen adaptively as:
	\begin{equation}
		h_j = \max\!\left(10^{-4},\; 0.05 \cdot \hat{\sigma}_j\right)
	\end{equation}
	where $\hat{\sigma}_j$ is the sample standard deviation of the $j$-th feature. Scaling $h_j$ to the spread of $x_j$ ensures the finite difference is neither so small as to introduce numerical instability nor so large as to obscure local behavior. The floor of $10^{-4}$ guards against degenerate cases where $\hat{\sigma}_j \approx 0$. Observations within $h_j$ of the boundary of the observed support receive trimming weight zero, so that $\hat{f}$ is never evaluated at points where the regression function is unidentified; the trimmed shell carries probability mass of order $h_j$, and with the default step this exclusion is small. Untrimmed evaluation --- the default in earlier implementations --- silently extrapolates the fitted model outside the support, and the trimming weight should be understood as the repair of that extrapolation rather than as a restriction of the method.
	
	For binary or categorical variables the partial derivative does not exist and the finite difference is replaced by a contrast. The marginal effect for observation $i$ is:
	\begin{equation}
		m_{ij} = \hat{f}\!\left(\mathbf{x}_i^{(j,1)}\right) - \hat{f}\!\left(\mathbf{x}_i^{(j,0)}\right)
	\end{equation}
	which measures the change in predicted outcome when $x_j$ is switched from 0 to 1 holding all other variables at their observed values. The sample average marginal effect is then $\widehat{\mathrm{AME}}_j = n^{-1}\sum_{i=1}^n m_{ij}$, with interpretation identical to the continuous case: a change in $x_j$ from 0 to 1 is associated with a $\widehat{\mathrm{AME}}_j$ change in the predicted outcome, on average across the sample.
	
	For binary classification models, $\hat{f}$ must return the estimated conditional probability $\hat{P}(Y=1 \mid \mathbf{x})$ rather than the predicted class label. Class labels are discrete and their finite differences are identically zero; the method requires a continuous output. This is the predicted probability of the positive class, available as \texttt{predict\_proba(X)[:,1]} in scikit-learn and equivalent interfaces in other frameworks. The AME is then interpreted as the change in predicted probability associated with a one-unit change in $x_j$, on average across the sample.
	
	\subsection{Relationship to Classical Average Marginal Effects}
	
	Average marginal effects are a standard tool in the econometrics of nonlinear models \citep{greene2012, wooldridge2010}. For parametric models the estimator $\widehat{\mathrm{AME}}_j$ defined above reduces exactly to the classical result. Consider a generalized linear model with inverse link function $\mu$ and linear predictor $\eta_i = \mathbf{x}_i^\top \boldsymbol{\beta}$, so that $f(\mathbf{x}_i) = \mu(\eta_i)$. The marginal effect of variable $j$ at observation $i$ is:
	\begin{equation}
		\frac{\partial f}{\partial x_j}\bigg|_{\mathbf{x}_i} = \beta_j \cdot \mu'(\eta_i)
	\end{equation}
	and the population average marginal effect is $\mathrm{AME}_j = \beta_j \cdot E[\mu'(\eta)]$. For the identity link $\mu'(\eta) = 1$ and $\mathrm{AME}_j = \beta_j$ --- the marginal effect is constant across observations and equals the OLS coefficient. For the logistic link $\mu'(\eta) = \Lambda(\eta)(1 - \Lambda(\eta))$ and the AME averages the logistic density over the sample \citep{greene2012}. For the probit link $\mu'(\eta) = \phi(\eta)$, the standard normal density evaluated at the linear predictor \citep{wooldridge2010}.
	
	In each case $\widehat{\mathrm{AME}}_j$ recovers the classical result automatically --- the finite difference applied to the model output numerically approximates $\beta_j \cdot \mu'(\eta_i)$ without requiring knowledge of the link function. The framework therefore nests the entire GLM family as a special case and reduces to familiar output when the model is parametric. This backward compatibility is what the \texttt{marginaleffects} package in R and Python implements \citep{arelbundock2024marginaleffects}; the present framework extends that implementation with formal asymptotic guarantees for machine learning models.
	
	\subsection{Inference}
	
	Inference is based on a debiased, cross-fitted estimator whose formal analysis is
	the subject of Section~3; here we describe the procedure. Partition the sample
	into $K$ folds ($K = 5$ by default). For each fold $k$, using only the
	observations outside the fold: (1) fit the learner $\hat{f}^{(-k)}$ exactly as it
	would be fit in any application, including hyperparameter selection by
	cross-validation; (2) fit an estimate $\hat{\alpha}^{(-k)}$ of the Riesz
	representer of the window effect --- an explicit weight function characterized in
	Section~3 --- by minimizing an empirical quadratic loss that requires only
	evaluating candidate functions at the shifted points $\mathbf{x}_i \pm h_j e_j$,
	with no density estimation of any kind. The estimate of $\theta_{j,h}$ averages,
	over the held-out observations of every fold, the finite difference of the
	out-of-fold learner plus a correction term:
	\begin{equation}
		\hat{\theta}_{j,h} = \frac{1}{n} \sum_{k=1}^{K} \sum_{i \in I_k} \left[
		w(\mathbf{x}_i)\,
		\frac{\hat{f}^{(-k)}\!\left(\mathbf{x}_i^{(j,\, x_{ij} + h_j)}\right)
		- \hat{f}^{(-k)}\!\left(\mathbf{x}_i^{(j,\, x_{ij} - h_j)}\right)}{2 h_j}
		+ \hat{\alpha}^{(-k)}(\mathbf{x}_i)\,
		\bigl(y_i - \hat{f}^{(-k)}(\mathbf{x}_i)\bigr)
		\right]
	\end{equation}
	The first term is the plug-in average of Section~2.2 computed on held-out data;
	the second term removes, at first order, the bias that regularization of the
	learner transmits to the plug-in average --- each held-out residual, weighted by
	the representer, testifies to how the learner is locally mis-calibrated where
	that mis-calibration matters for the marginal effect. For binary and categorical
	variables the same formula applies with the contrast in place of the finite
	difference and inverse propensity weights in the role of the representer.
	
	Standard errors require no resampling: the summands above, centered at
	$\hat{\theta}_{j,h}$, are estimates of the influence function values, and the
	standard error is the sample standard deviation of those values divided by
	$\sqrt{n}$. The coefficient table --- estimate, standard error, $t$-statistic,
	and $p$-value for each feature --- is therefore produced by $K$ model fits in
	total, in place of the hundreds of refits that a bootstrap requires; Section~3
	establishes that the resulting confidence intervals are asymptotically valid and
	that the estimator attains the semiparametric efficiency bound. When a full
	resampling distribution is wanted --- for simultaneous confidence bands across
	features, for example --- a multiplier bootstrap that perturbs the estimated
	influence function values with i.i.d.\ mean-zero weights is available at no
	additional model fits.
	
	The package additionally retains a nonparametric model-refitting bootstrap
	\citep{efron1979, efrontibshirani1993} as a \emph{diagnostic} for specification
	and hyperparameter-selection uncertainty, which the asymptotic theory conditions
	away. Let $\hat{\lambda}$ denote the hyperparameters selected by cross-validation
	on the observed sample $\mathcal{D}$. For $b = 1, \dots, B$: draw
	$\mathcal{D}^{(b)}$ by resampling $(\mathbf{x}_i, y_i)$ pairs jointly with
	replacement; refit $\hat{f}^{(b)}$ on $\mathcal{D}^{(b)}$ with hyperparameter
	search warm-started at $\hat{\lambda}$; and recompute the plug-in average
	marginal effect. The dispersion of the replicates measures how sensitive the
	reported effects are to the joint variability of resampling, refitting, and
	hyperparameter selection, and large dispersion is a useful warning that the
	fitted model is unstable. It is not, however, a substitute for the influence
	function standard errors above, for reasons Section~3 makes precise: the
	refitting bootstrap recenters at the fitted model and is therefore blind to the
	regularization bias that the correction term removes --- a blindness whose
	consequences for coverage are documented for random forests in Section~5 --- and,
	because each replicate is evaluated at the original sample, it omits the
	component of sampling variance that comes from averaging a heterogeneous effect
	over a finite sample of covariate values.
	
	\section{Asymptotic Theory}
	\label{sec:theory}
	
	\subsection{Overview}
	
	The asymptotic theory must deliver three things. First, a population estimand that
	is well defined for any supervised learner, including learners such as trees and
	forests whose prediction functions are not differentiable anywhere. Second, valid
	root-$n$ inference on that estimand under conditions that flexible machine learning
	estimators can actually satisfy --- in particular, without requiring the learner
	itself to converge at the parametric rate, a condition that fails whenever the
	complexity of the function class grows with the sample size. Third, an account of
	what the gradient of a fitted model measures: when the average marginal effect of a
	fitted forest is computed, what population quantity is that number tracking? This
	section delivers the three in turn.
	
	The organizing observation is structural: the average marginal effect is a
	\emph{linear} functional of the regression function. Linearity has three
	consequences that together carry the entire theory. First, the functional admits an
	exact Riesz representation --- there is an explicit, bounded weight function
	$\alpha_h$ such that averaging the finite difference of any function against the
	covariate distribution equals averaging the function itself against $\alpha_h$
	(Lemma~\ref{lem:riesz}). The representation is obtained by a change of variables
	rather than integration by parts, so no boundary terms arise and no derivative of
	the covariate density is ever needed. Second, the estimation bias of a debiased
	moment built from this representation is \emph{exactly} a product of two estimation
	errors (Lemma~\ref{lem:orth}) --- not a product plus a linearization remainder, as
	for nonlinear functionals, but an identity. Root-$n$ inference therefore requires
	only that the product of the $L^2$ errors of two nuisance estimates be
	$o_p(n^{-1/2})$ --- for instance, both converging faster than $n^{-1/4}$ --- rather
	than root-$n$ convergence of the learner in any norm
	(Theorem~\ref{thm:main}). Third, the resulting estimator attains the semiparametric
	efficiency bound for the estimand (Theorem~\ref{thm:efficiency}): among all regular
	estimators, none has smaller asymptotic variance.
	
	The estimand the theory analyzes is exactly the quantity the estimator of Section~2
	computes: the average of a centered finite difference at a fixed step $h$, which we
	call the \emph{window average marginal effect}. Treating $h$ as part of the estimand
	rather than as a numerical approximation to a derivative has two payoffs. The window
	effect is well defined for any bounded prediction function --- no differentiability
	anywhere --- which is what makes a genuinely model-agnostic theory possible. And the
	classical average derivative is recovered as the $h \to 0$ limit whenever the true
	regression function has one weak derivative (Proposition~\ref{prop:hlimit}), so
	nothing classical is given up. A trimming weight confines the evaluation points to
	the interior of the support; the binary and categorical contrasts of Section~2 fit
	the same template, with the change of variables running over a two-point support and
	the representer reducing to inverse-probability weights
	(Remark~\ref{rem:binary}).
	
	The interpretation question is answered by a pairing identity
	(Theorem~\ref{thm:pairing}). For any estimator of bounded variation --- which
	includes every tree-based model --- the window average marginal effect of the
	\emph{fitted} model equals the integral of a locally averaged covariate density
	against the distributional gradient of the fitted model: for a tree, a sum of leaf
	boundary jumps weighted by the boundary's share of smoothed covariate mass. This
	makes precise, apparently for the first time, which population quantity the marginal
	effect of a piecewise constant learner measures. The Sobolev and BV lemmas of
	Section~\ref{sec:interpretation} then characterize the regime in which the learner's own
	gradient converges to the population gradient --- the strongest interpretability
	guarantee available, and the condition under which the uncorrected plug-in estimator
	of Section~2 is already consistent. Those results describe when the correction term
	is a refinement; the inference theory of Theorem~\ref{thm:main} does not depend on
	them.
	
	The relationship between the debiased estimator and the plug-in-plus-bootstrap
	procedure of Section~2 is worth stating plainly. The plug-in average is the first
	component of the orthogonal score; the correction term
	$\alpha_h(x_i)\,(y_i - \hat{f}(x_i))$ removes, at first order, the bias transmitted
	to the average marginal effect by regularization of the learner. The simulation
	evidence in Section~5 displays exactly this mechanism: random forest confidence
	intervals built from the uncorrected plug-in collapse to near-zero coverage at large
	$n$ because the intervals concentrate around a regularization-biased center --- the
	bias term the correction is constructed to remove. The model-refitting bootstrap of
	Section~2 remains a useful diagnostic for specification uncertainty, but the
	inferential guarantees of this section attach to the debiased estimator and its
	influence function, not to the uncorrected plug-in.
	
	\subsection{Notation and Preliminaries}
	
	Let $\Omega \subset \R^d$ be a bounded Lipschitz domain and let $f: \Omega \to \R$
	be the true underlying function relating covariates to outcomes. Let
	$\fn: \Omega \to \R$ denote an estimator of $f$ constructed from $n$ observations
	$\{(x_i, y_i)\}_{i=1}^n$. We assume throughout that $\Omega$ is bounded and
	Lipschitz, which ensures the compactness and embedding results we rely on are
	available. These are mild regularity conditions satisfied by any compact covariate
	space with reasonable boundary behavior.
	
	Throughout, $(\mathbf{x}_i, y_i)_{i=1}^n$ are independent draws from $P$, the joint
	distribution of $(X, Y)$; with mild abuse of notation $P$ also denotes the marginal
	distribution of the covariates, with $\norm{g}_{P,2}^2 = \E[g(X)^2]$ the
	corresponding $L^2$ norm. We fix a coordinate $j \in \{1, \dots, d\}$ for the
	remainder of the section and suppress it from the notation where no confusion
	arises. For a step size $h > 0$, define the centered difference operator:
	\begin{equation}
		\Dh g(x) = \frac{g(x + h e_j) - g(x - h e_j)}{2h}
	\end{equation}
	where $e_j$ is the $j$-th coordinate vector, and let
	$\Omega_{j,h} = \{x \in \Omega : x + h e_j \in \Omega \text{ and }
	x - h e_j \in \Omega\}$ denote the set of points whose shifted evaluations remain in
	the domain. A \emph{trimming weight} is a measurable function
	$w : \Omega \to [0,1]$ with $\operatorname{supp}(w) \subseteq \Omega_{j,h}$; it
	confines the finite difference to points where both shifted arguments lie inside the
	support, so that the prediction function is never evaluated where the regression
	function is unidentified. The following conditions are maintained throughout the
	section.
	
	\begin{assumption}
		\label{ass:sampling}
		The data are i.i.d.\ draws from $P$ with $\E[Y^2] < \infty$ and
		$\mathrm{Var}(Y \mid X = x) \leq \bar{\sigma}^2 < \infty$ for $P$-almost every
		$x$. The covariate distribution admits a density $p$ on $\Omega$ with
		$0 < p_{\min} \leq p(x) \leq p_{\max} < \infty$ for all $x \in \Omega$.
	\end{assumption}
	
	The density bounds are the same condition previously used to define the average
	marginal effect functional and are standard in the average derivative literature
	\citep{powell1989, newey1993}. They enter the theory only through the ratio
	$p_{\max}/p_{\min}$, which bounds the Riesz representer below.
	
	We work with two function spaces depending on the estimator class. The choice of
	function space is not arbitrary---each space is the natural habitat of the estimator
	class in question, in the sense that estimators of that class belong to the space
	with probability one under mild conditions.
	
	The \textbf{Sobolev space} $W^{k,p}(\Omega)$ for $k \geq 1$ and $p \geq 1$
	consists of functions whose weak derivatives up to order $k$ are in $\Lp$, with
	norm:
	\begin{equation}
		\norm{g}_{W^{k,p}} = \sum_{|\alpha| \leq k} \norm{\partial^\alpha g}_{L^p}
	\end{equation}
	The gradient $\nabla g \in L^p(\Omega)^d$ is a vector valued $L^p$ function for any
	$g \in \Wonep$. Sobolev spaces are the natural setting for smooth estimators because
	they quantify not just how well the function is approximated in $L^p$ but how well
	its derivatives are approximated simultaneously. This is exactly the property needed
	to justify gradient based interpretation---$L^p$ consistency of $\fn$ alone is
	insufficient because differentiation is an unbounded operator on $L^p$, and
	convergence of a function does not in general imply convergence of its derivatives
	without additional regularity \citep[Chapter~5]{evans2010pde}.
	
	The \textbf{space of functions of bounded variation} $\BV$ consists of functions
	$g \in \Lone$ whose distributional gradient $Dg$ is a vector valued Radon measure
	with finite total variation $|Dg|(\Omega) < \infty$. The $BV$ norm is:
	\begin{equation}
		\norm{g}_{BV} = \norm{g}_{L^1} + |Dg|(\Omega)
	\end{equation}
	The $BV$ framework extends the notion of gradient to functions that are not
	differentiable in the classical or even weak sense---precisely the situation with
	tree based estimators, which have jump discontinuities at leaf boundaries. For
	piecewise constant functions the distributional gradient is a sum of surface measures
	concentrated on the boundaries between constant regions, capturing the intuition that
	the ``derivative'' of a step function is a spike at the jump location. This makes
	$BV$ the natural space for random forests and boosting with tree base learners. The
	total variation $|Dg|(\Omega)$ measures the total size of all jumps, weighted by the
	boundary surface area---a quantity that is finite and well controlled for any
	piecewise constant function with finitely many pieces. The standard reference for
	the theory of $BV$ functions is \citet{ambrosio2000bv}.
	
	The relationship between the two spaces is important. If $f \in \Wonep$ for
	$p \geq 1$, then $f \in \BV$ with $Df = \nabla f \cdot \mathcal{L}^d$---that is,
	the distributional gradient of $f$ is absolutely continuous with respect to Lebesgue
	measure with density $\nabla f$. This embedding $W^{1,p} \hookrightarrow BV$ means
	that assuming $f \in \Wonep$ is sufficient for both the Sobolev and BV results below.
	We make this assumption throughout.
	
	\begin{definition}[Well-Behaved Estimator, Sobolev Sense]
		\label{def:wellbehaved_sobolev}
		We say that $\fn$ is \emph{well-behaved in the Sobolev sense} if:
		\begin{equation}
			\sup_n \norm{\fn}_{W^{2,p}} < \infty \quad \text{and} \quad
			\norm{\fn - f}_{L^p} \xrightarrowp 0
		\end{equation}
	\end{definition}
	
	\begin{definition}[Well-Behaved Estimator, BV Sense]
		\label{def:wellbehaved_bv}
		We say that $\fn$ is \emph{well-behaved in the BV sense} if:
		\begin{equation}
			\sup_n \E\left[|D\fn|(\Omega)\right] < \infty \quad \text{and} \quad
			\norm{\fn - f}_{L^1} \xrightarrowp 0
		\end{equation}
	\end{definition}
	
	In both cases the well-behaved condition has the same structure: a uniform
	boundedness condition in a strong space combined with consistency in a weaker space.
	The uniform boundedness condition is the key regularity requirement---it says that
	the estimator sequence does not become increasingly irregular as $n$ grows, even as
	it fits the data more closely. This condition is verifiable from properties of the
	estimator class and is established for each class in the theorems below. The
	consistency condition is the standard requirement that the estimator converges to the
	truth and is available from existing results in each case.
	
	\subsection{The Window Estimand and Its Riesz Representation}
	\label{sec:window}
	
	\begin{definition}[Window Average Marginal Effect]
		\label{def:window}
		Fix $h > 0$ and a trimming weight $w$ with
		$\operatorname{supp}(w) \subseteq \Omega_{j,h}$. The \emph{window average
		marginal effect} of variable $j$ at scale $h$ is:
		\begin{equation}
			\theta_h(f) = \E\left[w(X)\, \Dh f(X)\right]
		\end{equation}
	\end{definition}
	
	Up to the trimming weight, $\theta_h(f)$ is exactly the population counterpart of
	the quantity the estimator of Section~2 computes: the finite difference at step $h$
	is not a numerical approximation to a derivative that the theory then ignores, but
	the defining operation of the estimand itself. The step size is a fixed feature of
	the estimand, chosen and reported by the analyst exactly as a bandwidth is, and the
	default $h_j = 0.05\,\hat{\sigma}_j$ of Section~2 corresponds to asking how the
	predicted outcome responds to displacements of a twentieth of a standard deviation
	--- an economically interpretable perturbation scale. Two features deserve emphasis.
	First, $\theta_h(f)$ is well defined for any bounded measurable $f$: no
	differentiability is assumed anywhere, which is what permits a single estimand to
	cover neural networks and regression trees simultaneously. Second, the trimming
	weight repairs a defect of the untrimmed implementation, which evaluates
	$\hat{f}$ at points $x_i \pm h e_j$ that may fall outside the support of the
	covariates, where the regression function is not identified; $w$ excises a boundary
	shell of probability mass $O(h)$ and confines every evaluation to the region where
	$f$ means something.

	\begin{remark}[Data-dependent step sizes]
		\label{rem:randomh}
		Definition~\ref{def:window} fixes $h$, and every result below is a
		statement about the estimand at that fixed value. The default of
		Section~2 sets $h_j = \max(10^{-4},\, 0.05\,\hat{\sigma}_j)$, which is a
		function of the data. The two are reconciled by reading all results
		conditional on the realized $\hat{h}$: the reported quantity is the
		window effect at the step actually used, which is why $h_j$ is reported
		alongside every estimate in Table~\ref{tab:window_settings} rather than
		being treated as a tuning parameter. Nothing in the argument requires
		$h$ to be non-random --- the representation of
		Lemma~\ref{lem:riesz}, the orthogonality of Lemma~\ref{lem:orth} and the
		expansion of Theorem~\ref{thm:main} all hold pointwise in $h$ --- and
		the conditional reading is the natural one, because the estimand is
		genuinely indexed by the perturbation scale the analyst chose.

		It is worth being explicit about what the conditional reading buys,
		since the alternative is not free. If one insists instead that the
		target is $\theta_{h_0}$ at the deterministic limit
		$h_0 = 0.05\,\sigma_j$, then whenever $h \mapsto \theta_h(f)$ is
		differentiable at $h_0$ the delta method gives
		$\theta_{\hat{h}} - \theta_{h_0} = \theta_{h_0}'\,(\hat{h} - h_0) +
		o_p(n^{-1/2})$, and since $\hat{\sigma}_j$ is root-$n$ consistent this
		is $O_p(n^{-1/2})$ --- the same order as the sampling error itself. The
		substitution is therefore \emph{not} asymptotically negligible: it
		contributes a term to the limiting variance, and an unconditional
		treatment would have to add that term to $V_h$. We do not pursue this,
		because the conditional statement is the one that matches how the
		estimand is used and reported. An analyst who wants the unconditional
		reading can fix $h$ from prior information or from a pilot sample, at
		which point the two coincide.
	\end{remark}
	
	The estimand admits an exact dual representation, which is the engine of everything
	that follows. Extend $q = wp$ by zero outside $\Omega$.
	
	\begin{lemma}[Riesz Representation]
		\label{lem:riesz}
		Let Assumption~\ref{ass:sampling} hold and define:
		\begin{equation}
			\alpha_h(u) = \frac{q(u - h e_j) - q(u + h e_j)}{2h\, p(u)},
			\qquad u \in \Omega
		\end{equation}
		Then for every $g \in L^1(P)$:
		\begin{equation}
			\E\left[w(X)\, \Dh g(X)\right] = \E\left[\alpha_h(X)\, g(X)\right]
			\label{eq:riesz}
		\end{equation}
		Moreover $\norm{\alpha_h}_\infty \leq p_{\max} / (h\, p_{\min})$ and
		$\E[\alpha_h(X)] = 0$.
	\end{lemma}
	
	The proof is a change of variables in each half of the finite difference and nothing
	more. This is worth dwelling on, because it is the point at which the present theory
	parts ways with the classical average derivative literature. For the derivative
	functional $\E[\partial_j f(X)]$, the analogous representation is obtained by
	integration by parts, the representer is the negative score $-\partial_j \log p$ of
	the covariate density, and a boundary term appears that must be killed by assuming
	the density vanishes on the boundary --- an assumption incompatible with a density
	bounded away from zero on a compact support, and the reason
	\citet{powell1989} pass to density-weighted variants. For the window functional the
	integration by parts is replaced by exact summation: no boundary term ever arises,
	because the trimming weight guarantees that every shifted evaluation point stays
	inside $\Omega$. The representer involves the density only through ratios of its
	values at points $2h$ apart --- no derivative of $p$, and hence no density
	derivative estimation, is ever required. The mean-zero property follows from
	applying \eqref{eq:riesz} to $g \equiv 1$, and the uniform bound makes $\alpha_h$ a
	bounded weight function, in contrast to the unbounded score-based representers of
	the derivative functional. The cost of the exactness is the factor $h^{-1}$ in the
	bound: smaller windows mean noisier representers, the familiar bias--variance logic
	of a bandwidth, here appearing in the variance of the efficient influence function
	rather than in a smoothing bias.
	
	The classical estimand is recovered in the small-window limit.
	
	\begin{proposition}[Recovery of the Average Derivative]
		\label{prop:hlimit}
		Let $f \in W^{1,1}(\Omega)$ and let $w$ be a trimming weight with
		$\operatorname{supp}(w) \subseteq \Omega_{j,h_0}$ for some $h_0 > 0$. Then:
		\begin{equation}
			\lim_{h \downarrow 0} \theta_h(f)
			= \E\left[w(X)\, \partial_j f(X)\right]
		\end{equation}
		If moreover $(w_m)_{m \geq 1}$ is a sequence of trimming weights with
		$w_m \uparrow 1$ $P$-almost everywhere, then
		$\E[w_m\, \partial_j f] \to \E[\partial_j f] = \mathrm{AME}_j$.
	\end{proposition}
	
	One weak derivative --- the minimal regularity under which the classical estimand
	exists at all --- is thus the only condition separating the window effect from the
	average derivative. Nothing about the fitted model appears in
	Proposition~\ref{prop:hlimit}: the double limit concerns the estimand, not the
	estimator, and in practice one fixes $h$ and reports the window effect, treating the
	proposition as the license to interpret it as a derivative summary when the truth is
	smooth.
	
	\begin{remark}[Binary and Categorical Variables]
		\label{rem:binary}
		For a binary $X_j \in \{0,1\}$ the contrast estimand of Section~2 is
		$\theta^{\mathrm{bin}} = \E[w(X_{-j})\{f(X_{-j}, 1) - f(X_{-j}, 0)\}]$, and the
		same change-of-variables argument --- with the integral over $x_j$ replaced by a
		sum over its two support points --- yields the representer:
		\begin{equation}
			\alpha^{\mathrm{bin}}(u) = w(u_{-j})\left[
			\frac{\mathbf{1}\{u_j = 1\}}{\pi(u_{-j})}
			- \frac{\mathbf{1}\{u_j = 0\}}{1 - \pi(u_{-j})}\right],
			\qquad \pi(x_{-j}) = \Prob(X_j = 1 \mid X_{-j} = x_{-j})
		\end{equation}
		These are inverse propensity weights, and the resulting orthogonal score is the
		augmented inverse probability weighting score of \citet{robins1994} with $X_j$
		in the role of the treatment. The overlap condition
		$\varepsilon \leq \pi \leq 1 - \varepsilon$ plays exactly the role the trimming
		weight plays in the continuous case, and every result of
		Section~\ref{sec:inference} applies verbatim with
		$(\alpha_h, \theta_h)$ replaced by
		$(\alpha^{\mathrm{bin}}, \theta^{\mathrm{bin}})$. Continuous, binary, and
		categorical effects are thus instances of one construction rather than a
		definition and two patches.
	\end{remark}
	
	\subsection{Orthogonal Score, Cross-Fitting, and Efficient Inference}
	\label{sec:inference}
	
	The Riesz representation converts the estimand into a moment condition with an
	exact second-order bias. Define, for candidate functions
	$\tilde{f}, \tilde{\alpha}$ and scalar $\theta$, the score:
	\begin{equation}
		\psi(x, y; \tilde{f}, \tilde{\alpha}, \theta)
		= w(x)\, \Dh \tilde{f}(x)
		+ \tilde{\alpha}(x)\,\bigl(y - \tilde{f}(x)\bigr) - \theta
		\label{eq:score}
	\end{equation}
	At the truth, $\E[\psi(X, Y; f, \alpha_h, \theta_h(f))] = 0$: the first term has
	mean $\theta_h(f)$ by definition and the correction term has conditional mean zero.
	The content of the construction is what happens away from the truth.
	
	\begin{lemma}[Exact Product Bias]
		\label{lem:orth}
		Let Assumption~\ref{ass:sampling} hold and let
		$\tilde{f} \in L^2(\Omega)$, $\tilde{\alpha} \in L^2(P)$. Then:
		\begin{equation}
			\E\left[\psi(X, Y; \tilde{f}, \tilde{\alpha}, \theta_h(f))\right]
			= \E\left[\bigl(\alpha_h(X) - \tilde{\alpha}(X)\bigr)
			\bigl(\tilde{f}(X) - f(X)\bigr)\right]
		\end{equation}
		In particular the bias is bounded by
		$\Ptwo{\alpha_h - \tilde{\alpha}} \cdot \Ptwo{\tilde{f} - f}$.
	\end{lemma}
	
	Lemma~\ref{lem:orth} is an identity, not an expansion: because $\theta_h$ is linear
	in $f$, the moment bias equals the product of the two nuisance errors with no
	linearization remainder whatsoever. Neyman orthogonality --- the vanishing of the
	Gateaux derivative of the moment in each nuisance direction at the truth
	\citep{chernozhukov2018dml} --- is the special case obtained by setting one error to
	zero, but the exact form is stronger and is what makes the subsequent analysis
	short. The practical meaning: an error in the fitted learner harms the estimate only
	to the extent that the representer is also misestimated, and vice versa. It is this
	multiplication of errors, not any smoothness of the learner, that buys the
	parametric rate.
	
	The estimator combines the score with cross-fitting. Partition
	$\{1, \dots, n\}$ into $K$ folds $I_1, \dots, I_K$ of approximately equal size
	($K = 5$ is the package default), and for each $k$ let
	$\hat{f}^{(-k)}$ and $\hat{\alpha}^{(-k)}$ be estimates of $f$ and $\alpha_h$
	computed from the observations outside fold $k$ --- the learner exactly as in
	Section~2, the representer as described in
	Section~\ref{sec:rates}. The debiased estimator is:
	\begin{equation}
		\hat{\theta}_h = \frac{1}{n} \sum_{k=1}^{K} \sum_{i \in I_k}
		\left[ w(\mathbf{x}_i)\, \Dh \hat{f}^{(-k)}(\mathbf{x}_i)
		+ \hat{\alpha}^{(-k)}(\mathbf{x}_i)\,
		\bigl(y_i - \hat{f}^{(-k)}(\mathbf{x}_i)\bigr) \right]
		\label{eq:dml}
	\end{equation}
	Cross-fitting decouples the randomness of the nuisance estimates from the
	observations at which the score is evaluated, so that no entropy or Donsker
	condition on the learner's function class is needed \citep{chernozhukov2018dml} ---
	a point of substance here, since the function classes of boosted trees and deep
	networks are exactly the classes for which such conditions fail.
	
	\begin{assumption}
		\label{ass:nuisance}
		For each fold $k$: (a)
		$\Ptwo{\hat{f}^{(-k)} - f} = o_p(1)$ and
		$\Ptwo{\hat{\alpha}^{(-k)} - \alpha_h} = o_p(1)$; (b)
		$\Ptwo{\hat{f}^{(-k)} - f} \cdot \Ptwo{\hat{\alpha}^{(-k)} - \alpha_h}
		= o_p(n^{-1/2})$; (c)
		$\norm{\hat{\alpha}^{(-k)}}_\infty \leq \bar{A} < \infty$ almost surely.
	\end{assumption}
	
	Condition (c) is without loss of generality: $\alpha_h$ obeys the known bound of
	Lemma~\ref{lem:riesz}, so any estimate may be truncated at that bound without
	degrading its $L^2$ error. Conditions (a) and (b) are rate conditions on
	\emph{regression-type} $L^2$ errors only. The sufficient condition for (b) that we
	verify for specific learners in Section~\ref{sec:rates} is that both nuisances
	converge faster than $n^{-1/4}$; the product form also permits trading accuracy in
	one nuisance against the other.
	
	\begin{theorem}[Asymptotic Linearity and Normality]
		\label{thm:main}
		Let Assumptions~\ref{ass:sampling} and~\ref{ass:nuisance} hold, write
		$\psi_i = \psi(\mathbf{x}_i, y_i; f, \alpha_h, \theta_h(f))$, and suppose
		$V_h = \mathrm{Var}(\psi_i) > 0$. Then:
		\begin{equation}
			\sqrt{n}\,\bigl(\hat{\theta}_h - \theta_h(f)\bigr)
			= \frac{1}{\sqrt{n}} \sum_{i=1}^n \psi_i + o_p(1)
			\;\xrightarrow{d}\; N(0, V_h)
		\end{equation}
		Moreover the plug-in variance estimator:
		\begin{equation}
			\hat{V}_h = \frac{1}{n} \sum_{k=1}^{K} \sum_{i \in I_k}
			\left[ w(\mathbf{x}_i)\, \Dh \hat{f}^{(-k)}(\mathbf{x}_i)
			+ \hat{\alpha}^{(-k)}(\mathbf{x}_i)\,
			\bigl(y_i - \hat{f}^{(-k)}(\mathbf{x}_i)\bigr)
			- \hat{\theta}_h \right]^2
		\end{equation}
		satisfies $\hat{V}_h \xrightarrowp V_h$.
	\end{theorem}
	
	No smoothness of the fitted learner appears anywhere in
	Theorem~\ref{thm:main} or its proof. The only property of the finite difference
	operator that is used is that it is bounded on $L^2(P)$ after trimming:
	\begin{equation}
		\Ptwo{w\, \Dh g} \leq
		\frac{1}{h} \left(\frac{p_{\max}}{p_{\min}}\right)^{1/2} \Ptwo{g}
		\label{eq:opbound}
	\end{equation}
	so $L^2$ convergence of the learner transfers to $L^2$ convergence of its window
	effects automatically --- precisely the step that fails for the derivative operator,
	which is unbounded on every $L^p$ space, and the reason the classical route
	requires function-space convergence of the learner's gradient. The condition that
	the earlier version of this theory imposed at this point --- root-$n$ convergence
	of $\fn$ to $f$ in a Sobolev or BV norm --- is not merely weakened but absent: it
	has been replaced by Assumption~\ref{ass:nuisance}(b), a product of $L^2$ rates
	that growing-complexity learners can attain.
	
	The exactness of Lemma~\ref{lem:orth} yields a second dividend beyond weakened
	rates: the score is doubly robust, and one direction of that robustness removes the
	consistency requirement on the learner altogether. Setting
	$\tilde{\alpha} = \alpha_h$ in Lemma~\ref{lem:orth} makes the moment bias vanish
	identically for \emph{every} $\tilde{f}$ --- not approximately, not at first order,
	but as an identity --- and cross-fitting converts this into exact finite-sample
	unbiasedness of $\hat{\theta}_h$. The learner's failures are then paid entirely in
	variance, never in location.
	
	\begin{proposition}[Double Robustness: Inference without Learner Consistency]
		\label{prop:robust}
		Let Assumption~\ref{ass:sampling} hold and suppose that for each fold:
		(a) $\Ptwo{\hat{f}^{(-k)} - g} = o_p(1)$ for some fixed $g \in L^2(\Omega)$,
		\emph{not necessarily equal to $f$};
		(b) $\Ptwo{\hat{\alpha}^{(-k)} - \alpha_h} = o_p(n^{-1/2})$ --- in particular
		$\hat{\alpha}^{(-k)} = \alpha_h$ when the covariate density is known; and
		(c) $\norm{\hat{\alpha}^{(-k)}}_\infty \leq \bar{A} < \infty$ almost surely.
		Define $m_g = w\, \Dh g + \alpha_h\,(f - g)$, which satisfies
		$\E[m_g(X)] = \theta_h(f)$ for every $g$, and:
		\begin{equation}
			V_h(g) = \mathrm{Var}\bigl(m_g(X)\bigr)
			+ \E\bigl[\alpha_h(X)^2\, \sigma^2(X)\bigr]
			\label{eq:vrobust}
		\end{equation}
		If $V_h(g) > 0$, then:
		\begin{equation}
			\sqrt{n}\,\bigl(\hat{\theta}_h - \theta_h(f)\bigr)
			\;\xrightarrow{d}\; N\bigl(0, V_h(g)\bigr),
			\qquad
			\hat{V}_h \xrightarrowp V_h(g)
		\end{equation}
		so the confidence intervals of Corollary~\ref{cor:ci} retain asymptotic
		coverage. When $\hat{\alpha}^{(-k)} = \alpha_h$ exactly,
		$\E[\hat{\theta}_h] = \theta_h(f)$ at every sample size.
	\end{proposition}
	
	The practical reach of Proposition~\ref{prop:robust} is wider than its hypotheses
	may suggest, because condition (a) is the condition that off-the-shelf machine
	learning defaults actually satisfy. A depth-capped forest, an early-stopped
	boosting run, a network of fixed width --- fixed-complexity learners generally ---
	converge not to $f$ but to a pseudo-true limit $g$ within their approximation
	class, so Assumption~\ref{ass:nuisance}(a) fails and Theorem~\ref{thm:main} is
	silent about them. The proposition covers exactly this case whenever the
	representer side is strong: in simulation designs and designed experiments the
	covariate density is known and $\alpha_h$ is computable in closed form; and in the
	binary case of Remark~\ref{rem:binary}, a randomized $X_j$ means a \emph{known}
	propensity score, so the marginal effect of a randomized treatment is validly
	covered for any black-box learner with no consistency requirement whatsoever ---
	the model-agnostic analog of the classical robustness of randomization
	\citep{robins1994}. The mirror-image direction of the double robustness is worth
	recording too: since the bias identity is symmetric, a learner converging at the
	parametric rate --- a correctly specified GLM --- tolerates an arbitrarily poor
	representer estimate, which is how the classical delta-method validity of
	parametric average marginal effects is nested here as the opposite corner of the
	same identity.
	
	Two honest qualifications. First, when $g \neq f$ the estimator is asymptotically
	linear with influence function different from the efficient one; there is no
	conflict with Theorem~\ref{thm:efficiency} below, because such an estimator is
	not regular, and efficiency is restored exactly when $g = f$, at which point
	\eqref{eq:vrobust} reduces to the bound. The noise component of the variance is
	invariant in $g$; misspecification of the learner is paid entirely through the
	dispersion of $m_g$, typically as wider intervals. Second, in observational data
	the density is unknown and condition (b) demands a parametric-rate representer,
	which Riesz regression delivers only under strong structure; the proposition is
	therefore a complement to Theorem~\ref{thm:main}, not a replacement --- the
	theorem splits the burden multiplicatively across the two nuisances, the
	proposition shifts it entirely onto one of them.
	
	\begin{theorem}[Semiparametric Efficiency]
		\label{thm:efficiency}
		Let Assumption~\ref{ass:sampling} hold. In the nonparametric model in which the
		distribution of $(X, Y)$ is restricted only by
		Assumption~\ref{ass:sampling}, the efficient influence function of the
		functional $P \mapsto \theta_h(f_P)$ is
		$\psi(x, y; f, \alpha_h, \theta_h(f))$, and the semiparametric efficiency
		bound is:
		\begin{equation}
			V_h = \mathrm{Var}\bigl(w(X)\, \Dh f(X)\bigr)
			+ \E\bigl[\alpha_h(X)^2\, \sigma^2(X)\bigr]
			\label{eq:effbound}
		\end{equation}
		where $\sigma^2(x) = \mathrm{Var}(Y \mid X = x)$. The estimator
		$\hat{\theta}_h$ of Theorem~\ref{thm:main} attains the bound.
	\end{theorem}
	
	The variance decomposition \eqref{eq:effbound} is the honest accounting of what
	inference on a marginal effect costs. The first term is the heterogeneity of the
	effect itself across the covariate distribution --- irreducible, present even if
	$f$ were known. The second is the price of not knowing $f$: the conditional noise
	amplified by the squared representer, which scales as $h^{-2}$ in the worst case
	through the bound of Lemma~\ref{lem:riesz}. Narrow windows purchase proximity to
	the derivative estimand at the cost of a larger efficiency bound --- the
	bias--variance tradeoff of the problem, located in the estimand rather than hidden
	in an estimator's tuning. No estimator that is regular for $\theta_h$ can do better
	than \eqref{eq:effbound}; the question the earlier version of this theory could not
	answer --- whether the procedure is optimal in any sense --- has the sharpest
	available affirmative answer for the window estimand.
	
	\begin{corollary}[Confidence Intervals]
		\label{cor:ci}
		Under the conditions of Theorem~\ref{thm:main},
		$\hat{\theta}_h \pm z_{1-\gamma/2}\, (\hat{V}_h / n)^{1/2}$ is an
		asymptotically valid $1 - \gamma$ confidence interval for $\theta_h(f)$, and
		the associated $t$-statistic is asymptotically standard normal under
		$H_0 : \theta_h(f) = 0$.
	\end{corollary}
	
	\begin{remark}[Resampling-Free and Resampling-Based Inference]
		\label{rem:bootstrap}
		Corollary~\ref{cor:ci} requires no resampling: the coefficient table of
		Section~2 is produced from a single pass of cross-fitting, with $K$ model fits
		in place of the $B$ refits of the bootstrap --- a reduction of one to two orders
		of magnitude in computation. When a resampling distribution is nonetheless
		wanted (for simultaneous bands across coefficients, for example), asymptotic
		linearity delivers one for free: the multiplier bootstrap that perturbs the
		estimated influence function values $\hat{\psi}_i$ by i.i.d.\ mean-zero
		weights is consistent \citep[Chapter~3.6]{vandervaart1996weak} under the
		conditions of Theorem~\ref{thm:main}, and costs no additional model fits.
	\end{remark}
	
	\begin{remark}[The Model-Refitting Bootstrap]
		\label{rem:refit}
		The refitting bootstrap of Section~2, applied to the uncorrected plug-in
		average, is not covered by these guarantees, for two reasons that the present
		framework makes visible. First, it recenters at the fitted model, so it
		measures variability around $\E[w \Dh \hat{f}]$ rather than around
		$\theta_h(f)$: the regularization bias
		$\E[\alpha_h (\hat{f} - f)]$ --- first order for slowly converging learners
		--- is invisible to it. This is the mechanism behind the collapse of random
		forest coverage documented in Section~5. Second, evaluating each refitted model
		at the original sample holds the empirical covariate distribution fixed and
		thereby omits the first variance component of \eqref{eq:effbound} from the
		bootstrap distribution. The refitting bootstrap retains a legitimate role as a
		diagnostic for specification and hyperparameter-selection variability, which
		the asymptotic theory conditions away; it should be read as that diagnostic,
		not as the inferential procedure.
	\end{remark}
	
	\subsection{What Learner Gradients Estimate}
	\label{sec:interpretation}
	
	The inference theory of Section~\ref{sec:inference} used no property of the fitted
	learner beyond an $L^2$ rate. The question it leaves open is interpretive: the
	first component of the score averages the finite differences of the fitted model,
	and one wants to know what that object measures about the model --- and when the
	model's gradient information converges to the population's. This subsection answers
	both questions in the function spaces the learners naturally inhabit. Four lemmas
	establish that gradients exist and converge in the Sobolev space of smooth
	estimators and the BV space of piecewise constant ones; a pairing identity
	(Theorem~\ref{thm:pairing}) then expresses the window effect of any BV estimator
	exactly as its jump measure integrated against a locally averaged covariate
	density, which is, to our knowledge, the first characterization of the population
	quantity that the marginal effect of a tree-based model measures. All proofs are in
	the Appendix.
	
	The parallel structure of Lemmas~\ref{lem:sobolev_existence}
	through~\ref{lem:bv_consistency} reflects the parallel structure of the Sobolev
	and BV frameworks and is the formal expression of the analogy between smooth and
	piecewise constant estimators. None of these results is needed for the validity of
	Theorem~\ref{thm:main}; they characterize the regime in which the \emph{uncorrected}
	plug-in average is already consistent, and hence in which the correction term of
	\eqref{eq:dml} is a refinement rather than a repair.
	
	\begin{lemma}[Derivative Existence in Sobolev Spaces]
		\label{lem:sobolev_existence}
		Let $\fn \in \Wtwop$ for some $p > 1$. Then $\nabla \fn \in W^{1,p}(\Omega)^d$
		exists as a well defined $L^p$ function and satisfies:
		\begin{equation}
			\norm{\nabla \fn}_{L^p} \leq \norm{\fn}_{W^{1,p}} \leq \norm{\fn}_{W^{2,p}}
		\end{equation}
	\end{lemma}
	
	Lemma~\ref{lem:sobolev_existence} is essentially definitional---membership in
	$\Wtwop$ implies by definition that the first and second weak derivatives exist and
	are in $L^p$. The content of the lemma is the norm bound, which establishes that
	the gradient operator $\nabla: W^{2,p} \to W^{1,p}$ is bounded and linear. This
	boundedness is what makes the subsequent consistency argument work---a bounded linear
	operator preserves convergence, so once we have convergence of $\fn$ in the right
	space, convergence of $\nabla \fn$ follows automatically. The reason we require
	$\fn \in \Wtwop$ rather than just $\Wonep$ is that we need one additional degree of
	differentiability to apply the compactness argument in
	Lemma~\ref{lem:sobolev_consistency}. The relevant theory of Sobolev spaces and their
	norms is developed in \citet[Chapter~5]{evans2010pde}.
	
	\begin{lemma}[Derivative Consistency in Sobolev Spaces]
		\label{lem:sobolev_consistency}
		Let $f \in \Wonep$ and suppose $\fn$ is well-behaved in the Sobolev sense. Then:
		\begin{equation}
			\norm{\nabla \fn - \nabla f}_{L^p} \xrightarrowp 0
		\end{equation}
	\end{lemma}
	
	Lemma~\ref{lem:sobolev_consistency} is the harder result and the one that does the
	real work in the Sobolev setting. The difficulty is that $L^p$ consistency of $\fn$
	alone does not imply $L^p$ consistency of $\nabla \fn$---differentiation is an
	unbounded operator on $L^p$ and limits and derivatives cannot be exchanged without
	additional structure. The uniform $\Wtwop$ boundedness condition provides that
	structure. The proof applies the Rellich--Kondrachov compact embedding theorem, which
	states that a uniformly bounded sequence in $\Wtwop$ has a subsequence converging
	strongly in $\Wonep$ \citep[Theorem~5,~Section~5.7]{evans2010pde}. Combined with
	the $L^p$ consistency condition identifying the limit as $f$, this gives strong
	convergence of $\fn$ to $f$ in $\Wonep$, from which $L^p$ convergence of the
	gradients follows immediately from the boundedness of $\nabla$ established in
	Lemma~\ref{lem:sobolev_existence}. The key insight is that the uniform boundedness
	condition ensures the estimator sequence lives in a compact subset of $\Wonep$,
	preventing the gradients from escaping to infinity even as the estimator fits the
	data more closely.
	
	\begin{lemma}[Derivative Existence in BV]
		\label{lem:bv_existence}
		Let $\fn \in \BV$. Then the distributional gradient $D\fn$ exists as a well defined
		vector valued Radon measure in $\cM$ with:
		\begin{equation}
			|D\fn|(\Omega) < \infty
		\end{equation}
		For piecewise constant estimators with leaf partition $\{R_j\}$ and leaf means
		$\{\bar{y}_j\}$, the distributional gradient takes the explicit form:
		\begin{equation}
			D\fn = \sum_{j \sim k} (\bar{y}_k - \bar{y}_j) \cdot \nu_{jk} \cdot
			\mathcal{H}^{d-1}\lfloor_{\partial R_j \cap \partial R_k}
		\end{equation}
		where $\nu_{jk}$ is the unit normal between adjacent leaves and $\mathcal{H}^{d-1}$
		is the $(d-1)$-dimensional Hausdorff measure on leaf boundaries.
	\end{lemma}
	
	Lemma~\ref{lem:bv_existence} is the BV analog of Lemma~\ref{lem:sobolev_existence}
	and is similarly structural in nature. The first part---that $BV$ membership implies
	the distributional gradient exists as a Radon measure---follows from the definition
	of $BV$ as developed in \citet[Chapter~3]{ambrosio2000bv}. The content of the lemma
	is the explicit formula for piecewise constant estimators, which makes concrete what
	the distributional gradient actually looks like for a tree based model. The formula
	says that the gradient is not a function at all but a measure concentrated on the
	boundaries between leaves, with mass proportional to the jump in response values
	across the boundary weighted by the boundary surface area. This is the precise sense
	in which a random forest or boosted tree has a gradient---not a pointwise derivative,
	which does not exist, but a distributional object that captures where and how much
	the estimator changes. This explicit representation is also what makes the total
	variation $|D\fn|(\Omega)$ computable in terms of observable quantities---the leaf
	means and the partition geometry---which is essential for verifying the well-behaved
	condition in the theorems below.
	
	\begin{lemma}[Derivative Consistency in BV]
		\label{lem:bv_consistency}
		Let $f \in \Wonep$ for some $p > 1$, so that $f \in \BV$ with
		$Df = \nabla f \cdot \mathcal{L}^d$. Suppose $\fn$ is well-behaved in the BV sense.
		Then:
		\begin{equation}
			D\fn \wstar Df \quad \text{in } \cM
		\end{equation}
		If additionally $\E[|D\fn|(\Omega)] \to |Df|(\Omega)$, then convergence is strict.
	\end{lemma}
	
	Lemma~\ref{lem:bv_consistency} is the hardest of the four function space lemmas and
	the BV analog of Lemma~\ref{lem:sobolev_consistency}. The convergence is necessarily
	weaker than in the Sobolev case---weak-$*$ convergence of measures rather than strong
	$L^p$ convergence---because the $BV$ space is not reflexive and the stronger norm
	convergence is not available in general. Weak-$*$ convergence means that for any
	continuous test function $\phi$:
	\begin{equation}
		\int_\Omega \phi \, d(D\fn) \to \int_\Omega \phi \, d(Df)
	\end{equation}
	which is the appropriate notion of convergence for the distributional gradients of
	piecewise constant estimators. In the context of average marginal effects, this means
	that integrated gradient information---which is exactly what the AME computes---
	converges to the correct population quantity. The proof applies the BV compactness
	theorem of \citet[Theorem~3.23]{ambrosio2000bv}, which is the analog of
	Rellich--Kondrachov in the BV setting: a uniformly bounded sequence in $BV$ has a
	subsequence converging in $L^1$, and the $L^1$ consistency condition identifies the
	limit as $f$. The strict convergence upgrade follows from the lower semicontinuity of
	the total variation functional and is the stronger result one would want for
	applications requiring convergence of the total size of gradient information rather
	than just its distribution.
	
	The parallel between Lemmas~\ref{lem:sobolev_consistency}
	and~\ref{lem:bv_consistency} is worth noting explicitly. In both cases the argument
	has the same structure: uniform boundedness in a strong space gives compactness,
	consistency in a weak space identifies the limit, and the combination gives
	convergence of the derivative in the appropriate topology. The difference is in the
	topology of the conclusion---strong $L^p$ convergence in the Sobolev case, weak-$*$
	convergence of measures in the BV case---and this difference is fundamental rather
	than technical. It reflects the fact that smooth estimators carry gradient
	information in a stronger sense than piecewise constant estimators, which is itself a
	reflection of the different inductive biases of the estimator classes.
	
	The pairing identity is the capstone of the interpretation theory. Recall from
	Lemma~\ref{lem:bv_existence} that the distributional gradient of a piecewise
	constant estimator is a measure concentrated on leaf boundaries. The window effect
	of such an estimator integrates that measure against a smoothed density: let
	$K_h \star q$ denote the one-dimensional moving average of $q = wp$ (extended by
	zero) along coordinate $j$:
	\begin{equation}
		(K_h \star q)(u) = \frac{1}{2h} \int_{-h}^{h} q(u + s\, e_j)\, ds
	\end{equation}
	
	\begin{theorem}[Pairing Identity]
		\label{thm:pairing}
		Let Assumption~\ref{ass:sampling} hold, let $w$ be a trimming weight with
		$\operatorname{supp}(w) \subseteq \Omega_{j,h}$, and let $g \in \BV$
		(identified with its precise representative). Then:
		\begin{equation}
			\E\left[w(X)\, \Dh g(X)\right]
			= \int_\Omega (K_h \star q)(u)\; dD_j g(u)
			\label{eq:pairing}
		\end{equation}
		where $D_j g$ is the $j$-th component of the distributional gradient of $g$.
		In particular, for a piecewise constant $g$ with leaf partition $\{R_l\}$ and
		leaf values $\{\bar{y}_l\}$, by Lemma~\ref{lem:bv_existence}:
		\begin{equation}
			\E\left[w(X)\, \Dh g(X)\right]
			= \sum_{l \sim m} (\bar{y}_m - \bar{y}_l)\, (\nu_{lm})_j
			\int_{\partial R_l \cap \partial R_m} (K_h \star q)\;
			d\mathcal{H}^{d-1}
			\label{eq:pairing_tree}
		\end{equation}
	\end{theorem}
	
	\begin{corollary}
		\label{cor:pairing_limit}
		If in addition $q$ extends to a continuous function of compact support (for
		instance $w$ continuous and $p$ continuous on $\Omega$), then for fixed
		$g \in \BV$:
		\begin{equation}
			\lim_{h \downarrow 0} \int_\Omega (K_h \star q)\; dD_j g
			= \int_\Omega q\; dD_j g
		\end{equation}
		and when $g \in W^{1,1}(\Omega)$ the right side equals
		$\E[w(X)\, \partial_j g(X)]$, consistent with
		Proposition~\ref{prop:hlimit}.
	\end{corollary}
	
	Theorem~\ref{thm:pairing} settles a question that the interpretability literature
	has left informal: what does the average marginal effect of a random forest
	\emph{measure}? The answer given by \eqref{eq:pairing_tree} is exact and
	geometric: it is the total jump of the fitted surface across leaf boundaries in
	direction $j$, with each boundary segment weighted by its share of the covariate
	mass in an $h$-window --- the places where the model changes its prediction,
	weighted by how many observations live near those places. The identity reconciles
	the two readings of a marginal effect that Section~2 could only juxtapose. Read
	with $g = \hat{f}$, \eqref{eq:pairing} is a descriptive statement about the fitted
	model, available in finite samples with no assumptions at all. Read through
	Theorem~\ref{thm:main}, the same plug-in quantity, corrected by the orthogonal
	score, estimates the population functional $\theta_h(f)$ at the parametric rate.
	The descriptive and the inferential interpretations are the two sides of
	\eqref{eq:dml}, and the pairing identity is what makes the first of them precise.
	
	\subsection{Gradient Convergence of Specific Learners}
	\label{sec:learners}
	
	The theorems of this subsection verify the well-behaved conditions for three
	estimator classes and conclude, via the lemma pairs above, that the learner's own
	gradient information converges to the population's. This is the strongest
	interpretability property a learner can have --- the fitted model's local behavior,
	not merely its window averages, tracks the truth --- and it implies consistency of
	the uncorrected plug-in average for the small-$h$ estimand. We emphasize the
	division of labor: these results calibrate how much interpretive weight the fitted
	gradients themselves can bear, while the inference of Section~\ref{sec:inference}
	requires none of them. In each case the argument has two components: verifying that
	the estimator satisfies the well-behaved condition, and citing the appropriate
	lemma pair. The first component is the estimator-specific contribution; the second
	is the functional analytic machinery established above.
	
	\begin{theorem}[Gradient Consistency for Neural Networks]
		\label{thm:neural_nets}
		Let $\fn$ be a neural network estimator with smooth activations $\sigma \in C^2$,
		bounded weight norms $\sup_n \prod_\ell \norm{W_\ell^{(n)}} < \infty$, and
		$\norm{\fn - f}_{L^p} \xrightarrowp 0$ for $p > d$. Then $\fn$ is well-behaved in
		the Sobolev sense and:
		\begin{equation}
			\norm{\nabla \fn - \nabla f}_{L^p} \xrightarrowp 0
		\end{equation}
	\end{theorem}
	
	The conditions of Theorem~\ref{thm:neural_nets} are natural and verifiable. Smooth
	activations --- sigmoid, tanh, GELU --- are standard in practice and ensure the network
	function is twice differentiable. The bounded weight norm condition formalizes the
	effect of $L^2$ regularization during training, which prevents the network weights from
	growing without bound. The $L^p$ consistency condition is the standard requirement that
	the network is learning the true function. \citet{hornik1990universal} establish that
	multilayer feedforward networks with smooth activation functions are capable of
	approximating an arbitrary function and its derivatives to arbitrary accuracy, providing
	the foundation for the consistency condition and establishing that the function class is
	rich enough to contain the estimators we consider. The condition $p > d$ is required by
	the pointwise differentiability theorem of \citet[Theorem~6.5]{evans2015measure} and
	is satisfied for all $p$ larger than the covariate dimension --- a mild regularity
	condition that is implicit in the classical Sobolev embedding theory.
	
	\begin{proof}[Proof sketch]
		The chain rule applied through the network layers gives $\fn \in \Wtwop$ with:
		\begin{equation}
			\norm{\fn}_{W^{2,p}} \lesssim \prod_\ell \norm{W_\ell^{(n)}}^2 \cdot
			\norm{\sigma''}_\infty
		\end{equation}
		which is uniformly bounded by the weight norm assumption. The $L^p$ consistency
		condition and Lemma~\ref{lem:sobolev_consistency} then give the result. Full proof
		in Appendix~\ref{app:neural_nets}.
	\end{proof}
	
	\begin{remark}
		Theorem~\ref{thm:neural_nets} is stated for smooth activations. ReLU networks
		require a modified argument because ReLU is not $C^2$ and the chain rule calculation
		does not apply directly. Extension to ReLU networks is possible through a
		distributional argument but complicates the presentation and is omitted here.
	\end{remark}
	
	\begin{theorem}[Gradient Consistency for Gradient Boosted Models]
		\label{thm:boosting}
		Let $\fn = \hat{F}_m = (I - (I-S)^{m+1})Y$ be an $L_2$Boost estimator with
		symmetric linear learner $S$ satisfying the eigenvalue condition $0 < \lambda_k \leq
		1$ for all $k$, and with base learners $h_m$ that are either smooth ($h_m \in
		W^{2,p}$) or tree stumps. Suppose $\norm{\fn - f}_{L^1} \xrightarrowp 0$. Then:
		\begin{enumerate}
			\item[(i)] If base learners are smooth, $\fn$ is well-behaved in the Sobolev
			sense and $\norm{\nabla \fn - \nabla f}_{L^p} \xrightarrowp 0$.
			\item[(ii)] If base learners are tree stumps, $\fn$ is well-behaved in the
			BV sense and $D\fn \wstar Df$ in $\cM$.
		\end{enumerate}
	\end{theorem}
	
	Theorem~\ref{thm:boosting} covers both cases that arise in practice and illustrates
	how the function space framework adapts to the base learner. XGBoost and
	LightGBM --- the implementations used in Section~3 --- use shallow trees as base
	learners and therefore fall under case~(ii). The two cases in the theorem are not
	merely technical alternatives --- they reflect a genuine difference in what the
	gradient of the estimator is. For smooth base learners the gradient is an ordinary
	function and $L^p$ convergence is available. For tree base learners the gradient is a
	measure and only weak-$*$ convergence can be established. The average marginal effects
	are well defined and consistent in both cases, but the sense in which they are
	consistent differs.
	
	The eigenvalue condition $0 < \lambda_k \leq 1$ is the key regularity condition on
	the base learner. \citet{buhlmann2003boosting} impose this condition explicitly in
	their Theorem~1 and establish that under it the bias decays exponentially fast while
	the variance increases only by exponentially diminishing amounts --- the exponential
	bias--variance trade-off that makes $L_2$Boost analytically tractable. The condition
	holds for smoothing spline learners and, generically, for tree stump learners whose
	splits have positive but incomplete explanatory power, which is the typical case in
	practice. Learners with any eigenvalue equal to zero are degenerate in the sense of
	Corollary~1 of \citet{buhlmann2003boosting}: $\mathcal{B}_m \equiv S$ for all $m$ and
	boosting does not improve on the base learner. The functional gradient descent
	perspective of \citet{mason2000boosting} provides the conceptual framework for
	understanding the boosting path as gradient descent in function space.
	
	\begin{proof}[Proof sketch]
		For smooth base learners the $\Wtwop$ bound follows from linearity of the estimator:
		\begin{equation}
			\norm{\fn}_{W^{2,p}} \leq \sum_m \eta_m \norm{h_m}_{W^{2,p}} < \infty
		\end{equation}
		For tree stump base learners the BV bound uses the eigenvalue structure of
		$L_2$Boost. By Proposition~1 of \citet{buhlmann2003boosting}, the residual at step
		$m$ is $u_m = (I-S)^m Y$. By Proposition~2, under the eigenvalue condition $0 <
		\lambda_k \leq 1$, Theorem~1 of \citet{buhlmann2003boosting} gives exponential decay
		of residuals:
		\begin{equation}
			\norm{u_m}_{L^2} \leq C\rho^m, \quad \rho = \max_k(1-\lambda_k) < 1
		\end{equation}
		Each stump's BV norm is controlled by its residual magnitude via Lemma~\ref{lem:bv_existence},
		and the total BV norm of the ensemble satisfies:
		\begin{equation}
			\mathbb{E}[|D\fn|(\Omega)] \leq C\sum_{m=1}^\infty \rho^m = \frac{C\rho}{1-\rho} < \infty
		\end{equation}
		The geometric series converges because $\rho < 1$, giving uniform BV boundedness
		without any step size condition. Lemma~\ref{lem:bv_consistency} then gives
		$D\fn \xrightarrow{w*} Df$. Full proof in Appendix~\ref{app:boosting}.
	\end{proof}
	
	\begin{theorem}[Gradient Consistency for Mondrian Forests]
		\label{thm:forests}
		Let $\fn$ be a Mondrian random forest estimator with budget $\lambda_n \asymp
		n^{1/(d+2)}$, $B \gg n^{2(d-\alpha)/(d+2)}$ trees where $\alpha = 1 - d/p$, $f$
		Lipschitz on $[0,1]^d$, and $f \in \Wonep$ for $p > d$. Then $\fn$ is well-behaved
		in the BV sense and:
		\begin{equation}
			D\fn \wstar Df \quad \text{in } \cM
		\end{equation}
	\end{theorem}
	
	Theorem~\ref{thm:forests} is the most technically demanding result in the paper.
	Random forests are genuinely different from neural networks and boosted models in a
	way that makes gradient consistency harder to establish --- the partition structure is
	random and the response averages within leaves depend on the data in a complex way.
	The restriction to Mondrian forests, where splits are made independently of the data,
	is the key simplification that makes the proof tractable. For Mondrian forests the
	partition geometry and the response variation can be controlled separately, whereas for
	Breiman forests with data-dependent splits the two are coupled in a way that resists
	clean analysis. The $L^2$ consistency result under the budget $\lambda_n \asymp
	n^{1/(d+2)}$ follows from \citet{mourtada2017mondrian}, who establish the minimax rate
	$O(n^{-2/(d+2)})$ for Lipschitz regression functions, and $L^1$ consistency follows
	by Cauchy--Schwarz on the bounded domain $\Omega$.
	
	The budget condition $\lambda_n \asymp n^{1/(d+2)}$ and the tree count condition $B
	\gg n^{2(d-\alpha)/(d+2)}$ are natural and interpretable. The budget determines the
	approximation rate --- balancing the bias from coarse partitions against the variance
	from fine partitions at the minimax optimal rate. The tree count condition ensures that
	the forest average has sufficient cancellation across trees to control the total
	variation; it is mild in practice, satisfied by standard forest sizes for moderate
	dimension and sample size.
	
	\begin{proof}[Proof sketch]
		Each tree is piecewise constant and therefore in $BV$ by
		Lemma~\ref{lem:bv_existence}. The key difficulty is that individual tree BV norms
		grow with $n$, so the triangle inequality applied to the forest average is too crude.
		Instead we work with the second moment of the total variation. Writing $D\fn =
		\frac{1}{B}\sum_b DT_b$:
		\begin{equation}
			\mathbb{E}[|D\fn|(\Omega)] \leq \sqrt{\mathbb{E}[|D\fn|(\Omega)^2]}
			= \sqrt{\frac{1}{B^2}\mathbb{E}\!\left[\left|\sum_b DT_b\right|(\Omega)^2\right]}
		\end{equation}
		The cross terms between trees $b \neq b'$ vanish by independence of the Mondrian
		partitions and the zero-mean condition on the jumps. The diagonal terms are bounded
		using the cell diameter result $\mathbb{E}[D_{\lambda_n}(x)^2] \leq 4d/\lambda_n^2$
		of \citet{roy2009mondrian} and the split count bound $\mathbb{E}[K_{\lambda_n}]
		\leq (e(\lambda_n+1))^d$ of \citet{mourtada2017mondrian}, giving:
		\begin{equation}
			\mathbb{E}[|D\fn|(\Omega)^2]
			= \frac{1}{B}\,\mathbb{E}[|DT_1|(\Omega)^2]
			\leq \frac{1}{B}\,O(n^{2(d-\alpha)/(d+2)})
		\end{equation}
		Under the condition $B \gg n^{2(d-\alpha)/(d+2)}$ this is uniformly bounded.
		Lemma~\ref{lem:bv_consistency} then gives $D\fn \xrightarrow{w*} Df$ in $\cM$.
		Full proof in Appendix~\ref{app:forests}.
	\end{proof}
	
	\begin{remark}
		Extension of Theorem~\ref{thm:forests} to Breiman forests with data-dependent splits
		requires additional control over the interaction between partition geometry and
		response variation. The data-dependent splitting rule couples these two quantities in
		a way that makes the cross-term cancellation argument unavailable --- splits are
		chosen precisely to maximize response differences across boundaries, which is the
		worst case for BV control. This extension is an important direction for future work
		and is discussed further in Section~\ref{sec:conclusion}.
	\end{remark}
	
	\subsection{Nuisance Estimation and Rates}
	\label{sec:rates}
	
	Assumption~\ref{ass:nuisance} asks for two regression-type $L^2$ rates whose
	product is $o_p(n^{-1/2})$; the benchmark sufficient condition is that each
	nuisance converge faster than $n^{-1/4}$. This subsection records when the
	learners of Section~2 meet the benchmark for $\hat{f}$, and how the representer
	$\alpha_h$ is estimated without ever estimating a density.
	
	For the regression function, the required rates are exactly the $L^2$ risk bounds
	of the modern nonparametric literature, and a single smoothness threshold
	organizes all three learner classes. For deep neural networks with ReLU
	activations, \citet{schmidthieber2020} establishes the rate
	$n^{-s/(2s + d)}$ up to logarithmic factors for $s$-H\"{o}lder regression
	functions (with the effective $d$ reduced under compositional structure), so the
	$n^{-1/4}$ benchmark is met precisely when:
	\begin{equation}
		\frac{s}{2s + d} > \frac{1}{4}
		\quad \Longleftrightarrow \quad
		s > \frac{d}{2}
		\label{eq:threshold}
	\end{equation}
	For Mondrian forests, the minimax-optimal rates of \citet{mourtada2017mondrian}
	take the same form $n^{-s/(2s + d)}$ for H\"{o}lder smoothness $s \leq 2$, so the
	benchmark is again \eqref{eq:threshold}: at Lipschitz smoothness $s = 1$ this
	requires $d < 2$, and exploiting forest smoothing up to $s = 2$ accommodates
	$d \leq 3$. For $L_2$Boost, \citet{buhlmann2003boosting} establish minimax rates
	for smoothing-spline base learners with the boosting iteration acting as the
	regularization parameter, and the benchmark corresponds to smoothness exceeding
	one half in their univariate and additive settings --- with additive structure,
	as usual, breaking the dependence on the ambient dimension. The reading we intend
	is not that these conditions are innocuous --- \eqref{eq:threshold} is a real
	restriction, and in high dimensions it binds --- but that the currency in which
	the conditions are denominated has changed: what inference requires of the learner
	is an $L^2$ error rate of the kind the learning-theory literature actually
	proves, rather than root-$n$ convergence in a Sobolev norm, which no
	growing-complexity learner satisfies. For Breiman forests with data-dependent
	splits, $L^2$ consistency is available under structural assumptions
	\citep{scornet2015forests}, and polynomial $L^2$ rates for the original
	algorithm --- sample CART splits and column subsampling included --- have been
	established under sufficient-impurity-decrease conditions by
	\citet{chi2022forests}; the resulting exponents can fall short of the
	$n^{-1/4}$ benchmark, but the product form of
	Assumption~\ref{ass:nuisance}(b) permits pairing a slow forest rate with a fast
	representer estimate, and Proposition~\ref{prop:robust} covers the
	known-representer case with no forest rate at all. The open question relevant to
	this paper is therefore an $L^2$ rate question, a substantially more tractable
	target than the gradient-consistency question that
	Section~\ref{sec:learners} leaves open for that class.
	
	The representer is estimated by Riesz regression
	\citep{chernozhukov2022riesz}, which exploits the same linearity that produced
	Lemma~\ref{lem:orth}. By Lemma~\ref{lem:riesz}, for any candidate
	$\alpha \in L^2(P)$:
	\begin{equation}
		\E\left[\alpha(X)^2\right] - 2\,\E\left[w(X)\, \Dh \alpha(X)\right]
		= \Ptwo{\alpha - \alpha_h}^2 - \Ptwo{\alpha_h}^2
		\label{eq:rieszloss}
	\end{equation}
	so the population loss on the left is minimized uniquely at $\alpha_h$, and its
	sample analog:
	\begin{equation}
		\hat{L}(\alpha) = \frac{1}{n} \sum_{i=1}^{n} \left[
		\alpha(\mathbf{x}_i)^2
		- 2\, w(\mathbf{x}_i)\,
		\frac{\alpha(\mathbf{x}_i + h e_j) - \alpha(\mathbf{x}_i - h e_j)}{2h}
		\right]
	\end{equation}
	is computable directly: minimizing $\hat{L}$ over a class of functions ---
	penalized linear sieves, neural networks, or forests --- requires only the ability
	to evaluate candidates at shifted points, and at no stage is the density $p$, let
	alone a derivative of it, estimated. \citet{chernozhukov2022riesz} provide
	$L^2$ convergence guarantees for such estimators over exactly these classes, with
	rates governed by the complexity of the class and the approximability of
	$\alpha_h$. Since $\alpha_h$ obeys the explicit bound of Lemma~\ref{lem:riesz},
	the minimizer is truncated at that bound, which enforces
	Assumption~\ref{ass:nuisance}(c) by construction without affecting the $L^2$
	error. Two further remarks. First, the product form of
	Assumption~\ref{ass:nuisance}(b) permits asymmetry: a learner with a fast rate
	for $f$ tolerates a slow Riesz estimate, and conversely. Second, in designed or
	simulated data where the covariate density is known --- as in the Monte Carlo
	study of Section~5 --- $\alpha_h$ is available in closed form and the representer
	error is identically zero; Proposition~\ref{prop:robust} then applies, and valid
	inference requires no rate --- indeed no consistency --- from the learner at all.
	
	\section{Empirical Examples}
	
	Across the three empirical applications we fit four model classes: logistic or linear regression as the parametric baseline, a random forest, XGBoost, and a feedforward neural network. The random forest uses 200 trees with a maximum depth of 8. XGBoost uses 200 boosting rounds with a maximum depth of 4 and a learning rate of 0.05. The neural network architecture and training schedule depend on the outcome type. For binary classification, the network is a feedforward architecture with two hidden layers of 64 and 32 units trained with the Adam optimizer at a learning rate of $10^{-3}$, with early stopping on a 10\% validation split with patience of 20 epochs. For continuous regression, the network is deeper --- four hidden layers of 256, 128, 64, and 32 units --- trained at a learning rate of $3 \times 10^{-4}$ with patience of 50 epochs, reflecting the larger capacity required to fit the continuous outcome surface and the longer training horizon needed for the optimization to converge. For regression outcomes, the network includes the input normalization layer and target standardization wrapper described in Section~2 to ensure scale-invariant gradient updates and predictions on the natural outcome scale. Estimation follows the debiased cross-fitted procedure of Section~2.4 with $K=5$ folds: within each fold the learner and the Riesz representer are fit on the other four folds and evaluated on the held-out fold, and standard errors come from the sample variance of the orthogonal score. No learner is ever differenced at an observation it was trained on. Every model class is refit from scratch on each fold, including the neural networks, whose early stopping schedule is applied within the fold. The fit statistics reported at the foot of each table are computed from a single full-sample fit and are descriptive of the learner rather than of the estimator; they are not used in the estimation. Where applicable, predicted probabilities for classification models are extracted from \texttt{predict\_proba(X)[:,1]} as discussed in Section~2.

The window estimand requires two settings to be fixed before estimation, and because $\theta_{j,h}$ is defined by them they are reported rather than treated as tuning. The step size is the default $h_j = \max(10^{-4},\, 0.05\hat{\sigma}_j)$, floored at $0.5$ for integer-valued features. The floor matters for count variables such as bedrooms, bathrooms, years of education and months delayed: a count has no mass strictly between its levels, so a window narrower than one unit has an empty interior, and for a piecewise constant learner the difference across it vanishes except where the window happens to straddle a split. At $h_j = 0.5$ the window is exactly the one-unit contrast, which is also the quantity of substantive interest for a count. The second setting is the trimming weight, which gives weight zero to observations lying within $h_j$ of the boundary of the observed support of $x_j$, where $x \pm h_j e_j$ would require the learner to extrapolate. For most features this removes a negligible share of the sample, but for variables with mass sitting exactly on a boundary it does not: capital gains and capital losses are zero for over ninety percent of the Adult sample, and zero is the minimum of their observed support, so those AMEs are effects on the subpopulation with nonzero gains or losses rather than on the full sample. Table~\ref{tab:window_settings} reports $h_j$ and the trimmed fraction for every feature in the three applications, and the discussion below flags the cases where the trimming is large enough to change the interpretation.
	
		\begin{table*}[htbp]
		\centering
		\small
		\begin{threeparttable}
		\caption{Window settings for the empirical applications. Step size $h_j$ and the fraction of observations receiving trimming weight zero.}
		\label{tab:window_settings}
		\begin{tabular}{llrr}
			\toprule
			Dataset & Feature & $h_j$ & Trimmed \\
			\midrule
			UCI Adult Income & age & 0.6855 & 1.3\% \\
			& female & --- & 0.0\% \\
			& education\_num & 0.5 & 1.4\% \\
			& hours\_per\_week & 0.6196 & 0.3\% \\
			& government & --- & 0.0\% \\
			& self\_employed & --- & 0.0\% \\
			& white\_collar & --- & 0.0\% \\
			& blue\_collar & --- & 0.0\% \\
			& service & --- & 0.0\% \\
			& capital\_gain & 373 & 92.3\% \\
			& capital\_loss & 20.15 & 95.3\% \\
			& married & --- & 0.0\% \\
			\midrule
			UCI Credit Card Default & age & 0.5 & 0.2\% \\
			& female & --- & 0.0\% \\
			& education & 0.5 & 36.8\% \\
			& married & --- & 0.0\% \\
			& avg\_pay\_status & 0.04911 & 7.0\% \\
			& months\_delayed & 0.5 & 70.9\% \\
			& avg\_bill\_amt & 3,163 & 0.0\% \\
			& avg\_pay\_amt & 507 & 11.7\% \\
			& pay\_ratio & 0.01839 & 18.4\% \\
			\midrule
			Ames Housing & GrLivArea & 26.27 & 0.1\% \\
			& BedroomAbvGr & 0.5 & 0.5\% \\
			& FullBath & 0.5 & 2.9\% \\
			& HalfBath & 0.5 & 63.4\% \\
			& OverallQual & 0.5 & 1.4\% \\
			& OverallCond & 0.5 & 1.6\% \\
			& YearBuilt & 1.51 & 1.4\% \\
			& YearRemodAdd & 1.032 & 14.5\% \\
			& LotArea & 499 & 1.2\% \\
			& GarageArea & 10.69 & 5.6\% \\
			& nbhd\_high & --- & 0.0\% \\
			& nbhd\_mid & --- & 0.0\% \\
			\bottomrule
			\end{tabular}
			\begin{tablenotes}
			\footnotesize
			\item \textit{Notes:} $h_j$ is the default adaptive step size, $\max(10^{-4},\, 0.05\hat{\sigma}_j)$, floored at $0.5$ for integer-valued features so that a count is contrasted over a whole unit. Features declared binary use the level contrast $f(x \mid x_j=1) - f(x \mid x_j=0)$, for which $h_j$ is undefined and no trimming applies. The trimmed fraction is the share of observations lying within $h_j$ of the boundary of the observed support of $x_j$, which receive weight zero; for those features the estimand is the window effect on the untrimmed subpopulation.
			\end{tablenotes}
		\end{threeparttable}
		\end{table*}
	
	\subsection{UCI Adult Income}
	
	The first dataset is the UCI Adult Income dataset \citep{UCI, kohavi1996}, a widely used benchmark for binary classification in the social sciences. We combine the training and test sets for a total of 48,842 observations. The outcome is whether an individual earns more than \$50,000 per year. Features are organized into three panels: demographics (age, sex, years of education), work characteristics (hours worked per week, occupation group, work class), and financial variables (capital gains, capital losses, marital status). The specification search considers logistic regression and XGBoost across four specifications, allowing the method to be evaluated both where the parametric model is approximately correctly specified and where the black-box learner is substantially better fitted. Years of education takes the floored step size $h_j = 0.5$, so its effect is the one-year contrast; age and hours per week have enough spread that the adaptive rule already exceeds the floor, at $h_j = 0.686$ and $0.620$. Trimming is negligible for every feature except capital gains and capital losses, which lose 92.3\% and 95.3\% of the sample for the reason discussed below.
	
	\begin{table*}[htbp]
		\centering
		\small
		
		\resizebox{\textwidth}{!}{%
		\begin{threeparttable}
			\caption{UCI Adult Income: Specification Search. Average marginal effects on P(Income $>$ \$50k).}
			\label{tab:adult_spec_search}
			
			\begin{tabular}{lrrrrrrrr}
				\toprule
				& \multicolumn{4}{c}{Logistic} & \multicolumn{4}{c}{XGBoost} \\
				\cmidrule(lr){2-5} \cmidrule(lr){6-9}
				& Spec A & A+B & A+C & A+B+C & Spec A & A+B & A+C & A+B+C \\
				\midrule
				\multicolumn{9}{l}{\textit{Panel A: Demographics}} \\
				\quad age & 0.0084$^{***}$ & 0.0081$^{***}$ & 0.0041$^{***}$ & 0.0043$^{***}$ & 0.0059$^{***}$ & 0.0059$^{***}$ & 0.0031$^{***}$ & 0.0033$^{***}$ \\
				& (0.0001) & (0.0001) & (0.0001) & (0.0001) & (0.0001) & (0.0001) & (0.0001) & (0.0001) \\
				\quad female & -0.1749$^{***}$ & -0.1670$^{***}$ & -0.0122$^{***}$ & -0.0153$^{**}$ & -0.1604$^{***}$ & -0.1507$^{***}$ & -0.0089$^{**}$ & -0.0128$^{**}$ \\
				& (0.0034) & (0.0045) & (0.0046) & (0.0061) & (0.0033) & (0.0041) & (0.0042) & (0.0056) \\
				\quad education\_num & 0.0464$^{***}$ & 0.0351$^{***}$ & 0.0371$^{***}$ & 0.0284$^{***}$ & 0.0606$^{***}$ & 0.0458$^{***}$ & 0.0417$^{***}$ & 0.0325$^{***}$ \\
				& (0.0007) & (0.0008) & (0.0006) & (0.0007) & (0.0007) & (0.0008) & (0.0006) & (0.0007) \\
				\midrule
				\multicolumn{9}{l}{\textit{Panel B: Work Characteristics}} \\
				\quad hours\_per\_week & -- & 0.0035$^{***}$ & -- & 0.0027$^{***}$ & -- & 0.0084$^{***}$ & -- & 0.0038$^{***}$ \\
				&  & (0.0002) &  & (0.0001) &  & (0.0002) &  & (0.0001) \\
				\quad government & -- & 0.0059 & -- & -0.0004 & -- & -0.0072 & -- & -0.0060 \\
				&  & (0.0067) &  & (0.0061) &  & (0.0066) &  & (0.0055) \\
				\quad self\_employed & -- & 0.0218$^{***}$ & -- & -0.0085$^{**}$ & -- & 0.0192$^{***}$ & -- & -0.0080$^{**}$ \\
				&  & (0.0058) &  & (0.0043) &  & (0.0056) &  & (0.0040) \\
				\quad white\_collar & -- & 0.0649$^{***}$ & -- & 0.0561$^{***}$ & -- & 0.0480$^{***}$ & -- & 0.0478$^{***}$ \\
				&  & (0.0053) &  & (0.0049) &  & (0.0051) &  & (0.0045) \\
				\quad blue\_collar & -- & -0.0246$^{***}$ & -- & -0.0248$^{***}$ & -- & -0.0482$^{***}$ & -- & -0.0300$^{***}$ \\
				&  & (0.0036) &  & (0.0031) &  & (0.0034) &  & (0.0028) \\
				\quad service & -- & -0.0744$^{***}$ & -- & -0.0456$^{***}$ & -- & -0.0795$^{***}$ & -- & -0.0342$^{***}$ \\
				&  & (0.0024) &  & (0.0023) &  & (0.0024) &  & (0.0020) \\
				\midrule
				\multicolumn{9}{l}{\textit{Panel C: Financial}} \\
				\quad capital\_gain & -- & -- & 0.0000$^{***}$ & 0.0000$^{***}$ & -- & -- & 0.0000$^{***}$ & 0.0000$^{***}$ \\
				&  &  & (0.0000) & (0.0000) &  &  & (0.0000) & (0.0000) \\
				\quad capital\_loss & -- & -- & 0.0000$^{***}$ & 0.0000$^{***}$ & -- & -- & -0.0001$^{***}$ & -0.0000$^{***}$ \\
				&  &  & (0.0000) & (0.0000) &  &  & (0.0000) & (0.0000) \\
				\quad married & -- & -- & 0.2842$^{***}$ & 0.2785$^{***}$ & -- & -- & 0.2442$^{***}$ & 0.2423$^{***}$ \\
				&  &  & (0.0048) & (0.0054) &  &  & (0.0045) & (0.0047) \\
				\midrule
				McFadden $R^2$ & 0.199 & 0.233 & 0.376 & 0.400 & 0.257 & 0.296 & 0.467 & 0.492 \\
				Accuracy & 0.797 & 0.806 & 0.841 & 0.847 & 0.805 & 0.818 & 0.864 & 0.872 \\
				\bottomrule
				\end{tabular}
				\begin{tablenotes}
				\footnotesize
				\item \textit{Notes:} $^{***}$p$<$0.01, $^{**}$p$<$0.05, $^{*}$p$<$0.10. Estimates are the debiased cross-fitted window AME with $K=5$ folds; standard errors from the orthogonal score in parentheses. Step sizes $h_j$ and trimmed fractions are reported in Table~\ref{tab:window_settings}.
				\end{tablenotes}
		\end{threeparttable}
		}%
	\end{table*}
	
	Table~\ref{tab:adult_spec_search} presents the specification search results for the UCI Adult Income dataset. This is the largest of the three applications at 48,842 observations, and the estimates are correspondingly the most precisely determined in the paper. The demographic panel tells a consistent story across all specifications and both model classes. Age and education have positive effects on the probability of earning above \$50,000, and the female indicator is large and negative in the specifications that omit marital status, at $-$0.175 for logistic regression and $-$0.160 for XGBoost, collapsing to $-$0.012 and $-$0.009 once marital status enters in specification A+C. That collapse is the single most dramatic omitted variable pattern in the paper, and it is legible only because the AME puts every specification on the same probability scale. Education runs from 0.046 to 0.028 for logistic regression and 0.061 to 0.033 for XGBoost as controls accumulate. Marital status itself carries much the largest effect in the table, 0.278 and 0.242 in the full specification.

	The two model classes track each other closely throughout, and their standard errors are now of the same order everywhere --- 0.0061 and 0.0056 on the female indicator in the full specification, 0.0007 and 0.0007 on education. Under the refitting bootstrap the black-box standard errors were sometimes an order of magnitude smaller than the parametric ones, which the earlier draft read as an efficiency gain from the better-fitting model. That reading was mistaken. The two models estimate the same $\theta_{j,h}$ at the same sample size, so their standard errors ought to be similar, and the earlier gap measured how stable each fitting procedure was under resampling rather than how precisely either pinned down the target. The influence-function standard errors do behave as sampling variability should: they shrink with $n$, they are largest for the features carrying the least independent variation, and they do not reward a model for being reproducibly wrong.

	The financial panel illustrates the limits of what this design can say about capital gains and losses. Both variables are measured in raw dollars, so the reported effect is the response to a one-dollar change, and it displays as zero to four decimals in every cell but one --- as it did under the plug-in. The estimates are nonetheless significant at the one percent level, because the standard errors are smaller still; with nearly fifty thousand observations a substantively negligible effect can be sharply distinguished from zero, and this panel is a reminder that the stars answer a different question from the one the reader usually cares about. They also carry by far the heaviest trimming in the paper. Both are zero for the overwhelming majority of the sample and zero is the minimum of their observed support, so 92.3\% and 95.3\% of observations receive weight zero and the reported effect is identified only among filers who have nonzero gains or losses. The sign reversal between logistic regression and XGBoost that the earlier draft flagged in specification A+C does not survive the correction; both are zero to the displayed precision. Nothing here should be read as evidence about capital income, and a study interested in that question would want the variable re-expressed on a scale where a unit change is meaningful.
	
	\begin{table*}[htbp]
		\centering
		\small
		\begin{threeparttable}
			\caption{UCI Adult Income: Model Comparison. Average marginal effects on P(Income $>$ \$50k), full specification (A+B+C).}
			\label{tab:adult_model_comparison}
			\begin{tabular}{lrrrr}
				\toprule
				Feature & Logistic & Random Forest & XGBoost & Neural Net \\
				\midrule
				age & 0.0043$^{***}$ & 0.0056$^{***}$ & 0.0033$^{***}$ & 0.0041$^{***}$ \\
				& (0.0001) & (0.0001) & (0.0001) & (0.0001) \\
				female & -0.0153$^{**}$ & -0.0210$^{***}$ & -0.0128$^{**}$ & -0.0166$^{***}$ \\
				& (0.0061) & (0.0058) & (0.0056) & (0.0059) \\
				education\_num & 0.0284$^{***}$ & 0.0336$^{***}$ & 0.0325$^{***}$ & 0.0284$^{***}$ \\
				& (0.0007) & (0.0007) & (0.0007) & (0.0007) \\
				\midrule
				hours\_per\_week & 0.0027$^{***}$ & 0.0077$^{***}$ & 0.0038$^{***}$ & 0.0049$^{***}$ \\
				& (0.0001) & (0.0001) & (0.0001) & (0.0001) \\
				government & -0.0004 & -0.0025 & -0.0060 & -0.0060 \\
				& (0.0061) & (0.0056) & (0.0055) & (0.0060) \\
				self\_employed & -0.0085$^{**}$ & -0.0103$^{**}$ & -0.0080$^{**}$ & -0.0071$^{*}$ \\
				& (0.0043) & (0.0041) & (0.0040) & (0.0042) \\
				white\_collar & 0.0561$^{***}$ & 0.0606$^{***}$ & 0.0478$^{***}$ & 0.0383$^{***}$ \\
				& (0.0049) & (0.0047) & (0.0045) & (0.0047) \\
				blue\_collar & -0.0248$^{***}$ & -0.0316$^{***}$ & -0.0300$^{***}$ & -0.0533$^{***}$ \\
				& (0.0031) & (0.0030) & (0.0028) & (0.0031) \\
				service & -0.0456$^{***}$ & -0.0262$^{***}$ & -0.0342$^{***}$ & -0.0758$^{***}$ \\
				& (0.0023) & (0.0021) & (0.0020) & (0.0023) \\
				\midrule
				capital\_gain & 0.0000$^{***}$ & 0.0000$^{***}$ & 0.0000$^{***}$ & 0.0000$^{***}$ \\
				& (0.0000) & (0.0000) & (0.0000) & (0.0000) \\
				capital\_loss & 0.0000$^{***}$ & -0.0000$^{***}$ & -0.0000$^{***}$ & 0.0000$^{***}$ \\
				& (0.0000) & (0.0000) & (0.0000) & (0.0000) \\
				married & 0.2785$^{***}$ & 0.2504$^{***}$ & 0.2423$^{***}$ & 0.2438$^{***}$ \\
				& (0.0054) & (0.0048) & (0.0047) & (0.0048) \\
				\midrule
				McFadden $R^2$ & 0.400 & 0.450 & 0.492 & 0.449 \\
				Accuracy & 0.847 & 0.863 & 0.872 & 0.858 \\
				\bottomrule
				\end{tabular}
				\begin{tablenotes}
				\footnotesize
				\item \textit{Notes:} $^{***}$p$<$0.01, $^{**}$p$<$0.05, $^{*}$p$<$0.10. Estimates are the debiased cross-fitted window AME with $K=5$ folds; standard errors from the orthogonal score in parentheses. Step sizes $h_j$ and trimmed fractions are reported in Table~\ref{tab:window_settings}.
				\end{tablenotes}
		\end{threeparttable}
	\end{table*}
	
	Table~\ref{tab:adult_model_comparison} compares all four model classes on the full specification, and with nearly fifty thousand observations it is the cleanest test in the paper of whether the correction delivers a common estimand. It does. The education AME is 0.028, 0.034, 0.033 and 0.028 for logistic regression, random forest, XGBoost and the neural network, and marital status is 0.279, 0.250, 0.242 and 0.244 --- a spread of 0.005 and 0.036 respectively across four architectures with McFadden $R^2$ ranging from 0.400 to 0.492. The female indicator, which the plug-in spread from $-$0.027 down to $-$0.009 and left insignificant for the neural network, is $-$0.015, $-$0.021, $-$0.013 and $-$0.017 after correction, all four significant and all four within a standard error or two of one another. The variable the earlier draft singled out as sensitive to model specification turns out not to have been.

	The standard errors tell the same story from the other side. For a common target at a common sample size, four estimators using different nuisances should report standard errors of similar magnitude, and here they do: 0.0056 to 0.0061 on the female indicator, 0.0047 to 0.0054 on marital status, 0.0020 to 0.0023 on the service indicator. Under the plug-in they did not. The random forest reported 0.0001 on the female indicator and 0.0006 on marital status, one to two orders of magnitude below every other model, which the earlier draft attributed to ensemble averaging producing a smoother prediction surface and therefore more stable finite differences across bootstrap replicates. That description was accurate as a description of the bootstrap and wrong as a statement about precision: the replicates agreed with each other because every one of them was regularized the same way, not because the quantity was well determined. The convergence of the four standard-error columns to a common magnitude is what the orthogonal score is supposed to produce, and it is the more convincing for happening at the sample size where the asymptotics have the best claim to apply.
	
	\begin{sidewaystable}[htbp]
		\centering
		\small
		\begin{threeparttable}
			\caption{UCI Adult Income: Method Comparison. AME, SHAP, and PDP estimates for P(Income $>$ \$50k), full specification (A+B+C).}
			\label{tab:adult_method_comparison}
			\begin{tabular}{lrrrrrrrrrrrr}
				\toprule
				& \multicolumn{3}{c}{Logistic} & \multicolumn{3}{c}{Random Forest} & \multicolumn{3}{c}{XGBoost} & \multicolumn{3}{c}{Neural Net} \\
				\cmidrule(lr){2-4} \cmidrule(lr){5-7} \cmidrule(lr){8-10} \cmidrule(lr){11-13}
				& AME & SHAP & PDP & AME & SHAP & PDP & AME & SHAP & PDP & AME & SHAP & PDP \\
				\midrule
				age & 0.0043 & -0.0270 & 0.0028 & 0.0056 & -0.0068 & 0.0024 & 0.0033 & -0.0097 & 0.0036 & 0.0041 & 0.0044 & 0.0033 \\
				female & -0.0153 & 0.0022 & -0.0270 & -0.0210 & 0.0011 & -0.0218 & -0.0128 & 0.0015 & -0.0122 & -0.0166 & -0.0001 & -0.0222 \\
				education\_num & 0.0284 & -0.1186 & 0.0269 & 0.0336 & -0.0077 & 0.0137 & 0.0325 & -0.0093 & 0.0170 & 0.0284 & -0.0051 & 0.0207 \\
				\midrule
				hours\_per\_week & 0.0027 & -0.0691 & 0.0033 & 0.0077 & -0.0054 & 0.0020 & 0.0038 & -0.0080 & 0.0028 & 0.0049 & -0.0032 & 0.0028 \\
				government & -0.0004 & -0.0009 & 0.0026 & -0.0025 & -0.0003 & 0.0051 & -0.0060 & -0.0008 & -0.0003 & -0.0060 & 0.0002 & 0.0082 \\
				self\_employed & -0.0085 & 0.0016 & -0.0265 & -0.0103 & -0.0003 & 0.0032 & -0.0080 & -0.0002 & -0.0055 & -0.0071 & -0.0008 & -0.0045 \\
				white\_collar & 0.0561 & 0.0127 & 0.0538 & 0.0606 & 0.0031 & 0.0600 & 0.0478 & 0.0025 & 0.0474 & 0.0383 & 0.0007 & 0.0312 \\
				blue\_collar & -0.0248 & 0.0136 & -0.0273 & -0.0316 & 0.0010 & -0.0240 & -0.0300 & 0.0004 & -0.0251 & -0.0533 & -0.0008 & -0.0531 \\
				service & -0.0456 & -0.0027 & -0.0427 & -0.0262 & 0.0003 & -0.0176 & -0.0342 & 0.0005 & -0.0343 & -0.0758 & 0.0011 & -0.0558 \\
				\midrule
				capital\_gain & 0.0000 & -0.2627 & 0.0000 & 0.0000 & -0.0224 & 0.0000 & 0.0000 & -0.0264 & -0.0001 & 0.0000 & 0.0042 & 0.0000 \\
				capital\_loss & 0.0000 & 0.0027 & -- & -0.0000 & -0.0034 & -- & -0.0000 & -0.0062 & -- & 0.0000 & -0.0115 & -- \\
				married & 0.2785 & -0.0768 & 0.2749 & 0.2504 & -0.0064 & 0.2446 & 0.2423 & -0.0049 & 0.2357 & 0.2438 & 0.0149 & 0.1958 \\
				\bottomrule
				\end{tabular}
				\begin{tablenotes}
				\footnotesize
				\item \textit{Notes:} AME = Average Marginal Effect (marginfx). SHAP = mean signed SHAP value (GradientExplainer). PDP = slope of partial dependence plot. Full specification (A+B+C) only. No standard errors reported for SHAP and PDP.
				\end{tablenotes}
		\end{threeparttable}
	\end{sidewaystable}
	
	Table~\ref{tab:adult_method_comparison} compares AMEs against SHAP values and partial dependence plot slopes across all four models. The SHAP column is the more straightforward contrast, and it is stark. Mean signed SHAP values are on a different scale entirely and in several cases carry the wrong sign: for logistic regression the SHAP value for education is $-$0.119 against an AME of 0.028, and for marital status it is $-$0.077 against an AME of 0.279 --- opposite in sign on the two largest effects in the application. The SHAP value for age is negative for three of the four models while every AME is positive. The most striking entry is capital gains, whose SHAP value of $-$0.263 for logistic regression is the largest SHAP magnitude anywhere in the table, attached to a variable whose AME is zero to four decimals and 92\% of whose observations are trimmed. None of this is a defect of the SHAP implementation; it is a consequence of mean signed SHAP values not being marginal effects and not admitting that interpretation.

	The partial dependence column behaves as it did in the other two applications: it agrees with the debiased AME where the underlying plug-in is nearly unbiased and falls short of it elsewhere. For logistic regression, the best-specified model here, agreement is close on the large effects --- marital status 0.279 against 0.275, the white collar indicator 0.056 against 0.054, service $-$0.046 against $-$0.043, education 0.028 against 0.027. For the two tree ensembles the same features separate substantially: the random forest education AME of 0.034 sits against a PDP slope of 0.014, and its hours-per-week AME of 0.008 against a slope of 0.002. A piecewise constant fit has little local slope to report between splits, and the partial dependence average inherits that flatness while the orthogonal correction does not. Partial dependence slopes for capital losses are unavailable for every model, a limitation of the grid construction when almost all mass sits at a single value. As with SHAP, no standard errors accompany either column, so a researcher working from them has point estimates of unknown precision and no way to test whether any variable has a statistically distinguishable effect.
	
	\subsection{UCI Credit Default}
	
	The second dataset is the UCI Default of Credit Card Clients dataset \citep{yeh2009, UCI}, comprising 30,000 observations of Taiwanese credit card holders. The outcome is whether a cardholder defaulted on their payment in the following month. Features are organized into demographics (age, sex, education, marital status), payment history (average payment status, months delayed), and bill and payment amounts (average bill amount, average payment amount, payment ratio). The specification search considers logistic regression and a TensorFlow neural network, providing a direct comparison between a classical parametric model and a deep learning approach on the same inferential task. Two features here carry substantial trimming. Months delayed takes seven integer values and 70.9\% of the sample sits on one of the two outer ones, so its effect is identified among cardholders strictly inside that range; education takes four levels and is trimmed at 36.8\%. Both take the floored step size $h_j = 0.5$, so their reported effects are one-level contrasts. The payment ratio, which turns out to carry the largest effect in the application, is continuous and trimmed at 18.4\%.
	
	\begin{table*}[htbp]
		\centering
		\small
		\resizebox{\textwidth}{!}{%
			\begin{threeparttable}
				\caption{UCI Credit Card Default: Specification Search. Average marginal effects on P(Default).}
				\label{tab:credit_default_spec_search}
				\begin{tabular}{lrrrrrrrr}
					\toprule
					& \multicolumn{4}{c}{Logistic} & \multicolumn{4}{c}{Neural Net} \\
					\cmidrule(lr){2-5} \cmidrule(lr){6-9}
					& Spec A & A+B & A+C & A+B+C & Spec A & A+B & A+C & A+B+C \\
					\midrule
					\multicolumn{9}{l}{\textit{Panel A: Demographics}} \\
					\quad age & -0.0015$^{***}$ & -0.0002 & -0.0005 & 0.0001 & -0.0026$^{***}$ & -0.0006$^{*}$ & -0.0010$^{***}$ & -0.0000 \\
					& (0.0004) & (0.0003) & (0.0004) & (0.0003) & (0.0004) & (0.0003) & (0.0004) & (0.0003) \\
					\quad female & -0.0358$^{***}$ & -0.0206$^{***}$ & -0.0291$^{***}$ & -0.0198$^{***}$ & -0.0364$^{***}$ & -0.0205$^{***}$ & -0.0265$^{***}$ & -0.0176$^{***}$ \\
					& (0.0050) & (0.0046) & (0.0049) & (0.0046) & (0.0050) & (0.0046) & (0.0049) & (0.0045) \\
					\quad education & -0.0001 & -0.0027 & -0.0054$^{**}$ & -0.0040$^{*}$ & -0.0008 & -0.0035 & -0.0048$^{**}$ & -0.0040$^{*}$ \\
					& (0.0022) & (0.0022) & (0.0023) & (0.0022) & (0.0021) & (0.0021) & (0.0022) & (0.0022) \\
					\quad married & 0.0253$^{***}$ & 0.0189$^{***}$ & 0.0303$^{***}$ & 0.0206$^{***}$ & 0.0260$^{***}$ & 0.0152$^{**}$ & 0.0269$^{***}$ & 0.0186$^{***}$ \\
					& (0.0068) & (0.0064) & (0.0066) & (0.0063) & (0.0068) & (0.0064) & (0.0067) & (0.0063) \\
					\midrule
					\multicolumn{9}{l}{\textit{Panel B: Payment History}} \\
					\quad avg\_pay\_status & -- & -0.0140$^{***}$ & -- & -0.0218$^{***}$ & -- & 0.0073$^{*}$ & -- & 0.0040 \\
					&  & (0.0041) &  & (0.0075) &  & (0.0041) &  & (0.0072) \\
					\quad months\_delayed & -- & 0.0347$^{***}$ & -- & 0.0301$^{***}$ & -- & 0.0298$^{***}$ & -- & 0.0254$^{***}$ \\
					&  & (0.0011) &  & (0.0012) &  & (0.0011) &  & (0.0012) \\
					\midrule
					\multicolumn{9}{l}{\textit{Panel C: Bill and Payment Amounts}} \\
					\quad avg\_bill\_amt & -- & -- & -0.0000$^{***}$ & -0.0000$^{**}$ & -- & -- & -0.0000 & -0.0000$^{*}$ \\
					&  &  & (0.0000) & (0.0000) &  &  & (0.0000) & (0.0000) \\
					\quad avg\_pay\_amt & -- & -- & -0.0000$^{***}$ & 0.0000 & -- & -- & -0.0000$^{***}$ & -0.0000$^{***}$ \\
					&  &  & (0.0000) & (0.0000) &  &  & (0.0000) & (0.0000) \\
					\quad pay\_ratio & -- & -- & -0.3521$^{***}$ & -0.2651$^{***}$ & -- & -- & -0.2572$^{***}$ & -0.1943$^{**}$ \\
					&  &  & (0.0731) & (0.0730) &  &  & (0.0792) & (0.0773) \\
					\midrule
					McFadden $R^2$ & 0.003 & 0.130 & 0.031 & 0.137 & 0.010 & 0.155 & 0.055 & 0.166 \\
					Accuracy & 0.779 & 0.803 & 0.779 & 0.803 & 0.779 & 0.808 & 0.779 & 0.808 \\
					\bottomrule
					\end{tabular}
					\begin{tablenotes}
					\footnotesize
					\item \textit{Notes:} $^{***}$p$<$0.01, $^{**}$p$<$0.05, $^{*}$p$<$0.10. Estimates are the debiased cross-fitted window AME with $K=5$ folds; standard errors from the orthogonal score in parentheses. Step sizes $h_j$ and trimmed fractions are reported in Table~\ref{tab:window_settings}.
					\end{tablenotes}
			\end{threeparttable}
		}%
	\end{table*}
	
	Table~\ref{tab:credit_default_spec_search} presents the specification search. The dominant finding is the importance of payment history and of the payment ratio. Months delayed enters at 0.035 for logistic regression and 0.030 for the neural network in specification A+B and remains at 0.030 and 0.025 in the full specification, while the payment ratio, once it is available in specifications A+C and A+B+C, is much the largest effect in the table: $-$0.265 for logistic regression and $-$0.194 for the neural network in the full specification. The demographic panel tells a different story from the Adult Income dataset: age is statistically indistinguishable from zero in nearly every specification for both models, and the female indicator, while consistently negative and significant, shrinks substantially once payment history is controlled for, falling from $-$0.036 to $-$0.020 for logistic regression and from $-$0.036 to $-$0.018 for the neural network. This pattern of attenuation is cleanly legible in the AME framework because the estimates are on a common probability scale across specifications.

	What is most striking in this table is how little separates the two columns. In specification A the female AME is $-$0.0358 for logistic regression and $-$0.0364 for the neural network; married is 0.0253 and 0.0260; in the full specification the pair is $-$0.0198 and $-$0.0176, and 0.0206 and 0.0186. A parametric model and a four-layer network, differing in McFadden $R^2$ by three percentage points, produce estimates that agree to the third decimal place across sixteen model-specification cells. This is what the orthogonal score is supposed to deliver and what the plug-in did not: since each column estimates the same $\theta_{j,h}$ and the correction absorbs the first-order effect of the fitted regression, the remaining gap between the two columns reflects only the second-order term and sampling noise. The standard errors are also now nearly identical between the two models --- 0.0046 and 0.0045 on the female indicator, 0.0012 and 0.0012 on months delayed --- rather than differing by the order of magnitude the refitting bootstrap reported, where the network's standard error for months delayed was 0.032 against the logistic model's 0.002. Both figures were describing the stability of a particular fitting procedure under resampling rather than the sampling variability of an estimate of a fixed target.
	
	\begin{table*}[htbp]
		\centering
		\small
		\begin{threeparttable}
			\caption{UCI Credit Card Default: Model Comparison. Average marginal effects on P(Default), full specification (A+B+C).}
			\label{tab:credit_default_model_comparison}
			\begin{tabular}{lrrrr}
				\toprule
				Feature & Logistic & Random Forest & XGBoost & Neural Net \\
				\midrule
				age & 0.0001 & -0.0002 & -0.0005 & -0.0000 \\
				& (0.0003) & (0.0003) & (0.0003) & (0.0003) \\
				female & -0.0198$^{***}$ & -0.0161$^{***}$ & -0.0175$^{***}$ & -0.0176$^{***}$ \\
				& (0.0046) & (0.0045) & (0.0045) & (0.0045) \\
				education & -0.0040$^{*}$ & -0.0018 & 0.0034 & -0.0040$^{*}$ \\
				& (0.0022) & (0.0022) & (0.0022) & (0.0022) \\
				married & 0.0206$^{***}$ & 0.0222$^{***}$ & 0.0205$^{***}$ & 0.0186$^{***}$ \\
				& (0.0063) & (0.0062) & (0.0062) & (0.0063) \\
				\midrule
				avg\_pay\_status & -0.0218$^{***}$ & 0.0061 & 0.0907$^{***}$ & 0.0040 \\
				& (0.0075) & (0.0074) & (0.0075) & (0.0072) \\
				months\_delayed & 0.0301$^{***}$ & 0.0134$^{***}$ & 0.0191$^{***}$ & 0.0254$^{***}$ \\
				& (0.0012) & (0.0011) & (0.0011) & (0.0012) \\
				\midrule
				avg\_bill\_amt & -0.0000$^{**}$ & -0.0000$^{***}$ & -0.0000$^{***}$ & -0.0000$^{*}$ \\
				& (0.0000) & (0.0000) & (0.0000) & (0.0000) \\
				avg\_pay\_amt & 0.0000 & -0.0000$^{***}$ & -0.0000$^{***}$ & -0.0000$^{***}$ \\
				& (0.0000) & (0.0000) & (0.0000) & (0.0000) \\
				pay\_ratio & -0.2651$^{***}$ & -0.4890$^{***}$ & -0.5836$^{***}$ & -0.1943$^{**}$ \\
				& (0.0730) & (0.0670) & (0.0685) & (0.0773) \\
				\midrule
				McFadden $R^2$ & 0.137 & 0.210 & 0.194 & 0.166 \\
				Accuracy & 0.803 & 0.821 & 0.815 & 0.808 \\
				\bottomrule
				\end{tabular}
				\begin{tablenotes}
				\footnotesize
				\item \textit{Notes:} $^{***}$p$<$0.01, $^{**}$p$<$0.05, $^{*}$p$<$0.10. Estimates are the debiased cross-fitted window AME with $K=5$ folds; standard errors from the orthogonal score in parentheses. Step sizes $h_j$ and trimmed fractions are reported in Table~\ref{tab:window_settings}.
				\end{tablenotes}
		\end{threeparttable}
	\end{table*}
	
	Table~\ref{tab:credit_default_model_comparison} compares all four models on the full specification, and it is the sharpest illustration in the paper of what the orthogonal correction changes. Under the plug-in the four models flatly contradicted one another on the two variables that matter. Months delayed was 0.083 for logistic regression, 0.018 for the random forest, 0.021 for XGBoost and 0.112 for the neural network, a sixfold spread; the payment ratio was $-$0.431 and $-$0.440 for the two tree ensembles, $-$0.045 for logistic regression, and 0.0009 --- indistinguishable from no effect at all --- for the neural network. A researcher reading that table had to conclude either that the models were detecting genuinely different structure, or that at least three of the four were wrong, with no way to tell which.

	The debiased estimates resolve the contradiction rather than splitting the difference. Months delayed becomes 0.030, 0.013, 0.019 and 0.025, and the payment ratio becomes $-$0.265, $-$0.489, $-$0.584 and $-$0.194: every model now agrees that a higher payment ratio is associated with a substantially lower probability of default, including the neural network, whose plug-in estimate had put the effect at zero. The demographic panel converges the same way. The female AME was $-$0.020, $-$0.006, $-$0.015 and $-$0.014 under the plug-in and is $-$0.0198, $-$0.0161, $-$0.0175 and $-$0.0176 under the correction; married was 0.022, 0.005, 0.014 and 0.019 and is now 0.0206, 0.0222, 0.0205 and 0.0186. The variation that the earlier draft read as four models weighing the covariates differently was in the main the regularization bias $\E[\alpha_h(\hat f - f)]$, which differs across learners because their regularization differs, and which the correction removes. What survives is a common estimate of a common target.

	Convergence is not complete, and the payment ratio shows where it stops. The four debiased estimates still span $-$0.194 to $-$0.584, a range far wider than their standard errors of roughly 0.07, so the models have not been brought into agreement to within sampling error on that variable. Payment ratio is the most skewed feature in the design and is trimmed at 18.4\%, and the second-order product term that Theorem~\ref{thm:main} requires to be $o_p(n^{-1/2})$ is plausibly not negligible here for the weaker nuisance fits. The correct reading is that the correction removes the first-order bias and leaves a residual that the theory bounds but does not eliminate, which is visible precisely where the nuisance problem is hardest. Finally, the random forest standard errors, which the plug-in bootstrap reported as 0.0000 to four decimal places on the female indicator, are 0.0045 under the influence-function formula --- essentially identical to the other three models, as they should be for a common estimand at a common sample size.
	
		\begin{sidewaystable}[htbp]
		\centering
		\small
		\begin{threeparttable}
			\caption{UCI Credit Card Default: Method Comparison. AME, SHAP, and PDP estimates for P(Default), full specification (A+B+C).}
			\label{tab:credit_default_method_comparison}
			\begin{tabular}{lrrrrrrrrrrrr}
				\toprule
				& \multicolumn{3}{c}{Logistic} & \multicolumn{3}{c}{Random Forest} & \multicolumn{3}{c}{XGBoost} & \multicolumn{3}{c}{Neural Net} \\
				\cmidrule(lr){2-4} \cmidrule(lr){5-7} \cmidrule(lr){8-10} \cmidrule(lr){11-13}
				& AME & SHAP & PDP & AME & SHAP & PDP & AME & SHAP & PDP & AME & SHAP & PDP \\
				\midrule
				age & 0.0001 & -0.0003 & 0.0002 & -0.0002 & 0.0011 & -0.0000 & -0.0005 & 0.0004 & -0.0002 & -0.0000 & -0.0010 & 0.0012 \\
				female & -0.0198 & -0.0005 & -0.0198 & -0.0161 & 0.0001 & -0.0062 & -0.0175 & 0.0002 & -0.0146 & -0.0176 & 0.0009 & -0.0200 \\
				education & -0.0040 & 0.0002 & 0.0002 & -0.0018 & 0.0006 & -0.0057 & 0.0034 & -0.0005 & -0.0356 & -0.0040 & -0.0016 & 0.0082 \\
				married & 0.0206 & -0.0038 & 0.0220 & 0.0222 & 0.0002 & 0.0054 & 0.0205 & 0.0002 & 0.0142 & 0.0186 & 0.0022 & 0.0217 \\
				\midrule
				avg\_pay\_status & -0.0218 & -0.0144 & -0.0246 & 0.0061 & 0.0029 & 0.0610 & 0.0907 & 0.0027 & 0.0404 & 0.0040 & -0.0116 & 0.0097 \\
				months\_delayed & 0.0301 & -0.0148 & 0.1107 & 0.0134 & -0.0008 & 0.0456 & 0.0191 & 0.0001 & 0.0404 & 0.0254 & -0.0154 & 0.0653 \\
				\midrule
				avg\_bill\_amt & -0.0000 & -0.0001 & 0.0000 & -0.0000 & -0.0018 & -0.0000 & -0.0000 & -0.0046 & -0.0000 & -0.0000 & 0.0036 & 0.0000 \\
				avg\_pay\_amt & 0.0000 & 0.0026 & -0.0000 & -0.0000 & -0.0006 & -0.0000 & -0.0000 & -0.0005 & -0.0000 & -0.0000 & -0.0025 & -0.0000 \\
				pay\_ratio & -0.2651 & 0.0100 & -0.0444 & -0.4890 & 0.0012 & -0.0998 & -0.5836 & 0.0005 & -0.0953 & -0.1943 & 0.0015 & -0.0170 \\
				\bottomrule
				\end{tabular}
				\begin{tablenotes}
				\footnotesize
				\item \textit{Notes:} AME = Average Marginal Effect (marginfx). SHAP = mean signed SHAP value (GradientExplainer). PDP = slope of partial dependence plot. Full specification (A+B+C) only. No standard errors reported for SHAP and PDP.
				\end{tablenotes}
		\end{threeparttable}
	\end{sidewaystable}
	
	Table~\ref{tab:credit_default_method_comparison} again reveals a systematic disconnect between SHAP values and AMEs. For the payment history variables the disagreement is stark: the AME for months delayed is 0.030 for logistic regression while the mean signed SHAP value is $-$0.015, opposite in sign. The payment ratio illustrates the most consequential divergence: the random forest and XGBoost AMEs are $-$0.489 and $-$0.584, indicating that a higher payment ratio is associated with a large reduction in default probability, while the corresponding SHAP values are 0.0012 and 0.0005 --- near zero and in the wrong direction. A researcher relying on SHAP alone would conclude that payment ratio is essentially irrelevant to default risk. With no standard errors attached to the SHAP estimates there is no statistical basis on which to adjudicate.

	The partial dependence column separates from the AME column here in the same pattern as in the Ames application, and for the same reason: a partial dependence slope is a plug-in average over a single full-sample fit, so it carries the regularization bias that the correction removes. The separation tracks the quality of the underlying fit. On the female indicator, where the logistic model is close to correctly specified, the AME of $-$0.0198 and the PDP slope of $-$0.0198 agree exactly; for the random forest on the same variable they are $-$0.0161 and $-$0.0062. On the payment ratio, the hardest nuisance problem in the design, the gap is largest and always toward zero: $-$0.265 against $-$0.044 for logistic regression, $-$0.489 against $-$0.100 for the random forest, $-$0.584 against $-$0.095 for XGBoost, $-$0.194 against $-$0.017 for the neural network. Months delayed runs the other way, with PDP slopes of 0.111, 0.046, 0.040 and 0.065 against debiased AMEs of 0.030, 0.013, 0.019 and 0.025, which is what one expects when an uncorrected estimator distributes a shared effect between two correlated features differently than the corrected one does. Taken across all three applications, the attenuation is the dominant pattern rather than a universal one: restricting attention to the effects the debiased estimator resolves away from zero, the partial dependence slope is smaller in absolute value than the AME in 62 of 84 cases, with a median ratio of 0.47 on Ames Housing, 0.60 here and 0.95 on Adult Income. It is most pronounced for the random forest, the learner whose regularization bias the simulations isolate.
	
	\subsection{Ames Housing}
	
	The third dataset is the Ames Housing dataset \citep{decock2011}, comprising approximately 1,460 observations of residential home sales in Ames, Iowa between 2006 and 2010. The outcome is sale price in dollars, so average marginal effects are interpreted in dollar units --- a one-unit change in a continuous feature is associated with a $\widehat{\theta}_{j,h}$ change in predicted sale price, on average across the untrimmed sample. Features are organized into physical characteristics (above-ground living area, number of bedrooms, number of bathrooms), quality and condition (overall quality rating, condition rating, year built, year remodeled), and lot and location (lot area, garage area, neighborhood tier indicators). The specification search considers linear regression and a random forest, covering the regression case and completing the sweep across outcome types. Six of the ten continuous features here are integer counts or ratings and so take the floored step size $h_j = 0.5$, which makes their reported effects one-unit contrasts: one more bedroom, one more point of overall quality. One feature requires a narrower reading than the rest. Half bathrooms take only the values zero, one and two, and both outer levels lie on the boundary of the observed support, so 63.4\% of the sample is trimmed and the half-bath effect is identified only among homes that have exactly one. The remaining features in this application are trimmed at 15\% or below.
	
	\begin{table*}[htbp]
		\centering
		\small
		\resizebox{\textwidth}{!}{%
			\begin{threeparttable}
				\caption{Ames Housing: Specification Search. Average marginal effects on Sale Price (USD).}
				\label{tab:ames_housing_spec_search}
				\begin{tabular}{lrrrrrrrr}
					\toprule
					& \multicolumn{4}{c}{Linear} & \multicolumn{4}{c}{Random Forest} \\
					\cmidrule(lr){2-5} \cmidrule(lr){6-9}
					& Spec A & A+B & A+C & A+B+C & Spec A & A+B & A+C & A+B+C \\
					\midrule
					\multicolumn{9}{l}{\textit{Panel A: Physical Characteristics}} \\
					\quad GrLivArea & 138.96$^{***}$ & 106.16$^{***}$ & 84.82$^{***}$ & 68.05$^{***}$ & 128.77$^{***}$ & 103.42$^{***}$ & 90.19$^{***}$ & 81.50$^{***}$ \\
					& (20.51) & (15.58) & (14.25) & (12.14) & (13.06) & (8.98) & (8.97) & (8.77) \\
					\quad BedroomAbvGr & -33209.66$^{***}$ & -14046.40$^{*}$ & -24282.55$^{***}$ & -26454.60$^{*}$ & -28632.01$^{***}$ & -11415.18$^{***}$ & -11292.79$^{**}$ & -9212.02 \\
					& (5594.51) & (7435.20) & (9027.62) & (14881.54) & (3382.73) & (3920.43) & (5576.82) & (8925.07) \\
					\quad FullBath & 18589.00$^{***}$ & -10533.55$^{**}$ & 7261.14 & -10578.42 & 16236.01$^{***}$ & -6616.40$^{*}$ & 5652.46 & 7858.42 \\
					& (5990.54) & (4912.27) & (4987.91) & (9969.12) & (3779.18) & (3670.56) & (3754.82) & (8988.79) \\
					\quad HalfBath & -11228.94$^{**}$ & -8180.81$^{**}$ & -9092.88$^{**}$ & 1028.31 & -7189.23$^{*}$ & -6638.05$^{***}$ & -11578.88$^{**}$ & -10317.96$^{**}$ \\
					& (4624.11) & (3242.49) & (4451.09) & (5697.41) & (4264.04) & (2525.05) & (4601.22) & (4713.92) \\
					\midrule
					\multicolumn{9}{l}{\textit{Panel B: Quality and Condition}} \\
					\quad OverallQual & -- & 15819.18$^{***}$ & -- & 5436.28 & -- & 11911.79$^{***}$ & -- & 16255.82$^{***}$ \\
					&  & (1589.02) &  & (4964.23) &  & (1342.87) &  & (4829.24) \\
					\quad OverallCond & -- & 6477.48$^{*}$ & -- & 2890.60 & -- & 8281.97$^{***}$ & -- & 9190.11$^{***}$ \\
					&  & (3687.03) &  & (5496.46) &  & (2434.71) &  & (3492.01) \\
					\quad YearBuilt & -- & 443.39 & -- & -141.08 & -- & 1033.46$^{***}$ & -- & 1323.39$^{***}$ \\
					&  & (277.71) &  & (404.21) &  & (190.55) &  & (419.89) \\
					\quad YearRemodAdd & -- & 192.69 & -- & 1743.45$^{**}$ & -- & 556.20$^{***}$ & -- & -311.44 \\
					&  & (161.82) &  & (740.09) &  & (173.98) &  & (753.53) \\
					\midrule
					\multicolumn{9}{l}{\textit{Panel C: Lot and Location}} \\
					\quad LotArea & -- & -- & 5.51 & 23.02$^{**}$ & -- & -- & -7.53 & -12.51 \\
					&  &  & (3.93) & (11.37) &  &  & (8.06) & (11.14) \\
					\quad GarageArea & -- & -- & 40.08$^{*}$ & -102.78 & -- & -- & 103.31$^{***}$ & 90.33 \\
					&  &  & (20.94) & (71.57) &  &  & (34.64) & (59.74) \\
					\quad nbhd\_high & -- & -- & 128801.76$^{***}$ & 89812.31$^{**}$ & -- & -- & 66978.55$^{***}$ & 43512.42$^{*}$ \\
					&  &  & (34103.92) & (38393.43) &  &  & (16110.93) & (24635.87) \\
					\quad nbhd\_mid & -- & -- & 77339.38$^{**}$ & 74336.31$^{*}$ & -- & -- & 42305.81$^{***}$ & 33767.43 \\
					&  &  & (33663.73) & (39182.81) &  &  & (15387.40) & (26676.67) \\
					\midrule
					$R^2$ & 0.584 & 0.758 & 0.739 & 0.802 & 0.829 & 0.937 & 0.926 & 0.953 \\
					RMSE & 51,243.9 & 39,080.1 & 40,596.4 & 35,355.5 & 32,830.9 & 19,881.6 & 21,534.2 & 17,244.4 \\
					\bottomrule
					\end{tabular}
					\begin{tablenotes}
					\footnotesize
					\item \textit{Notes:} $^{***}$p$<$0.01, $^{**}$p$<$0.05, $^{*}$p$<$0.10. Estimates are the debiased cross-fitted window AME with $K=5$ folds; standard errors from the orthogonal score in parentheses. Step sizes $h_j$ and trimmed fractions are reported in Table~\ref{tab:window_settings}.
					\end{tablenotes}
			\end{threeparttable}
		}%
	\end{table*}
	
	Table~\ref{tab:ames_housing_spec_search} presents the specification search. The regression setting allows AMEs to be interpreted directly in dollar units, which makes the substantive findings immediately legible. Above-ground living area carries an AME of \$139 per square foot in the physical-characteristics specification for linear regression, falling to \$106, \$85 and finally \$68 as quality, condition, and location controls are added --- a pattern of positive omitted variable bias that makes intuitive sense, since larger homes tend to be in better neighborhoods and of higher quality. The random forest traces the same path from \$129 to \$81. That the two model classes track each other this closely across the sequence is the first thing to note about the debiased estimates: because the orthogonal correction removes the first-order dependence on the fitted regression, the two columns are estimating the same window effect rather than reporting the slope each model happens to believe, and the linear column is no longer simply the ordinary least squares coefficient. The standard errors are now of comparable size across the two model classes --- \$12 for linear regression and \$9 for the random forest on GrLivArea in the full specification --- rather than differing by the order of magnitude that the refitting bootstrap produced, where the ensemble's smooth prediction surface made every replicate agree with every other and drove the reported standard error to \$1.93. Those intervals were narrow around a biased center; the influence-function intervals are wider and honest.

	The cost of that honesty is visible in the quality panel. The overall quality rating is estimated at \$15,819 per point for linear regression and \$11,912 for the random forest in specification A+B, both precisely determined, but in the full specification the linear estimate falls to \$5,436 with a standard error of \$4,964 and is no longer distinguishable from zero, while the random forest gives \$16,256 with a standard error of \$4,829. With $n = 1{,}460$ observations, twelve features and a second-order sieve for the representer, the full specification is close to the limit of what this sample can support: the estimator reports that limit rather than concealing it. The neighborhood tier indicators show the same attenuation as before --- the high neighborhood effect falls from \$128,802 to \$89,812 for linear regression and from \$66,979 to \$43,512 for the random forest as quality controls are added --- but with standard errors near \$25{,}000 to \$38{,}000, the attenuation is a pattern in the point estimates rather than a difference the data can resolve.
	
	\begin{table*}[htbp]
		\centering
		\small
		\begin{threeparttable}
			\caption{Ames Housing: Model Comparison. Average marginal effects on Sale Price (USD), full specification (A+B+C).}
			\label{tab:ames_housing_model_comparison}
			\begin{tabular}{lrrrr}
				\toprule
				Feature & Linear & Random Forest & XGBoost & Neural Net \\
				\midrule
				GrLivArea & 68.05$^{***}$ & 81.50$^{***}$ & 70.27$^{***}$ & 66.20$^{***}$ \\
				& (12.14) & (8.77) & (6.60) & (6.24) \\
				BedroomAbvGr & -26454.60$^{*}$ & -9212.02 & -3464.31 & -12638.60$^{***}$ \\
				& (14881.54) & (8925.07) & (4508.19) & (3830.14) \\
				FullBath & -10578.42 & 7858.42 & 4761.27 & -5055.91 \\
				& (9969.12) & (8988.79) & (7661.72) & (5080.88) \\
				HalfBath & 1028.31 & -10317.96$^{**}$ & -10906.00$^{***}$ & -2175.44 \\
				& (5697.41) & (4713.92) & (3940.81) & (3286.03) \\
				\midrule
				OverallQual & 5436.28 & 16255.82$^{***}$ & 14913.45$^{***}$ & 8524.39$^{***}$ \\
				& (4964.23) & (4829.24) & (4271.98) & (2394.41) \\
				OverallCond & 2890.60 & 9190.11$^{***}$ & 12102.59$^{***}$ & 7234.42$^{**}$ \\
				& (5496.46) & (3492.01) & (2519.57) & (2920.49) \\
				YearBuilt & -141.08 & 1323.39$^{***}$ & 1352.73$^{***}$ & 230.04 \\
				& (404.21) & (419.89) & (366.04) & (267.61) \\
				YearRemodAdd & 1743.45$^{**}$ & -311.44 & -289.89 & 942.55$^{***}$ \\
				& (740.09) & (753.53) & (673.39) & (351.83) \\
				\midrule
				LotArea & 23.02$^{**}$ & -12.51 & -10.83 & 10.04$^{**}$ \\
				& (11.37) & (11.14) & (9.91) & (4.81) \\
				GarageArea & -102.78 & 90.33 & 92.29$^{*}$ & -6.11 \\
				& (71.57) & (59.74) & (49.36) & (27.32) \\
				nbhd\_high & 89812.31$^{**}$ & 43512.42$^{*}$ & 33476.93$^{**}$ & 35458.39$^{**}$ \\
				& (38393.43) & (24635.87) & (16103.49) & (17705.63) \\
				nbhd\_mid & 74336.31$^{*}$ & 33767.43 & 17468.12 & 36951.86$^{**}$ \\
				& (39182.81) & (26676.67) & (18286.73) & (17998.56) \\
				\midrule
				$R^2$ & 0.802 & 0.953 & 0.954 & 0.932 \\
				RMSE & 35,355.5 & 17,244.4 & 17,036.6 & 20,699.9 \\
				\bottomrule
				\end{tabular}
				\begin{tablenotes}
				\footnotesize
				\item \textit{Notes:} $^{***}$p$<$0.01, $^{**}$p$<$0.05, $^{*}$p$<$0.10. Estimates are the debiased cross-fitted window AME with $K=5$ folds; standard errors from the orthogonal score in parentheses. Step sizes $h_j$ and trimmed fractions are reported in Table~\ref{tab:window_settings}.
				\end{tablenotes}
		\end{threeparttable}
	\end{table*}
	
	Table~\ref{tab:ames_housing_model_comparison} compares all four models on the full specification. All four are well fitted, with $R^2$ ranging from 0.802 for linear regression to 0.954 for XGBoost, the random forest at 0.953 and the neural network at 0.932 with normalization. The GrLivArea AME is \$68, \$81, \$70, and \$66 for linear regression, random forest, XGBoost, and neural network respectively --- a range of \$15 across four models that differ fundamentally in their functional form, and all four intervals overlap comfortably. Under the debiased estimator this agreement means something different from what it meant under the plug-in. Every column is now an estimate of the same $\theta_{j,h}$, so agreement is evidence that the four nuisance estimates are all good enough for the correction to do its work, not evidence that four different models happen to have similar slopes. Disagreement, correspondingly, is now diagnostic of nuisance quality rather than of functional form.

	The most consequential change from the plug-in results is in the standard errors, and it runs in the opposite direction to the one the refitting bootstrap suggested. Under that bootstrap the neural network's standard errors were implausibly tight --- on the order of single dollars for several variables, \$7 for bedrooms and \$8 for overall quality --- because resampling propagated uncertainty through the model fitting but the smooth differentiable surface produced nearly identical finite differences across replicates. The influence-function standard errors for the same quantities are \$3,830 and \$2,394. The random forest shows the same correction: the high-neighborhood standard error rises from \$256 to \$24,636. Those earlier figures were not precision but the bootstrap's inability to see a bias that every replicate shared, which is the failure mode of Remark~\ref{rem:refit}, and their disappearance is the clearest single indication in the empirical results that the orthogonal score is doing what it was constructed to do.

	Read with those intervals attached, the bedroom variable is more equivocal than the plug-in made it look. The AME is $-$\$26,455 for linear regression, $-$\$9,212 for the random forest, $-$\$3,464 for XGBoost, and $-$\$12,639 for the neural network. All four are negative, consistent with the well-known hedonic finding that holding square footage fixed, more bedrooms means smaller rooms, but only the neural network estimate is significant at conventional levels and the linear estimate only marginally so. The plug-in reported all four as significant with standard errors between \$7 and \$3,042. The honest reading is that the sign is stable across architectures while this sample cannot pin down the magnitude, which is a weaker conclusion than the paper could previously state and a better supported one.
	
	\begin{sidewaystable}[htbp]
		\centering
		\small
		\begin{threeparttable}
			\caption{Ames Housing: Method Comparison. AME, SHAP, and PDP estimates for Sale Price (USD), full specification (A+B+C).}
			\label{tab:ames_housing_method_comparison}
			\begin{tabular}{lrrrrrrrrrrrr}
				\toprule
				& \multicolumn{3}{c}{Linear} & \multicolumn{3}{c}{Random Forest} & \multicolumn{3}{c}{XGBoost} & \multicolumn{3}{c}{Neural Net} \\
				\cmidrule(lr){2-4} \cmidrule(lr){5-7} \cmidrule(lr){8-10} \cmidrule(lr){11-13}
				& AME & SHAP & PDP & AME & SHAP & PDP & AME & SHAP & PDP & AME & SHAP & PDP \\
				\midrule
				GrLivArea & 68.05 & 3423.49 & 65.57 & 81.50 & 61.69 & 38.64 & 70.27 & -971.88 & 39.06 & 66.20 & -- & 58.62 \\
				BedroomAbvGr & -26454.60 & 145.82 & -6189.02 & -9212.02 & 90.87 & -1272.27 & -3464.31 & 158.44 & -4456.61 & -12638.60 & -- & -9294.43 \\
				FullBath & -10578.42 & 31.01 & -6287.42 & 7858.42 & 34.06 & 543.63 & 4761.27 & 54.86 & 193.04 & -5055.91 & -- & 2427.62 \\
				HalfBath & 1028.31 & -298.76 & -6967.90 & -10317.96 & 50.46 & -125.75 & -10906.00 & 30.53 & -1495.20 & -2175.44 & -- & -2723.69 \\
				\midrule
				OverallQual & 5436.28 & 1635.24 & 16465.22 & 16255.82 & 596.70 & 20776.14 & 14913.45 & 713.13 & 20130.92 & 8524.39 & -- & 10802.73 \\
				OverallCond & 2890.60 & 980.44 & 6311.48 & 9190.11 & 161.26 & 2030.44 & 12102.59 & -481.60 & 5071.47 & 7234.42 & -- & 6365.50 \\
				YearBuilt & -141.08 & -415.01 & 440.47 & 1323.39 & 288.29 & 209.05 & 1352.73 & 569.00 & 397.33 & 230.04 & -- & 421.27 \\
				YearRemodAdd & 1743.45 & 86.86 & 101.51 & -311.44 & 373.04 & 118.06 & -289.89 & 499.84 & 152.00 & 942.55 & -- & 561.16 \\
				\midrule
				LotArea & 23.02 & 267.44 & 0.61 & -12.51 & -183.50 & 1.33 & -10.83 & 216.28 & 2.38 & 10.04 & -- & 2.35 \\
				GarageArea & -102.78 & 93.22 & 38.68 & 90.33 & -884.31 & 23.77 & 92.29 & -704.59 & 28.98 & -6.11 & -- & 41.48 \\
				nbhd\_high & 89812.31 & 77.50 & 51433.02 & 43512.42 & -506.36 & 5058.64 & 33476.93 & 50.70 & 20307.34 & 35458.39 & -- & 12399.98 \\
				nbhd\_mid & 74336.31 & 266.38 & 7508.07 & 33767.43 & 40.68 & 5949.97 & 17468.12 & -134.71 & 9744.97 & 36951.86 & -- & 11392.44 \\
				\bottomrule
				\end{tabular}
				\begin{tablenotes}
				\footnotesize
				\item \textit{Notes:} AME = Average Marginal Effect (marginfx). SHAP = mean signed SHAP value (GradientExplainer). PDP = slope of partial dependence plot. Full specification (A+B+C) only. No standard errors reported for SHAP and PDP.
				\end{tablenotes}
		\end{threeparttable}
	\end{sidewaystable}
	
	Table~\ref{tab:ames_housing_method_comparison} brings the SHAP comparison into sharpest relief in the regression setting, where the units are dollars and the magnitudes are large enough to make the divergences concrete. SHAP values are unavailable for the neural network in this table, reflecting a limitation of the GradientExplainer implementation for this architecture. For the other three models the divergence between SHAP and AME is stark. The SHAP value for GrLivArea is \$3,423 for the linear model, which attributes to that variable an effect fifty times larger than the AME of \$68 --- a number that is not interpretable as a marginal effect in any standard sense. The SHAP values for garage area ($-$\$884), lot area ($-$\$183), and the high neighborhood indicator ($-$\$506) are negative for the random forest while the AME estimates are positive for all three. And as before, no standard errors attach to the SHAP column, so there is no basis on which to adjudicate any of these conflicts.

	The partial dependence column requires a different reading than it did under the plug-in, and the change is itself informative. A partial dependence slope is computed by perturbing a single full-sample fit and averaging, which makes it a plug-in average of exactly the kind the orthogonal correction is designed to repair; it carries no correction term and no standard error. Where the plug-in is nearly unbiased the two therefore continue to agree closely: on above-ground living area, the best supported and smoothest feature in the data, the AME is \$68 and the PDP slope \$66 for linear regression, \$66 and \$59 for the neural network. Where the regularization bias bites, they separate, and they separate in the direction the theory predicts --- toward zero for the uncorrected quantity. The bedroom AME of $-$\$26,455 for linear regression sits against a PDP slope of $-$\$6,189; the high-neighborhood AME of \$43,512 for the random forest sits against a PDP slope of \$5,059, and the mid-neighborhood AME of \$33,767 against \$5,950. These are the count and binary contrasts, where the fitted surface is flattest between levels and where the plug-in has the most to lose. The pattern that once read as mutual corroboration between two estimators now reads as a measurement of the bias separating them, which is the same quantity Simulation~2 isolates in the coverage of the refitting bootstrap. The Ames Housing results complete the sweep across outcome types: the framework produces dollar-denominated marginal effects with valid standard errors for a continuous regression outcome, using exactly the same estimation procedure as the binary classification examples, with no modification required.
	
	\section{Simulation}
	
	We evaluate the estimator through three simulation studies, which together separate what the window average recovers from what can be said about it. Simulation 1 examines the bias and root mean squared error of the uncorrected plug-in window average as a function of sample size, across three data generating processes and two outcome types. Simulation 2 asks whether the refitting bootstrap of Section~\ref{sec:inference} delivers valid intervals for that plug-in, comparing empirical coverage of nominal 95\% intervals against the target. Simulation 3 repeats both exercises for the debiased cross-fitted estimator of Section~\ref{sec:inference}, holding the designs, learners and windows fixed so that the two estimators differ only in the orthogonal correction. The three are best read in order: the first locates a bias, the second shows that resampling cannot see it, and the third removes it. All features $x_1, x_2, x_3, x_4$ are drawn independently from $N(0,1)$. Features $x_1$ and $x_2$ have nonzero true effects; $x_3$ and $x_4$ are pure noise with true AME of zero. Three data generating processes are considered. The linear DGP sets $\eta = 2x_1 + 3x_2$. The nonlinear DGP sets $\eta = 2x_1^2 + 3x_2$, so that the true AME of $x_1$ is zero in the regression case by symmetry of $N(0,1)$ --- $\mathbb{E}[4x_1] = 0$ --- but nonzero in the classification case where the sigmoid breaks that symmetry. The interaction DGP sets $\eta = 2x_1 + 3x_2 + 2x_1 x_2$, which preserves the marginal AMEs of $x_1$ and $x_2$ at 2.0 and 3.0 respectively in the regression case since $\mathbb{E}[2x_2] = \mathbb{E}[2x_1] = 0$. For regression, the outcome is $y = \eta + \varepsilon$ where $\varepsilon \sim N(0,1)$. For classification, $P(y=1) = \sigma(\eta)$ where $\sigma$ denotes the logistic function. True AMEs for all DGPs and outcome types are computed via Monte Carlo integration using $n = 1{,}000{,}000$ observations drawn from the true prediction function directly, treating the result as ground truth. Sample sizes of $n \in \{250, 500, 1\,000, 2\,500, 5\,000\}$ are evaluated over 500 Monte Carlo iterations for classification and 1{,}000 iterations for regression.
	
	The simulation study fits four model classes with hyperparameters chosen to balance computational tractability against representativeness of typical applied use. The random forest uses 100 trees with a maximum depth of 6, and XGBoost uses 100 boosting rounds with a maximum depth of 4 and a learning rate of 0.05 --- conservative settings that prevent overfitting on the small sample sizes considered while remaining flexible enough to capture the nonlinear and interaction structures in the data generating processes. The neural network is a feedforward architecture with two hidden layers of 32 and 16 units, trained for 50 epochs with batch size 64 and learning rate $10^{-3}$ using the Adam optimizer. For regression outcomes, the network includes the input normalization layer and target standardization wrapper described in Section~2 to ensure scale-invariant gradient updates. Logistic and linear regression are included as baselines and use no hyperparameter tuning. All model fits are performed independently on each Monte Carlo iteration with no warm starting between iterations, so the variance estimates reflect the full sampling distribution of the AME estimator including model fitting variability. The Monte Carlo ground truth AMEs are computed once at the start of the study using $n=1{,}000{,}000$ observations drawn from the true prediction function directly, with the resulting AMEs treated as known quantities for the bias and coverage calculations. Simulation 2 uses 1{,}000 bootstrap resamples per Monte Carlo iteration and 1{,}000 Monte Carlo iterations per cell for the linear, logistic, random forest and XGBoost learners. The neural network uses 200 resamples and 500 iterations: every resample is a fresh training run, which makes the network roughly nine tenths of the total cost of the study, and the calibration result the argument rests on is the forest's. Percentile intervals are used throughout, which is why the resample count is set well above the couple of hundred that would suffice for a standard error alone. Simulation 3 uses 1{,}000 Monte Carlo iterations per cell for every learner; it needs no resampling, since its standard errors come from the influence function.
	
	\subsection{Simulation 1: Recovering the Window Effect}
	
	\begin{table*}[htbp]
		\centering
		\small
		\begin{threeparttable}
			\caption{Simulation 1: AME Recovery --- Linear DGP (Classification). Mean bias and RMSE (in parentheses) across 500 Monte Carlo iterations.}
			\label{tab:sim1_classification_linear}
			\begin{tabular}{lrrrrrr}
				\toprule
				& & \multicolumn{5}{c}{Bias (RMSE)} \\
				\cmidrule(lr){3-7}
				Variable & True AME & $n=250$ & $n=500$ & $n=1,000$ & $n=2,500$ & $n=5,000$ \\
				\midrule
				\multicolumn{7}{l}{\textit{Panel: Logistic}} \\
				\quad x1 & 0.199 & -0.0058 & -0.0034 & -0.0014 & -0.0008 & -0.0006 \\
				&  & (0.0201) & (0.0146) & (0.0106) & (0.0064) & (0.0048) \\
				\quad x2 & 0.298 & -0.0109 & -0.0057 & -0.0026 & -0.0010 & -0.0005 \\
				&  & (0.0229) & (0.0153) & (0.0105) & (0.0066) & (0.0046) \\
				\quad x3 & 0.000 & 0.0012 & 0.0006 & 0.0006 & 0.0003 & 0.0002 \\
				&  & (0.0198) & (0.0136) & (0.0099) & (0.0063) & (0.0045) \\
				\quad x4 & 0.000 & 0.0008 & -0.0001 & -0.0008 & -0.0001 & 0.0003 \\
				&  & (0.0197) & (0.0136) & (0.0100) & (0.0062) & (0.0045) \\
				\midrule
				\multicolumn{7}{l}{\textit{Panel: Random Forest}} \\
				\quad x1 & 0.199 & 0.0126 & 0.0044 & -0.0018 & -0.0111 & -0.0157 \\
				&  & (0.0412) & (0.0239) & (0.0160) & (0.0142) & (0.0168) \\
				\quad x2 & 0.298 & 0.0199 & 0.0113 & -0.0008 & -0.0102 & -0.0158 \\
				&  & (0.0500) & (0.0309) & (0.0174) & (0.0143) & (0.0172) \\
				\quad x3 & 0.000 & 0.0007 & 0.0010 & 0.0005 & 0.0002 & 0.0002 \\
				&  & (0.0284) & (0.0171) & (0.0104) & (0.0056) & (0.0034) \\
				\quad x4 & 0.000 & -0.0008 & -0.0013 & -0.0012 & -0.0001 & 0.0002 \\
				&  & (0.0267) & (0.0177) & (0.0106) & (0.0055) & (0.0034) \\
				\midrule
				\multicolumn{7}{l}{\textit{Panel: XGBoost}} \\
				\quad x1 & 0.199 & 0.0590 & 0.0265 & 0.0107 & -0.0016 & -0.0049 \\
				&  & (0.0782) & (0.0401) & (0.0223) & (0.0108) & (0.0083) \\
				\quad x2 & 0.298 & 0.0643 & 0.0315 & 0.0085 & -0.0015 & -0.0046 \\
				&  & (0.0893) & (0.0490) & (0.0218) & (0.0116) & (0.0083) \\
				\quad x3 & 0.000 & 0.0008 & -0.0002 & 0.0006 & 0.0002 & 0.0002 \\
				&  & (0.0332) & (0.0193) & (0.0110) & (0.0056) & (0.0033) \\
				\quad x4 & 0.000 & 0.0008 & -0.0009 & -0.0006 & 0.0000 & 0.0002 \\
				&  & (0.0344) & (0.0198) & (0.0116) & (0.0054) & (0.0032) \\
				\midrule
				\multicolumn{7}{l}{\textit{Panel: Neural Net}} \\
				\quad x1 & 0.199 & -0.0078 & -0.0007 & 0.0005 & 0.0003 & -0.0002 \\
				&  & (0.0222) & (0.0156) & (0.0115) & (0.0074) & (0.0058) \\
				\quad x2 & 0.298 & -0.0162 & -0.0017 & 0.0004 & 0.0004 & 0.0002 \\
				&  & (0.0276) & (0.0174) & (0.0116) & (0.0073) & (0.0052) \\
				\quad x3 & 0.000 & 0.0014 & 0.0010 & 0.0007 & 0.0003 & 0.0004 \\
				&  & (0.0204) & (0.0144) & (0.0103) & (0.0069) & (0.0057) \\
				\quad x4 & 0.000 & 0.0009 & -0.0003 & -0.0007 & -0.0001 & 0.0004 \\
				&  & (0.0203) & (0.0143) & (0.0105) & (0.0070) & (0.0055) \\
				\bottomrule
				\end{tabular}
				\begin{tablenotes}
				\footnotesize
				\item \textit{Notes:} Mean bias and RMSE (in parentheses) computed over 500 Monte Carlo iterations. True AMEs computed via Monte Carlo integration with $n=1{,}000{,}000$ observations. Noise features x3 and x4 have true AME of zero.
				\end{tablenotes}
		\end{threeparttable}
	\end{table*}
	
	Table~\ref{tab:sim1_classification_linear} reports bias and RMSE for the linear classification DGP, where the true AMEs of $x_1$ and $x_2$ are 0.199 and 0.298 --- these are not 2.0 and 3.0 because the sigmoid compresses the scale. Logistic regression recovers the true AMEs with negligible bias at all sample sizes, and RMSE declines from around 0.020 at $n=250$ to 0.005 at $n=5{,}000$ at the expected $O(n^{-1/2})$ rate. The neural network behaves similarly, with small bias at $n=250$ that disappears by $n=500$. XGBoost shows larger bias at small samples --- around 0.06 for both $x_1$ and $x_2$ at $n=250$ --- reflecting the well-known tendency of gradient boosted models to underfit at small sample sizes, but bias falls sharply and RMSE converges by $n=2{,}500$. Random forests show a different pattern: bias is small at $n=250$ and $n=500$ but turns negative and persistent at $n=2{,}500$ and $n=5{,}000$, reaching around $-0.016$ for both signal features at the largest sample size. This reflects the regularization bias inherent in random forests --- the tree depth and subsampling constraints impose shrinkage that does not vanish with sample size at the default tuning settings. The noise features $x_3$ and $x_4$ have bias and RMSE near zero throughout for all model classes, confirming that the estimator does not spuriously attribute effects to irrelevant features.
	
	\begin{table*}[htbp]
		\centering
		\small
		\begin{threeparttable}
			\caption{Simulation 1: AME Recovery --- Nonlinear DGP (Classification). Mean bias and RMSE (in parentheses) across 500 Monte Carlo iterations.}
			\label{tab:sim1_classification_nonlinear}
			\begin{tabular}{lrrrrrr}
				\toprule
				& & \multicolumn{5}{c}{Bias (RMSE)} \\
				\cmidrule(lr){3-7}
				Variable & True AME & $n=250$ & $n=500$ & $n=1,000$ & $n=2,500$ & $n=5,000$ \\
				\midrule
				\multicolumn{7}{l}{\textit{Panel: Logistic}} \\
				\quad x1 & 0.000 & 0.0002 & 0.0015 & 0.0009 & 0.0002 & -0.0001 \\
				&  & (0.0230) & (0.0166) & (0.0123) & (0.0078) & (0.0055) \\
				\quad x2 & 0.298 & -0.0077 & -0.0042 & -0.0017 & -0.0007 & -0.0005 \\
				&  & (0.0215) & (0.0152) & (0.0109) & (0.0065) & (0.0048) \\
				\quad x3 & 0.000 & 0.0003 & 0.0001 & -0.0002 & -0.0001 & 0.0003 \\
				&  & (0.0229) & (0.0158) & (0.0112) & (0.0074) & (0.0052) \\
				\quad x4 & 0.000 & -0.0012 & 0.0005 & 0.0006 & -0.0001 & -0.0000 \\
				&  & (0.0252) & (0.0159) & (0.0108) & (0.0075) & (0.0051) \\
				\midrule
				\multicolumn{7}{l}{\textit{Panel: Random Forest}} \\
				\quad x1 & 0.000 & 0.0014 & 0.0001 & 0.0000 & -0.0003 & -0.0000 \\
				&  & (0.0407) & (0.0256) & (0.0173) & (0.0114) & (0.0082) \\
				\quad x2 & 0.298 & 0.0177 & 0.0060 & -0.0040 & -0.0120 & -0.0155 \\
				&  & (0.0521) & (0.0278) & (0.0182) & (0.0157) & (0.0171) \\
				\quad x3 & 0.000 & 0.0019 & -0.0000 & 0.0003 & 0.0001 & 0.0003 \\
				&  & (0.0285) & (0.0166) & (0.0100) & (0.0052) & (0.0031) \\
				\quad x4 & 0.000 & -0.0026 & -0.0007 & 0.0005 & 0.0000 & -0.0001 \\
				&  & (0.0295) & (0.0165) & (0.0099) & (0.0054) & (0.0031) \\
				\midrule
				\multicolumn{7}{l}{\textit{Panel: XGBoost}} \\
				\quad x1 & 0.000 & 0.0039 & 0.0003 & 0.0004 & -0.0000 & -0.0003 \\
				&  & (0.0549) & (0.0324) & (0.0209) & (0.0123) & (0.0079) \\
				\quad x2 & 0.298 & 0.0471 & 0.0221 & 0.0066 & -0.0028 & -0.0054 \\
				&  & (0.0787) & (0.0396) & (0.0208) & (0.0114) & (0.0087) \\
				\quad x3 & 0.000 & 0.0038 & 0.0007 & 0.0007 & 0.0002 & 0.0003 \\
				&  & (0.0307) & (0.0172) & (0.0108) & (0.0051) & (0.0028) \\
				\quad x4 & 0.000 & -0.0016 & 0.0010 & 0.0006 & 0.0001 & 0.0001 \\
				&  & (0.0329) & (0.0184) & (0.0105) & (0.0052) & (0.0027) \\
				\midrule
				\multicolumn{7}{l}{\textit{Panel: Neural Net}} \\
				\quad x1 & 0.000 & 0.0006 & 0.0015 & 0.0002 & 0.0001 & 0.0000 \\
				&  & (0.0234) & (0.0189) & (0.0140) & (0.0094) & (0.0069) \\
				\quad x2 & 0.298 & -0.0132 & -0.0026 & 0.0002 & 0.0005 & 0.0001 \\
				&  & (0.0249) & (0.0155) & (0.0112) & (0.0077) & (0.0064) \\
				\quad x3 & 0.000 & 0.0001 & 0.0002 & 0.0007 & 0.0003 & 0.0003 \\
				&  & (0.0215) & (0.0138) & (0.0101) & (0.0068) & (0.0053) \\
				\quad x4 & 0.000 & -0.0007 & 0.0004 & 0.0005 & -0.0001 & 0.0002 \\
				&  & (0.0230) & (0.0144) & (0.0100) & (0.0070) & (0.0056) \\
				\bottomrule
				\end{tabular}
				\begin{tablenotes}
				\footnotesize
				\item \textit{Notes:} Mean bias and RMSE (in parentheses) computed over 500 Monte Carlo iterations. True AMEs computed via Monte Carlo integration with $n=1{,}000{,}000$ observations. Noise features x3 and x4 have true AME of zero.
				\end{tablenotes}
		\end{threeparttable}
	\end{table*}
	
	Table~\ref{tab:sim1_classification_nonlinear} reports results for the nonlinear classification DGP, where $P(y=1) = \sigma(2x_1^2 + 3x_2)$. The true AME of $x_1$ is zero despite the quadratic term, because the sigmoid and quadratic effects cancel in expectation at these parameter values as confirmed by the Monte Carlo ground truth. All four model classes recover this near-zero AME accurately, with bias well below 0.005 across all sample sizes. The more important result is $x_2$, which has true AME of 0.298. Logistic regression misspecifies the DGP --- the true relationship is nonlinear in $x_1$ --- but still recovers the $x_2$ AME accurately because $x_2$ enters linearly and the misspecification in $x_1$ does not propagate to the $x_2$ marginal effect. This is an important practical result: AMEs from a misspecified parametric model can still be reliable for variables that enter the true DGP in a manner consistent with the model's functional form. The black-box models show the same pattern as the linear DGP: XGBoost has elevated RMSE at small samples, random forests develop negative bias at large samples, and neural networks perform comparably to logistic regression throughout.
	
	\begin{table*}[htbp]
		\centering
		\small
		\begin{threeparttable}
			\caption{Simulation 1: AME Recovery --- Interaction DGP (Classification). Mean bias and RMSE (in parentheses) across 500 Monte Carlo iterations.}
			\label{tab:sim1_classification_interaction}
			\begin{tabular}{lrrrrrr}
				\toprule
				& & \multicolumn{5}{c}{Bias (RMSE)} \\
				\cmidrule(lr){3-7}
				Variable & True AME & $n=250$ & $n=500$ & $n=1,000$ & $n=2,500$ & $n=5,000$ \\
				\midrule
				\multicolumn{7}{l}{\textit{Panel: Logistic}} \\
				\quad x1 & 0.199 & -0.0065 & -0.0042 & -0.0021 & -0.0008 & -0.0007 \\
				&  & (0.0219) & (0.0153) & (0.0112) & (0.0071) & (0.0051) \\
				\quad x2 & 0.298 & -0.0083 & -0.0035 & -0.0018 & -0.0008 & -0.0007 \\
				&  & (0.0222) & (0.0154) & (0.0102) & (0.0071) & (0.0050) \\
				\quad x3 & 0.000 & 0.0009 & -0.0002 & 0.0005 & 0.0003 & -0.0001 \\
				&  & (0.0206) & (0.0155) & (0.0107) & (0.0070) & (0.0046) \\
				\quad x4 & 0.000 & 0.0007 & -0.0000 & -0.0001 & 0.0001 & 0.0002 \\
				&  & (0.0212) & (0.0159) & (0.0106) & (0.0066) & (0.0047) \\
				\midrule
				\multicolumn{7}{l}{\textit{Panel: Random Forest}} \\
				\quad x1 & 0.199 & 0.0114 & 0.0051 & -0.0031 & -0.0110 & -0.0154 \\
				&  & (0.0423) & (0.0264) & (0.0169) & (0.0144) & (0.0167) \\
				\quad x2 & 0.298 & 0.0178 & 0.0104 & -0.0023 & -0.0101 & -0.0140 \\
				&  & (0.0498) & (0.0320) & (0.0200) & (0.0148) & (0.0158) \\
				\quad x3 & 0.000 & 0.0019 & 0.0004 & 0.0002 & 0.0002 & -0.0000 \\
				&  & (0.0276) & (0.0170) & (0.0104) & (0.0059) & (0.0034) \\
				\quad x4 & 0.000 & -0.0005 & -0.0011 & -0.0001 & 0.0001 & -0.0000 \\
				&  & (0.0267) & (0.0177) & (0.0111) & (0.0057) & (0.0033) \\
				\midrule
				\multicolumn{7}{l}{\textit{Panel: XGBoost}} \\
				\quad x1 & 0.199 & 0.0511 & 0.0225 & 0.0079 & -0.0014 & -0.0046 \\
				&  & (0.0737) & (0.0392) & (0.0219) & (0.0105) & (0.0082) \\
				\quad x2 & 0.298 & 0.0496 & 0.0228 & 0.0042 & -0.0021 & -0.0047 \\
				&  & (0.0796) & (0.0439) & (0.0227) & (0.0125) & (0.0091) \\
				\quad x3 & 0.000 & 0.0004 & 0.0005 & 0.0004 & 0.0000 & 0.0000 \\
				&  & (0.0310) & (0.0205) & (0.0111) & (0.0064) & (0.0035) \\
				\quad x4 & 0.000 & -0.0012 & -0.0015 & -0.0004 & 0.0001 & 0.0001 \\
				&  & (0.0338) & (0.0205) & (0.0118) & (0.0060) & (0.0037) \\
				\midrule
				\multicolumn{7}{l}{\textit{Panel: Neural Net}} \\
				\quad x1 & 0.199 & -0.0089 & -0.0013 & 0.0001 & 0.0000 & -0.0004 \\
				&  & (0.0233) & (0.0158) & (0.0117) & (0.0079) & (0.0061) \\
				\quad x2 & 0.298 & -0.0114 & -0.0002 & 0.0007 & 0.0005 & -0.0002 \\
				&  & (0.0257) & (0.0170) & (0.0131) & (0.0085) & (0.0067) \\
				\quad x3 & 0.000 & 0.0008 & 0.0004 & 0.0008 & 0.0001 & 0.0000 \\
				&  & (0.0193) & (0.0148) & (0.0099) & (0.0070) & (0.0053) \\
				\quad x4 & 0.000 & 0.0005 & -0.0002 & 0.0000 & -0.0001 & 0.0006 \\
				&  & (0.0206) & (0.0151) & (0.0103) & (0.0069) & (0.0057) \\
				\bottomrule
				\end{tabular}
				\begin{tablenotes}
				\footnotesize
				\item \textit{Notes:} Mean bias and RMSE (in parentheses) computed over 500 Monte Carlo iterations. True AMEs computed via Monte Carlo integration with $n=1{,}000{,}000$ observations. Noise features x3 and x4 have true AME of zero.
				\end{tablenotes}
		\end{threeparttable}
	\end{table*}
	
	Table~\ref{tab:sim1_classification_interaction} reports results for the interaction DGP, where $P(y=1) = \sigma(2x_1 + 3x_2 + 2x_1 x_2)$. The true AMEs of $x_1$ and $x_2$ remain 0.199 and 0.298 --- the same as the linear DGP --- because the interaction term $2x_1 x_2$ has zero expectation under $N(0,1)$ and the Monte Carlo ground truth confirms near-identical values. The key finding is that all four model classes recover these AMEs with bias and RMSE patterns essentially identical to the linear DGP, despite the true DGP including a multiplicative interaction. This is by design of the AME estimand: the average marginal effect averages over the joint distribution of the features, so the interaction term is absorbed into the average derivative rather than requiring explicit modeling. A logistic regression that omits the interaction term recovers the correct AME because the interaction's contribution to the marginal effect of $x_1$ is $\mathbb{E}[2x_2] = 0$ and similarly for $x_2$. This is an important robustness result --- AMEs are not sensitive to interaction misspecification when the interactions involve zero-mean features, which is the common case after centering.
	
	\begin{table*}[htbp]
		\centering
		\small
		\begin{threeparttable}
			\caption{Simulation 1: AME Recovery --- Linear DGP (Regression). Mean bias and RMSE (in parentheses) across 1,000 Monte Carlo iterations.}
			\label{tab:sim1_regression_linear}
			\begin{tabular}{lrrrrrr}
				\toprule
				& & \multicolumn{5}{c}{Bias (RMSE)} \\
				\cmidrule(lr){3-7}
				Variable & True AME & $n=250$ & $n=500$ & $n=1,000$ & $n=2,500$ & $n=5,000$ \\
				\midrule
				\multicolumn{7}{l}{\textit{Panel: Linear}} \\
				\quad x1 & 2.000 & -0.0017 & -0.0001 & -0.0000 & -0.0006 & -0.0003 \\
				&  & (0.0659) & (0.0445) & (0.0322) & (0.0207) & (0.0137) \\
				\quad x2 & 3.000 & -0.0043 & -0.0017 & 0.0007 & -0.0004 & -0.0003 \\
				&  & (0.0636) & (0.0453) & (0.0316) & (0.0197) & (0.0140) \\
				\quad x3 & 0.000 & -0.0001 & -0.0006 & -0.0005 & 0.0012 & -0.0003 \\
				&  & (0.0650) & (0.0457) & (0.0312) & (0.0199) & (0.0140) \\
				\quad x4 & 0.000 & -0.0034 & 0.0006 & -0.0002 & -0.0003 & -0.0001 \\
				&  & (0.0642) & (0.0452) & (0.0315) & (0.0196) & (0.0139) \\
				\midrule
				\multicolumn{7}{l}{\textit{Panel: Random Forest}} \\
				\quad x1 & 2.000 & -0.0973 & -0.0332 & -0.0387 & -0.0703 & -0.0911 \\
				&  & (0.2064) & (0.1189) & (0.0850) & (0.0824) & (0.0956) \\
				\quad x2 & 3.000 & -0.0215 & -0.0105 & -0.0433 & -0.0797 & -0.0953 \\
				&  & (0.2338) & (0.1456) & (0.0997) & (0.0958) & (0.1025) \\
				\quad x3 & 0.000 & -0.0015 & -0.0005 & -0.0005 & 0.0001 & 0.0000 \\
				&  & (0.0562) & (0.0290) & (0.0137) & (0.0036) & (0.0011) \\
				\quad x4 & 0.000 & -0.0021 & -0.0010 & -0.0005 & -0.0002 & 0.0000 \\
				&  & (0.0574) & (0.0297) & (0.0138) & (0.0037) & (0.0011) \\
				\midrule
				\multicolumn{7}{l}{\textit{Panel: XGBoost}} \\
				\quad x1 & 2.000 & 0.4114 & 0.1780 & 0.0366 & -0.0337 & -0.0467 \\
				&  & (0.4644) & (0.2227) & (0.0917) & (0.0554) & (0.0547) \\
				\quad x2 & 3.000 & 0.3122 & 0.0948 & -0.0172 & -0.0558 & -0.0601 \\
				&  & (0.4103) & (0.1884) & (0.0947) & (0.0754) & (0.0677) \\
				\quad x3 & 0.000 & 0.0007 & 0.0016 & -0.0003 & 0.0011 & -0.0001 \\
				&  & (0.1090) & (0.0550) & (0.0313) & (0.0153) & (0.0083) \\
				\quad x4 & 0.000 & -0.0080 & 0.0002 & 0.0011 & -0.0006 & 0.0001 \\
				&  & (0.1077) & (0.0597) & (0.0314) & (0.0149) & (0.0081) \\
				\midrule
				\multicolumn{7}{l}{\textit{Panel: Neural Net}} \\
				\quad x1 & 2.000 & -0.0659 & -0.0208 & -0.0063 & -0.0021 & -0.0008 \\
				&  & (0.1161) & (0.0544) & (0.0344) & (0.0293) & (0.0329) \\
				\quad x2 & 3.000 & -0.1231 & -0.0349 & -0.0102 & -0.0031 & -0.0006 \\
				&  & (0.1765) & (0.0638) & (0.0361) & (0.0324) & (0.0372) \\
				\quad x3 & 0.000 & 0.0014 & -0.0006 & -0.0003 & 0.0009 & -0.0006 \\
				&  & (0.0788) & (0.0484) & (0.0330) & (0.0250) & (0.0280) \\
				\quad x4 & 0.000 & -0.0002 & 0.0008 & -0.0003 & -0.0004 & -0.0017 \\
				&  & (0.0752) & (0.0496) & (0.0336) & (0.0247) & (0.0287) \\
				\bottomrule
				\end{tabular}
				\begin{tablenotes}
				\footnotesize
				\item \textit{Notes:} Mean bias and RMSE (in parentheses) computed over 1,000 Monte Carlo iterations. True AMEs computed via Monte Carlo integration with $n=1{,}000{,}000$ observations. Noise features x3 and x4 have true AME of zero.
				\end{tablenotes}
		\end{threeparttable}
	\end{table*}
	
	Table~\ref{tab:sim1_regression_linear} turns to the regression setting with the linear DGP, $y = 2x_1 + 3x_2 + \varepsilon$. Linear regression recovers the true AMEs of 2.0 and 3.0 with negligible bias across all sample sizes, and RMSE declines from around 0.066 at $n=250$ to 0.014 at $n=5{,}000$, consistent with $O(n^{-1/2})$ convergence. The regression setting reveals the regularization bias of random forests more starkly than the classification setting: bias reaches $-0.091$ for $x_1$ and $-0.095$ for $x_2$ at $n=5{,}000$, and RMSE is not declining at the largest sample sizes, indicating that the bias is not vanishing with $n$ at default tuning. This is a known property of random forests in the regression case --- the ensemble averaging and subsampling impose shrinkage toward zero that is not corrected by the hyperparameter search at these sample sizes. XGBoost shows large positive bias at small samples --- 0.41 for $x_1$ at $n=250$ --- but converges quickly, with bias falling below 0.05 by $n=2{,}500$. The neural network shows moderate negative bias at small samples that falls toward zero by $n=2{,}500$, though RMSE does not decline monotonically at the largest sample sizes, suggesting residual variance from the optimization landscape. The noise features have RMSE near zero for linear regression and declining rapidly for the black-box models, confirming that the estimator correctly attributes near-zero effects to irrelevant features in the regression setting.
	
	\begin{table*}[htbp]
		\centering
		\small
		\begin{threeparttable}
			\caption{Simulation 1: AME Recovery --- Nonlinear DGP (Regression). Mean bias and RMSE (in parentheses) across 1,000 Monte Carlo iterations.}
			\label{tab:sim1_regression_nonlinear}
			\begin{tabular}{lrrrrrr}
				\toprule
				& & \multicolumn{5}{c}{Bias (RMSE)} \\
				\cmidrule(lr){3-7}
				Variable & True AME & $n=250$ & $n=500$ & $n=1,000$ & $n=2,500$ & $n=5,000$ \\
				\midrule
				\multicolumn{7}{l}{\textit{Panel: Linear}} \\
				\quad x1 & 0.000 & 0.0241 & 0.0261 & 0.0148 & 0.0048 & 0.0015 \\
				&  & (0.3868) & (0.2820) & (0.1966) & (0.1272) & (0.0904) \\
				\quad x2 & 3.000 & -0.0145 & -0.0041 & -0.0012 & -0.0002 & 0.0001 \\
				&  & (0.1921) & (0.1276) & (0.0930) & (0.0572) & (0.0417) \\
				\quad x3 & 0.000 & -0.0025 & 0.0001 & -0.0018 & -0.0020 & -0.0009 \\
				&  & (0.1889) & (0.1343) & (0.0941) & (0.0576) & (0.0414) \\
				\quad x4 & 0.000 & -0.0114 & -0.0003 & 0.0023 & 0.0006 & -0.0018 \\
				&  & (0.1857) & (0.1328) & (0.0927) & (0.0594) & (0.0430) \\
				\midrule
				\multicolumn{7}{l}{\textit{Panel: Random Forest}} \\
				\quad x1 & 0.000 & 0.0071 & 0.0097 & 0.0053 & 0.0006 & 0.0007 \\
				&  & (0.3433) & (0.2547) & (0.1836) & (0.1185) & (0.0903) \\
				\quad x2 & 3.000 & 0.0142 & -0.1152 & -0.1701 & -0.1867 & -0.1849 \\
				&  & (0.2704) & (0.2008) & (0.2004) & (0.1974) & (0.1916) \\
				\quad x3 & 0.000 & 0.0001 & 0.0008 & 0.0003 & -0.0002 & 0.0000 \\
				&  & (0.0549) & (0.0274) & (0.0138) & (0.0055) & (0.0027) \\
				\quad x4 & 0.000 & -0.0009 & -0.0007 & -0.0001 & -0.0000 & -0.0000 \\
				&  & (0.0573) & (0.0278) & (0.0138) & (0.0054) & (0.0025) \\
				\midrule
				\multicolumn{7}{l}{\textit{Panel: XGBoost}} \\
				\quad x1 & 0.000 & 0.0552 & 0.0038 & -0.0117 & -0.0137 & -0.0057 \\
				&  & (0.4662) & (0.2744) & (0.1829) & (0.1040) & (0.0735) \\
				\quad x2 & 3.000 & 0.2725 & 0.0816 & -0.0183 & -0.0702 & -0.0794 \\
				&  & (0.3937) & (0.1981) & (0.1058) & (0.0885) & (0.0862) \\
				\quad x3 & 0.000 & -0.0040 & -0.0022 & -0.0020 & -0.0011 & -0.0006 \\
				&  & (0.1124) & (0.0704) & (0.0446) & (0.0191) & (0.0097) \\
				\quad x4 & 0.000 & -0.0052 & -0.0028 & -0.0021 & 0.0003 & 0.0006 \\
				&  & (0.1121) & (0.0636) & (0.0381) & (0.0189) & (0.0096) \\
				\midrule
				\multicolumn{7}{l}{\textit{Panel: Neural Net}} \\
				\quad x1 & 0.000 & 0.0201 & 0.0178 & 0.0065 & -0.0003 & -0.0013 \\
				&  & (0.3086) & (0.1925) & (0.1390) & (0.0911) & (0.0760) \\
				\quad x2 & 3.000 & -0.1545 & -0.0202 & -0.0150 & -0.0098 & -0.0022 \\
				&  & (0.2601) & (0.0710) & (0.0461) & (0.0343) & (0.0395) \\
				\quad x3 & 0.000 & 0.0017 & 0.0004 & 0.0007 & 0.0009 & -0.0004 \\
				&  & (0.1419) & (0.0634) & (0.0400) & (0.0293) & (0.0321) \\
				\quad x4 & 0.000 & -0.0028 & -0.0001 & -0.0013 & -0.0004 & -0.0016 \\
				&  & (0.1410) & (0.0626) & (0.0406) & (0.0281) & (0.0323) \\
				\bottomrule
				\end{tabular}
				\begin{tablenotes}
				\footnotesize
				\item \textit{Notes:} Mean bias and RMSE (in parentheses) computed over 1,000 Monte Carlo iterations. True AMEs computed via Monte Carlo integration with $n=1{,}000{,}000$ observations. Noise features x3 and x4 have true AME of zero.
				\end{tablenotes}
		\end{threeparttable}
	\end{table*}
	
	Table~\ref{tab:sim1_regression_nonlinear} reports results for the nonlinear regression DGP, $y = 2x_1^2 + 3x_2 + \varepsilon$, where the true AME of $x_1$ is zero by symmetry of $N(0,1)$ and the true AME of $x_2$ is 3.0. Linear regression recovers the $x_2$ AME accurately with RMSE declining at the standard rate, and correctly attributes a near-zero AME to $x_1$ despite misspecifying its functional form --- the same robustness result observed in the classification case. However, the RMSE for $x_1$ under linear regression is substantially larger than in the linear DGP --- 0.387 at $n=250$ versus 0.066 --- reflecting increased estimation variance from the misspecified model. The most striking result is random forests: the bias for $x_2$ reaches $-0.185$ at $n=5{,}000$ and is not declining, indicating that the regularization bias is amplified in the nonlinear setting where the forest must capture the quadratic relationship in $x_1$ while simultaneously estimating the linear effect of $x_2$. XGBoost and the neural network both show elevated RMSE relative to the linear DGP but declining bias, suggesting that these models are better able to adapt to the nonlinear DGP with sufficient data. The $x_1$ AME of zero is recovered accurately by all model classes at moderate sample sizes, which is a reassuring result for practitioners who want to test whether nonlinear features have nonzero average effects.
	
	\begin{table*}[htbp]
		\centering
		\small
		\begin{threeparttable}
			\caption{Simulation 1: AME Recovery --- Interaction DGP (Regression). Mean bias and RMSE (in parentheses) across 1,000 Monte Carlo iterations.}
			\label{tab:sim1_regression_interaction}
			\begin{tabular}{lrrrrrr}
				\toprule
				& & \multicolumn{5}{c}{Bias (RMSE)} \\
				\cmidrule(lr){3-7}
				Variable & True AME & $n=250$ & $n=500$ & $n=1,000$ & $n=2,500$ & $n=5,000$ \\
				\midrule
				\multicolumn{7}{l}{\textit{Panel: Linear}} \\
				\quad x1 & 2.000 & -0.0041 & 0.0005 & 0.0000 & 0.0006 & 0.0011 \\
				&  & (0.2284) & (0.1531) & (0.1129) & (0.0715) & (0.0495) \\
				\quad x2 & 3.000 & 0.0040 & 0.0087 & 0.0075 & 0.0012 & 0.0009 \\
				&  & (0.2194) & (0.1627) & (0.1135) & (0.0727) & (0.0514) \\
				\quad x3 & 0.000 & -0.0114 & -0.0088 & -0.0054 & 0.0010 & -0.0009 \\
				&  & (0.1388) & (0.1002) & (0.0725) & (0.0450) & (0.0329) \\
				\quad x4 & 0.000 & -0.0073 & 0.0004 & -0.0006 & 0.0013 & 0.0005 \\
				&  & (0.1443) & (0.1018) & (0.0690) & (0.0453) & (0.0321) \\
				\midrule
				\multicolumn{7}{l}{\textit{Panel: Random Forest}} \\
				\quad x1 & 2.000 & -0.0755 & -0.0310 & -0.0433 & -0.0724 & -0.0804 \\
				&  & (0.2566) & (0.1663) & (0.1169) & (0.0979) & (0.0925) \\
				\quad x2 & 3.000 & -0.0634 & -0.0690 & -0.0922 & -0.1112 & -0.1143 \\
				&  & (0.2806) & (0.1913) & (0.1497) & (0.1306) & (0.1244) \\
				\quad x3 & 0.000 & -0.0031 & -0.0004 & -0.0009 & 0.0003 & -0.0001 \\
				&  & (0.0612) & (0.0352) & (0.0182) & (0.0067) & (0.0029) \\
				\quad x4 & 0.000 & -0.0023 & 0.0004 & -0.0001 & -0.0002 & 0.0000 \\
				&  & (0.0665) & (0.0349) & (0.0184) & (0.0069) & (0.0030) \\
				\midrule
				\multicolumn{7}{l}{\textit{Panel: XGBoost}} \\
				\quad x1 & 2.000 & 0.3184 & 0.1016 & -0.0009 & -0.0406 & -0.0429 \\
				&  & (0.4539) & (0.2184) & (0.1163) & (0.0758) & (0.0602) \\
				\quad x2 & 3.000 & 0.3672 & 0.1135 & -0.0083 & -0.0572 & -0.0589 \\
				&  & (0.4920) & (0.2373) & (0.1227) & (0.0898) & (0.0746) \\
				\quad x3 & 0.000 & -0.0031 & 0.0011 & -0.0017 & 0.0005 & -0.0001 \\
				&  & (0.1040) & (0.0551) & (0.0319) & (0.0151) & (0.0078) \\
				\quad x4 & 0.000 & -0.0082 & 0.0003 & -0.0003 & -0.0004 & -0.0001 \\
				&  & (0.1028) & (0.0564) & (0.0314) & (0.0149) & (0.0080) \\
				\midrule
				\multicolumn{7}{l}{\textit{Panel: Neural Net}} \\
				\quad x1 & 2.000 & -0.0842 & -0.0083 & -0.0056 & 0.0003 & 0.0022 \\
				&  & (0.2380) & (0.1114) & (0.0779) & (0.0573) & (0.0536) \\
				\quad x2 & 3.000 & -0.1231 & -0.0068 & -0.0052 & 0.0004 & 0.0026 \\
				&  & (0.2467) & (0.1120) & (0.0783) & (0.0579) & (0.0533) \\
				\quad x3 & 0.000 & -0.0026 & -0.0016 & -0.0024 & 0.0004 & -0.0004 \\
				&  & (0.0986) & (0.0551) & (0.0383) & (0.0293) & (0.0316) \\
				\quad x4 & 0.000 & -0.0079 & 0.0015 & 0.0002 & -0.0005 & -0.0019 \\
				&  & (0.0981) & (0.0565) & (0.0387) & (0.0284) & (0.0331) \\
				\bottomrule
				\end{tabular}
				\begin{tablenotes}
				\footnotesize
				\item \textit{Notes:} Mean bias and RMSE (in parentheses) computed over 1,000 Monte Carlo iterations. True AMEs computed via Monte Carlo integration with $n=1{,}000{,}000$ observations. Noise features x3 and x4 have true AME of zero.
				\end{tablenotes}
		\end{threeparttable}
	\end{table*}
	
	Table~\ref{tab:sim1_regression_interaction} reports results for the interaction regression DGP, $y = 2x_1 + 3x_2 + 2x_1 x_2 + \varepsilon$, with true AMEs of 2.0 and 3.0 for $x_1$ and $x_2$. Linear regression recovers these AMEs accurately despite omitting the interaction term, confirming the robustness result from the classification case: the interaction contribution to the marginal effect of $x_1$ is $\mathbb{E}[2x_2] = 0$ and similarly for $x_2$, so the misspecified linear model nevertheless recovers the correct AME. The RMSE for the noise features under linear regression is elevated relative to the linear DGP --- around 0.14 at $n=250$ --- because the interaction term inflates the residual variance and increases estimation uncertainty for all coefficients. Random forests again show persistent negative bias that grows rather than shrinks with sample size, reaching $-0.080$ and $-0.114$ for $x_1$ and $x_2$ respectively at $n=5{,}000$. XGBoost shows the same large positive bias at small samples followed by convergence, but the bias does not fully vanish at $n=5{,}000$ --- around $-0.043$ for $x_1$ and $-0.059$ for $x_2$ --- suggesting that the interaction structure requires larger samples for the gradient boosting to fully adapt. The neural network performs well, with bias below 0.003 for both signal features at $n=2{,}500$ and $n=5{,}000$, and RMSE declining steadily. Taken together the six bias and RMSE tables establish that the AME estimator is consistent for logistic regression, neural networks, and XGBoost across all three DGPs and both outcome types, with random forests showing persistent regularization bias in the regression setting that practitioners should be aware of when reporting AMEs from default-tuned forest models.
	
	\subsection{Simulation 2: The Refitting Bootstrap}
	
	\begin{table*}[htbp]
		\centering
		\footnotesize
		\begin{threeparttable}
			\caption{Simulation 2: Bootstrap SE Calibration --- Linear DGP (Classification). Nominal coverage target: 95\%. Coverage rate and mean 95\% CI width (in parentheses) across 500 Monte Carlo iterations.}
			\label{tab:sim2_calibration_linear}
			\begin{tabular}{lrrrr}
				\toprule
				& & \multicolumn{3}{c}{Coverage (CI Width)} \\
				\cmidrule(lr){3-5}
				Variable & True AME & $n=250$ & $n=1,000$ & $n=5,000$ \\
				\midrule
				\multicolumn{5}{l}{\textit{Panel: Logistic}} \\
				\quad x1 & 0.199 & 0.881 & 0.878 & 0.886 \\
				&  & (0.0652) & (0.0334) & (0.0150) \\
				\quad x2 & 0.298 & 0.726 & 0.766 & 0.781 \\
				&  & (0.0505) & (0.0257) & (0.0115) \\
				\quad x3 & 0.000 & 0.944 & 0.939 & 0.946 \\
				&  & (0.0752) & (0.0386) & (0.0174) \\
				\quad x4 & 0.000 & 0.933 & 0.954 & 0.944 \\
				&  & (0.0754) & (0.0386) & (0.0174) \\
				\midrule
				\multicolumn{5}{l}{\textit{Panel: Random Forest}} \\
				\quad x1 & 0.199 & 0.968 & 0.933 & 0.186 \\
				&  & (0.1343) & (0.0553) & (0.0205) \\
				\quad x2 & 0.298 & 0.961 & 0.958 & 0.268 \\
				&  & (0.1500) & (0.0614) & (0.0219) \\
				\quad x3 & 0.000 & 0.962 & 0.942 & 0.949 \\
				&  & (0.0976) & (0.0386) & (0.0132) \\
				\quad x4 & 0.000 & 0.969 & 0.964 & 0.953 \\
				&  & (0.0978) & (0.0386) & (0.0132) \\
				\midrule
				\multicolumn{5}{l}{\textit{Panel: XGBoost}} \\
				\quad x1 & 0.199 & 0.926 & 0.947 & 0.877 \\
				&  & (0.1864) & (0.0702) & (0.0230) \\
				\quad x2 & 0.298 & 0.908 & 0.954 & 0.910 \\
				&  & (0.2136) & (0.0794) & (0.0239) \\
				\quad x3 & 0.000 & 0.976 & 0.961 & 0.960 \\
				&  & (0.1139) & (0.0436) & (0.0141) \\
				\quad x4 & 0.000 & 0.973 & 0.969 & 0.959 \\
				&  & (0.1149) & (0.0436) & (0.0142) \\
				\midrule
				\multicolumn{5}{l}{\textit{Panel: Neural Net}} \\
				\quad x1 & 0.199 & 0.870 & 0.870 & 0.958 \\
				&  & (0.0679) & (0.0351) & (0.0203) \\
				\quad x2 & 0.298 & 0.658 & 0.774 & 0.894 \\
				&  & (0.0586) & (0.0287) & (0.0154) \\
				\quad x3 & 0.000 & 0.936 & 0.936 & 0.986 \\
				&  & (0.0764) & (0.0391) & (0.0225) \\
				\quad x4 & 0.000 & 0.938 & 0.936 & 0.982 \\
				&  & (0.0772) & (0.0392) & (0.0225) \\
				\bottomrule
				\end{tabular}
				\begin{tablenotes}
				\footnotesize
				\item \textit{Notes:} Coverage rate and mean 95\% CI width (in parentheses). Intervals are bootstrap percentile intervals. Logistic, random forest and XGBoost use 1{,}000 Monte Carlo iterations and 1{,}000 resamples; the neural network uses 500 and 200, since every one of its resamples is a fresh training run. True AMEs computed via Monte Carlo integration with $n=1{,}000{,}000$ observations. Noise features x3 and x4 have true AME of zero.
				\end{tablenotes}
		\end{threeparttable}
	\end{table*}
	
	Table~\ref{tab:sim2_calibration_linear} reports bootstrap standard error calibration for the linear classification DGP, evaluating empirical coverage of nominal 95\% confidence intervals at $n \in \{250, 1{\,}000, 5{\,}000\}$ using 200 bootstrap replicates. Logistic regression achieves coverage between 0.896 and 0.950 across all variables and sample sizes, with the slight undercoverage at $n=250$ for $x_2$ reflecting finite-sample bias rather than bootstrap miscalibration --- the bootstrap interval is correctly sized but the point estimate is slightly off-center. Coverage improves toward the nominal level as sample size increases and is well-calibrated by $n=5{,}000$. XGBoost shows modest overcoverage for the noise features at all sample sizes --- around 0.978 to 0.986 --- reflecting the wide bootstrap intervals produced by the model's high variance at small samples; the intervals are conservative rather than anti-conservative, which is the safer failure mode for inference. The neural network shows undercoverage for the noise features at $n=250$ (0.866 and 0.854) that recovers toward and above the nominal level by $n=5{,}000$, with the overcoverage at large samples (0.990 and 0.994) consistent with the smooth differentiable surface producing stable finite differences across bootstrap replicates. Random forests show a more concerning pattern: coverage collapses dramatically at $n=5{,}000$, falling to 0.130 for $x_1$, 0.274 for $x_2$, and 0.392 for the noise features. This is a direct consequence of the regularization bias documented in Simulation 1 --- the bootstrap intervals shrink at the expected rate but they shrink around a biased point estimate, so as the intervals narrow they fail to cover the true AME. The intervals at $n=5{,}000$ have mean width 0.0165 for $x_1$ but the bias is around $-0.016$, so the interval is roughly the width of the bias itself and coverage of the true value becomes nearly impossible. This is the first failure mode identified in Remark~\ref{rem:refit}: the refitting bootstrap is consistent for the variability of the plug-in average around its own expectation, but that expectation is displaced from the true AME by the regularization bias $\mathbb{E}[\alpha_h(\hat{f} - f)]$, which resampling cannot see because every replicate recenters at a similarly regularized fit. It is exactly the first-order term that the orthogonal correction of Section~2.4 is constructed to remove.
	
	\begin{table*}[htbp]
		\centering
		\footnotesize
		\begin{threeparttable}
			\caption{Simulation 2: Bootstrap SE Calibration --- Linear DGP (Regression). Nominal coverage target: 95\%. Coverage rate and mean 95\% CI width (in parentheses) across 500 Monte Carlo iterations.}
			\label{tab:sim2_calibration_regression_linear}
			\begin{tabular}{lrrrr}
				\toprule
				& & \multicolumn{3}{c}{Coverage (CI Width)} \\
				\cmidrule(lr){3-5}
				Variable & True AME & $n=250$ & $n=1,000$ & $n=5,000$ \\
				\midrule
				\multicolumn{5}{l}{\textit{Panel: Linear}} \\
				\quad x1 & 2.000 & 0.928 & 0.945 & 0.954 \\
				&  & (0.2478) & (0.1235) & (0.0553) \\
				\quad x2 & 3.000 & 0.955 & 0.948 & 0.956 \\
				&  & (0.2489) & (0.1239) & (0.0552) \\
				\quad x3 & 0.000 & 0.951 & 0.947 & 0.948 \\
				&  & (0.2495) & (0.1234) & (0.0552) \\
				\quad x4 & 0.000 & 0.949 & 0.948 & 0.957 \\
				&  & (0.2489) & (0.1234) & (0.0552) \\
				\midrule
				\multicolumn{5}{l}{\textit{Panel: Random Forest}} \\
				\quad x1 & 2.000 & 0.870 & 0.886 & 0.092 \\
				&  & (0.7447) & (0.2990) & (0.1170) \\
				\quad x2 & 3.000 & 0.967 & 0.962 & 0.323 \\
				&  & (0.9327) & (0.3632) & (0.1441) \\
				\quad x3 & 0.000 & 0.996 & 0.996 & 0.999 \\
				&  & (0.2555) & (0.0707) & (0.0082) \\
				\quad x4 & 0.000 & 0.992 & 0.995 & 1.000 \\
				&  & (0.2566) & (0.0708) & (0.0083) \\
				\midrule
				\multicolumn{5}{l}{\textit{Panel: XGBoost}} \\
				\quad x1 & 2.000 & 0.849 & 0.982 & 0.521 \\
				&  & (0.8174) & (0.3009) & (0.1041) \\
				\quad x2 & 3.000 & 0.917 & 0.990 & 0.503 \\
				&  & (0.9709) & (0.3552) & (0.1184) \\
				\quad x3 & 0.000 & 0.988 & 0.969 & 0.971 \\
				&  & (0.3863) & (0.1217) & (0.0382) \\
				\quad x4 & 0.000 & 0.978 & 0.975 & 0.969 \\
				&  & (0.3778) & (0.1210) & (0.0382) \\
				\midrule
				\multicolumn{5}{l}{\textit{Panel: Neural Net}} \\
				\quad x1 & 2.000 & 0.950 & 0.978 & 1.000 \\
				&  & (0.3442) & (0.1497) & (0.1938) \\
				\quad x2 & 3.000 & 0.936 & 0.974 & 1.000 \\
				&  & (0.4000) & (0.1537) & (0.2004) \\
				\quad x3 & 0.000 & 0.972 & 0.970 & 1.000 \\
				&  & (0.3075) & (0.1471) & (0.1887) \\
				\quad x4 & 0.000 & 0.968 & 0.972 & 1.000 \\
				&  & (0.3116) & (0.1474) & (0.1869) \\
				\bottomrule
				\end{tabular}
				\begin{tablenotes}
				\footnotesize
				\item \textit{Notes:} Coverage rate and mean 95\% CI width (in parentheses). Intervals are bootstrap percentile intervals. Logistic, random forest and XGBoost use 1{,}000 Monte Carlo iterations and 1{,}000 resamples; the neural network uses 500 and 200, since every one of its resamples is a fresh training run. True AMEs computed via Monte Carlo integration with $n=1{,}000{,}000$ observations. Noise features x3 and x4 have true AME of zero.
				\end{tablenotes}
		\end{threeparttable}
	\end{table*}
	
	Table~\ref{tab:sim2_calibration_regression_linear} reports the corresponding calibration results for the linear regression DGP, where the magnitudes of the AMEs are larger and the bias-coverage interaction becomes even more visible. Linear regression achieves coverage between 0.918 and 0.952 across all variables and sample sizes, well-calibrated throughout. The neural network performs well at small and moderate samples, with coverage approaching the nominal level by $n=1{,}000$, before drifting upward to exactly 1.000 at $n=5{,}000$ for all four variables --- the same conservatism observed in the classification case, amplified in the regression setting where the smooth prediction surface produces highly stable finite differences. XGBoost shows pronounced undercoverage for the signal features at $n=250$ --- 0.526 for $x_1$ and 0.752 for $x_2$ --- driven by the large positive bias documented in Simulation 1 at small samples, but coverage recovers toward and above the nominal level by $n=1{,}000$ and is well-calibrated at $n=5{,}000$. The most striking result is again random forests: coverage at $n=5{,}000$ falls to 0.104 for $x_1$ and 0.204 for $x_2$, the worst calibration in either table. The mechanism is identical to the classification case --- the bootstrap intervals narrow at the expected rate but the bias does not vanish, so the intervals fail to cover the true AME at large samples. The noise features under random forests show coverage of around 0.6 across all sample sizes, indicating that even where the bias is small, the bootstrap intervals are systematically too narrow to achieve nominal coverage. Taken together, the two calibration tables show that the refitting bootstrap applied to the uncorrected plug-in tracks sampling variability adequately when the learner's regularization bias is negligible at the sample sizes considered --- logistic and linear regression, the neural network, and XGBoost beyond small samples --- and fails without warning when the bias term dominates, as it does for random forests at large $n$. The tables should therefore be read as a controlled demonstration of the mechanism analyzed in Remark~\ref{rem:refit}, not as a validation of the plug-in bootstrap as an inferential procedure: the failure is not specific to forests but to any learner whose $L^2$ bias shrinks more slowly than its intervals, and forests at default tuning are simply the class for which this occurs within the sample sizes examined. In the designs of this section the covariate density is known, so the representer $\alpha_h$ is available in closed form and the debiased estimator of Section~2.4 requires no representer estimation at all. The forests here are moreover depth-capped, and a depth-capped forest converges to a depth-limited approximation of the regression function rather than to the regression function itself; Proposition~\ref{prop:robust} applies verbatim to this case, and asserts that coverage of the true window effect requires no consistency from the forest --- only the known representer --- with the forest's approximation error paid in interval width rather than in location.
	\subsection{Simulation 3: The Debiased Estimator}
	
	Simulation 3 repeats the design of Simulation 2 for the debiased cross-fitted estimator of Section~\ref{sec:inference}. The covariates are independent standard normals, so the Riesz representer of Lemma~\ref{lem:riesz} is available in closed form, $\alpha_h(u) = e^{-h^2/2}\sinh(h u_j)/h$, and the estimator requires no representer estimation at all; a second arm reported at the end of this subsection replaces it with a sieve estimate, to show what is lost when the covariate density is unknown. Cross-fitting uses $K=5$ folds. The data generating processes, learners, hyperparameters, sample sizes, window $h$ and Monte Carlo ground truth are all held at the values used in the two preceding studies, so any difference between Simulation 3 and its predecessors is attributable to the orthogonal correction and to nothing else. Trimming is inactive throughout because the covariate support is unbounded and the trimming weight of Section~2.1 is therefore identically one.
	
	\begin{table*}[htbp]
		\centering
		\small
		\begin{threeparttable}
			\caption{Simulation 3: debiased estimator, linear regression DGP. Bias and RMSE (in parentheses) by sample size.}
			\label{tab:sim3_bias_regression_linear}
			\begin{tabular}{lrrrrrr}
				\toprule
				& & \multicolumn{5}{c}{Bias (RMSE)} \\
				\cmidrule(lr){3-7}
				Variable & True AME & $n=250$ & $n=500$ & $n=1,000$ & $n=2,500$ & $n=5,000$ \\
				\midrule
				\multicolumn{7}{l}{\textit{Panel: Linear}} \\
				\quad x1 & 2.000 & -0.002 & -0.000 & -0.000 & -0.001 & -0.000 \\
				&  & (0.067) & (0.044) & (0.032) & (0.021) & (0.014) \\
				\quad x2 & 3.000 & -0.004 & -0.002 & 0.001 & -0.000 & -0.000 \\
				&  & (0.064) & (0.045) & (0.032) & (0.020) & (0.014) \\
				\quad x3 & 0.000 & -0.000 & -0.001 & -0.000 & 0.001 & -0.000 \\
				&  & (0.065) & (0.046) & (0.031) & (0.020) & (0.014) \\
				\quad x4 & 0.000 & -0.003 & 0.001 & -0.000 & -0.000 & -0.000 \\
				&  & (0.065) & (0.045) & (0.031) & (0.020) & (0.014) \\
				\midrule
				\multicolumn{7}{l}{\textit{Panel: Random Forest}} \\
				\quad x1 & 2.000 & -0.002 & 0.002 & -0.002 & 0.000 & 0.000 \\
				&  & (0.171) & (0.106) & (0.074) & (0.044) & (0.031) \\
				\quad x2 & 3.000 & -0.001 & 0.010 & 0.003 & 0.001 & -0.000 \\
				&  & (0.228) & (0.147) & (0.094) & (0.056) & (0.038) \\
				\quad x3 & 0.000 & 0.002 & 0.001 & -0.001 & 0.001 & -0.000 \\
				&  & (0.087) & (0.054) & (0.035) & (0.022) & (0.015) \\
				\quad x4 & 0.000 & -0.002 & 0.000 & -0.001 & -0.001 & -0.000 \\
				&  & (0.089) & (0.056) & (0.036) & (0.021) & (0.015) \\
				\midrule
				\multicolumn{7}{l}{\textit{Panel: XGBoost}} \\
				\quad x1 & 2.000 & -0.002 & 0.001 & -0.001 & -0.002 & 0.000 \\
				&  & (0.230) & (0.142) & (0.088) & (0.048) & (0.030) \\
				\quad x2 & 3.000 & -0.002 & 0.007 & 0.001 & 0.001 & 0.001 \\
				&  & (0.280) & (0.173) & (0.101) & (0.054) & (0.034) \\
				\quad x3 & 0.000 & -0.001 & 0.000 & -0.000 & 0.001 & -0.000 \\
				&  & (0.095) & (0.058) & (0.036) & (0.022) & (0.015) \\
				\quad x4 & 0.000 & -0.004 & 0.001 & 0.000 & -0.000 & -0.000 \\
				&  & (0.099) & (0.061) & (0.036) & (0.021) & (0.014) \\
				\midrule
				\multicolumn{7}{l}{\textit{Panel: Neural Net}} \\
				\quad x1 & 2.000 & -0.002 & -0.001 & 0.000 & -0.001 & -0.000 \\
				&  & (0.084) & (0.051) & (0.034) & (0.022) & (0.014) \\
				\quad x2 & 3.000 & -0.004 & -0.002 & 0.001 & -0.001 & -0.000 \\
				&  & (0.090) & (0.055) & (0.036) & (0.021) & (0.015) \\
				\quad x3 & 0.000 & 0.000 & -0.001 & -0.000 & 0.001 & -0.000 \\
				&  & (0.079) & (0.049) & (0.033) & (0.021) & (0.014) \\
				\quad x4 & 0.000 & -0.003 & 0.001 & -0.001 & 0.000 & -0.000 \\
				&  & (0.080) & (0.051) & (0.033) & (0.020) & (0.014) \\
				\bottomrule
				\end{tabular}
				\begin{tablenotes}
				\footnotesize
				\item \textit{Notes:} Debiased cross-fitted estimator with $K=5$ folds and the closed-form Riesz representer $\alpha_h(u) = e^{-h^2/2}\sinh(h u_j)/h$, exact for independent standard normal covariates. Trimming is inactive because the support is unbounded. Bias is the mean of $\hat{\theta}_{j,h} - \theta_{j,h}$ and RMSE its root mean square, over the Monte Carlo iterations. True values computed by Monte Carlo integration at $n=1{,}000{,}000$. Noise features $x_3$ and $x_4$ have a true effect of zero.
				\end{tablenotes}
		\end{threeparttable}
	\end{table*}
	
	Table~\ref{tab:sim3_bias_regression_linear} reports bias and RMSE for the linear regression DGP. The result to compare against is the random forest panel of Table~\ref{tab:sim1_regression_linear}, where the plug-in's bias was $-0.097$ at $n=250$ and, rather than shrinking, grew to $-0.091$ for $x_1$ and $-0.095$ for $x_2$ at $n=5{,}000$. Under the correction that bias is $-0.002$ at $n=250$ and $0.000$ at $n=5{,}000$, and the RMSE now falls monotonically from $0.171$ to $0.031$, at the $O(n^{-1/2})$ rate the plug-in never achieved. The correction does not merely shrink the bias at a given sample size; it restores the rate, which is what inference requires. The other learners were already close to unbiased on this design and are essentially unchanged, as they should be: the orthogonal score corrects a first-order term that is near zero for them.
	
	\begin{table*}[htbp]
		\centering
		\small
		\begin{threeparttable}
			\caption{Coverage of nominal 95\% confidence intervals, linear regression DGP: refitting bootstrap against the influence function.}
			\label{tab:sim3_coverage_regression_linear}
			\begin{tabular}{lrrrrrrr}
				\toprule
				& & \multicolumn{3}{c}{Refitting bootstrap} & \multicolumn{3}{c}{Influence function} \\
				\cmidrule(lr){3-5} \cmidrule(lr){6-8}
				Variable & True AME & $n=250$ & $n=1,000$ & $n=5,000$ & $n=250$ & $n=1,000$ & $n=5,000$ \\
				\midrule
				\multicolumn{8}{l}{\textit{Panel: Linear}} \\
				\quad x1 & 2.000 & 0.928 & 0.945 & 0.954 & 0.930 & 0.951 & 0.957 \\
				\quad x2 & 3.000 & 0.955 & 0.948 & 0.956 & 0.955 & 0.946 & 0.960 \\
				\quad x3 & 0.000 & 0.951 & 0.947 & 0.948 & 0.944 & 0.949 & 0.950 \\
				\quad x4 & 0.000 & 0.949 & 0.948 & 0.957 & 0.941 & 0.949 & 0.956 \\
				\midrule
				\multicolumn{8}{l}{\textit{Panel: Random Forest}} \\
				\quad x1 & 2.000 & 0.870 & 0.886 & 0.092 & 0.936 & 0.939 & 0.946 \\
				\quad x2 & 3.000 & 0.967 & 0.962 & 0.323 & 0.936 & 0.952 & 0.952 \\
				\quad x3 & 0.000 & 0.996 & 0.996 & 0.999 & 0.953 & 0.963 & 0.950 \\
				\quad x4 & 0.000 & 0.992 & 0.995 & 1.000 & 0.940 & 0.950 & 0.964 \\
				\midrule
				\multicolumn{8}{l}{\textit{Panel: XGBoost}} \\
				\quad x1 & 2.000 & 0.849 & 0.982 & 0.521 & 0.932 & 0.942 & 0.938 \\
				\quad x2 & 3.000 & 0.917 & 0.990 & 0.503 & 0.948 & 0.946 & 0.941 \\
				\quad x3 & 0.000 & 0.988 & 0.969 & 0.971 & 0.957 & 0.952 & 0.952 \\
				\quad x4 & 0.000 & 0.978 & 0.975 & 0.969 & 0.940 & 0.957 & 0.962 \\
				\midrule
				\multicolumn{8}{l}{\textit{Panel: Neural Net}} \\
				\quad x1 & 2.000 & 0.950 & 0.978 & 1.000 & 0.940 & 0.948 & 0.956 \\
				\quad x2 & 3.000 & 0.936 & 0.974 & 1.000 & 0.950 & 0.939 & 0.959 \\
				\quad x3 & 0.000 & 0.972 & 0.970 & 1.000 & 0.947 & 0.951 & 0.948 \\
				\quad x4 & 0.000 & 0.968 & 0.972 & 1.000 & 0.949 & 0.952 & 0.946 \\
				\bottomrule
				\end{tabular}
				\begin{tablenotes}
				\footnotesize
				\item \textit{Notes:} Empirical coverage of nominal 95\% intervals. The bootstrap columns resample and refit the learner, centring on the plug-in average; the influence-function columns use the debiased cross-fitted estimator with $K=5$ folds and the closed-form representer. Both are evaluated against the same window estimand $\theta_{j,h}$ and the same Monte Carlo ground truth. Noise features $x_3$ and $x_4$ have a true effect of zero.
				\end{tablenotes}
		\end{threeparttable}
	\end{table*}
	
	Table~\ref{tab:sim3_coverage_regression_linear} places the two inference procedures side by side on the same designs and the same ground truth. The bootstrap columns range from $0.092$ to $1.000$; the influence-function columns range from $0.930$ to $0.964$. Across all 300 cells of the closed-form arm --- three data generating processes, two outcome types, five learners, five sample sizes and four features --- empirical coverage has mean $0.948$ and lies between $0.920$ and $0.970$, with $98.7\%$ of cells inside $[0.93, 0.97]$. No learner is an exception: mean coverage by model runs from $0.947$ to $0.949$. The forest cell that motivates the paper moves from $0.092$ to $0.946$.
	
	The forest panel carries a second reading. These forests are capped at depth six, and a depth-capped forest does not converge to the regression function; it converges to the best depth-limited approximation of it. The estimator is therefore consistent for neither nuisance in the usual sense, and yet its intervals cover at the nominal rate. This is Proposition~\ref{prop:robust} operating exactly as stated: with the representer known, $\sqrt{n}$-normality survives an outcome model that converges to the wrong function, and the learner's approximation error is paid in the width of the interval rather than in its location. The forest's RMSE at $n=5{,}000$ is $0.031$ against the linear model's $0.014$ --- wider, and correctly centred. That is the practical content of double robustness here, and it is why the choice of learner becomes a question of efficiency rather than of validity.
	
	\begin{table*}[htbp]
		\centering
		\small
		\begin{threeparttable}
			\caption{Simulation 3: debiased estimator, linear classification DGP. Bias and RMSE (in parentheses) by sample size.}
			\label{tab:sim3_bias_classification_linear}
			\begin{tabular}{lrrrrrr}
				\toprule
				& & \multicolumn{5}{c}{Bias (RMSE)} \\
				\cmidrule(lr){3-7}
				Variable & True AME & $n=250$ & $n=500$ & $n=1,000$ & $n=2,500$ & $n=5,000$ \\
				\midrule
				\multicolumn{7}{l}{\textit{Panel: Logistic}} \\
				\quad x1 & 0.1990 & 0.0019 & 0.0005 & 0.0006 & 0.0000 & -0.0002 \\
				&  & (0.0199) & (0.0144) & (0.0105) & (0.0064) & (0.0048) \\
				\quad x2 & 0.2980 & 0.0004 & -0.0001 & 0.0003 & 0.0002 & 0.0002 \\
				&  & (0.0206) & (0.0142) & (0.0102) & (0.0065) & (0.0046) \\
				\quad x3 & 0.0000 & 0.0013 & 0.0007 & 0.0007 & 0.0003 & 0.0001 \\
				&  & (0.0208) & (0.0139) & (0.0100) & (0.0063) & (0.0045) \\
				\quad x4 & 0.0000 & 0.0008 & -0.0002 & -0.0008 & -0.0001 & 0.0003 \\
				&  & (0.0206) & (0.0139) & (0.0101) & (0.0062) & (0.0045) \\
				\midrule
				\multicolumn{7}{l}{\textit{Panel: Random Forest}} \\
				\quad x1 & 0.1990 & 0.0034 & -0.0009 & 0.0008 & 0.0003 & -0.0003 \\
				&  & (0.0339) & (0.0223) & (0.0152) & (0.0091) & (0.0060) \\
				\quad x2 & 0.2980 & -0.0003 & 0.0024 & 0.0004 & 0.0006 & 0.0004 \\
				&  & (0.0437) & (0.0269) & (0.0176) & (0.0101) & (0.0063) \\
				\quad x3 & 0.0000 & 0.0003 & 0.0009 & 0.0003 & 0.0003 & 0.0003 \\
				&  & (0.0246) & (0.0162) & (0.0111) & (0.0068) & (0.0047) \\
				\quad x4 & 0.0000 & -0.0001 & -0.0010 & -0.0010 & -0.0002 & 0.0002 \\
				&  & (0.0249) & (0.0168) & (0.0113) & (0.0065) & (0.0046) \\
				\midrule
				\multicolumn{7}{l}{\textit{Panel: XGBoost}} \\
				\quad x1 & 0.1990 & 0.0050 & -0.0025 & 0.0013 & 0.0001 & -0.0003 \\
				&  & (0.0495) & (0.0304) & (0.0193) & (0.0107) & (0.0069) \\
				\quad x2 & 0.2980 & -0.0003 & 0.0035 & 0.0005 & 0.0010 & 0.0003 \\
				&  & (0.0630) & (0.0395) & (0.0220) & (0.0119) & (0.0069) \\
				\quad x3 & 0.0000 & 0.0011 & 0.0007 & 0.0004 & 0.0002 & 0.0002 \\
				&  & (0.0260) & (0.0171) & (0.0112) & (0.0068) & (0.0047) \\
				\quad x4 & 0.0000 & 0.0004 & -0.0003 & -0.0004 & -0.0001 & 0.0003 \\
				&  & (0.0268) & (0.0174) & (0.0117) & (0.0065) & (0.0046) \\
				\midrule
				\multicolumn{7}{l}{\textit{Panel: Neural Net}} \\
				\quad x1 & 0.1990 & 0.0019 & 0.0005 & 0.0004 & 0.0001 & -0.0001 \\
				&  & (0.0217) & (0.0158) & (0.0115) & (0.0068) & (0.0049) \\
				\quad x2 & 0.2980 & 0.0004 & -0.0003 & 0.0000 & 0.0004 & 0.0002 \\
				&  & (0.0229) & (0.0171) & (0.0117) & (0.0072) & (0.0048) \\
				\quad x3 & 0.0000 & 0.0012 & 0.0010 & 0.0005 & 0.0004 & 0.0002 \\
				&  & (0.0216) & (0.0144) & (0.0101) & (0.0063) & (0.0045) \\
				\quad x4 & 0.0000 & 0.0008 & -0.0005 & -0.0008 & -0.0001 & 0.0003 \\
				&  & (0.0216) & (0.0141) & (0.0105) & (0.0063) & (0.0045) \\
				\bottomrule
				\end{tabular}
				\begin{tablenotes}
				\footnotesize
				\item \textit{Notes:} Debiased cross-fitted estimator with $K=5$ folds and the closed-form Riesz representer $\alpha_h(u) = e^{-h^2/2}\sinh(h u_j)/h$, exact for independent standard normal covariates. Trimming is inactive because the support is unbounded. Bias is the mean of $\hat{\theta}_{j,h} - \theta_{j,h}$ and RMSE its root mean square, over the Monte Carlo iterations. True values computed by Monte Carlo integration at $n=1{,}000{,}000$. Noise features $x_3$ and $x_4$ have a true effect of zero.
				\end{tablenotes}
		\end{threeparttable}
	\end{table*}
	
	\begin{table*}[htbp]
		\centering
		\small
		\begin{threeparttable}
			\caption{Coverage of nominal 95\% confidence intervals, linear classification DGP: refitting bootstrap against the influence function.}
			\label{tab:sim3_coverage_classification_linear}
			\begin{tabular}{lrrrrrrr}
				\toprule
				& & \multicolumn{3}{c}{Refitting bootstrap} & \multicolumn{3}{c}{Influence function} \\
				\cmidrule(lr){3-5} \cmidrule(lr){6-8}
				Variable & True AME & $n=250$ & $n=1,000$ & $n=5,000$ & $n=250$ & $n=1,000$ & $n=5,000$ \\
				\midrule
				\multicolumn{8}{l}{\textit{Panel: Logistic}} \\
				\quad x1 & 0.199 & 0.881 & 0.878 & 0.886 & 0.944 & 0.934 & 0.940 \\
				\quad x2 & 0.298 & 0.726 & 0.766 & 0.781 & 0.924 & 0.956 & 0.954 \\
				\quad x3 & 0.000 & 0.944 & 0.939 & 0.946 & 0.934 & 0.950 & 0.952 \\
				\quad x4 & 0.000 & 0.933 & 0.954 & 0.944 & 0.948 & 0.950 & 0.944 \\
				\midrule
				\multicolumn{8}{l}{\textit{Panel: Random Forest}} \\
				\quad x1 & 0.199 & 0.968 & 0.933 & 0.186 & 0.942 & 0.956 & 0.944 \\
				\quad x2 & 0.298 & 0.961 & 0.958 & 0.268 & 0.932 & 0.952 & 0.946 \\
				\quad x3 & 0.000 & 0.962 & 0.942 & 0.949 & 0.944 & 0.960 & 0.948 \\
				\quad x4 & 0.000 & 0.969 & 0.964 & 0.953 & 0.944 & 0.944 & 0.944 \\
				\midrule
				\multicolumn{8}{l}{\textit{Panel: XGBoost}} \\
				\quad x1 & 0.199 & 0.926 & 0.947 & 0.877 & 0.932 & 0.952 & 0.950 \\
				\quad x2 & 0.298 & 0.908 & 0.954 & 0.910 & 0.934 & 0.968 & 0.972 \\
				\quad x3 & 0.000 & 0.976 & 0.961 & 0.960 & 0.964 & 0.948 & 0.946 \\
				\quad x4 & 0.000 & 0.973 & 0.969 & 0.959 & 0.954 & 0.938 & 0.946 \\
				\midrule
				\multicolumn{8}{l}{\textit{Panel: Neural Net}} \\
				\quad x1 & 0.199 & 0.870 & 0.870 & 0.958 & 0.940 & 0.930 & 0.942 \\
				\quad x2 & 0.298 & 0.658 & 0.774 & 0.894 & 0.934 & 0.950 & 0.958 \\
				\quad x3 & 0.000 & 0.936 & 0.936 & 0.986 & 0.932 & 0.946 & 0.954 \\
				\quad x4 & 0.000 & 0.938 & 0.936 & 0.982 & 0.946 & 0.932 & 0.944 \\
				\bottomrule
				\end{tabular}
				\begin{tablenotes}
				\footnotesize
				\item \textit{Notes:} Empirical coverage of nominal 95\% intervals. The bootstrap columns resample and refit the learner, centring on the plug-in average; the influence-function columns use the debiased cross-fitted estimator with $K=5$ folds and the closed-form representer. Both are evaluated against the same window estimand $\theta_{j,h}$ and the same Monte Carlo ground truth. Noise features $x_3$ and $x_4$ have a true effect of zero.
				\end{tablenotes}
		\end{threeparttable}
	\end{table*}
	
	Tables~\ref{tab:sim3_bias_classification_linear} and~\ref{tab:sim3_coverage_classification_linear} report the corresponding classification results, where the same pattern holds on a compressed scale: the forest's coverage at $n=5{,}000$ moves from $0.186$ to $0.944$ for $x_1$ and from $0.268$ to $0.946$ for $x_2$, and the debiased intervals are centred for every learner at every sample size. The nonlinear and interaction designs behave identically and are reported in the supplement.
	
	The closed-form representer used above is a luxury of the simulation design, and the natural question is what the estimator costs when $\alpha_h$ must itself be estimated. The second arm answers it by replacing the closed form with a sieve estimate, holding everything else fixed. Coverage survives, with a mild conservatism: over the same cells the sieve arm has mean coverage $0.958$ against the closed form's $0.948$. The cost appears where theory says it should, in the product bias of Assumption~\ref{ass:nuisance}(b) and in efficiency, and it vanishes with sample size. Mean absolute bias is $3.3$ times the closed form's at $n=250$, then $1.6$, $1.2$, $1.1$, and $0.9$ at $n=5{,}000$ --- indistinguishable by the largest sample size. Standard errors are $15\%$ wider at $n=250$ and $1\%$ wider at $n=5{,}000$. Estimating the representer costs something real at small samples and asymptotically nothing, which is the behaviour the product structure of the bias predicts: neither nuisance need converge quickly on its own, so long as the product of their errors is $o(n^{-1/2})$.
	
	\section{Conclusion}
	\label{sec:conclusion}
	
	This paper has proposed a framework for interpreting any supervised machine learning
	model through average marginal effects. The central claim is that the traditional
	accuracy-interpretability tradeoff is not fundamental. The move that carries the
	argument is to make the window the estimand rather than a numerical approximation to
	something else: the window average marginal effect is the average centered difference of
	the regression function at a fixed, analyst-chosen step, and it is well defined for a
	learner that is differentiable nowhere. A single estimand therefore covers trees,
	ensembles and networks without differentiability assumptions of any kind, and the
	classical average derivative is recovered in the small-window limit whenever one weak
	derivative exists. The estimator applies uniformly across model classes and outcome
	types and produces a coefficient table --- estimate, standard error, and significance
	level --- directly analogous to the output of a parametric regression. Standard errors
	come from the estimated influence function of a debiased, cross-fitted estimator in a
	single pass, with a model-refitting bootstrap retained as a diagnostic for specification
	uncertainty. These are capabilities that existing interpretability methods, including
	SHAP \citep{lundberg2017shap}, LIME \citep{ribeiro2016lime}, and partial dependence
	plots \citep{friedman2001pdp}, do not provide.
	
	The theoretical contribution proceeds in three parts. The first is the estimand
	and its dual: the window average marginal effect is linear in the regression
	function and admits an exact Riesz representation by a change of variables, with a
	bounded representer, no boundary terms, and no density derivatives --- and the
	classical average derivative is recovered in the small-window limit under one weak
	derivative. The second is inference and optimality: the orthogonal score built
	from the representer has estimation bias exactly equal to a product of two $L^2$
	nuisance errors, so the cross-fitted estimator is root-$n$ consistent and
	asymptotically normal under $n^{-1/4}$-type $L^2$ rates that flexible learners can
	satisfy, and it attains the semiparametric efficiency bound --- replacing the
	parametric-rate function-space condition that a functional delta method treatment
	would require and that growing-complexity learners cannot meet; a double
	robustness proposition extends validity to inconsistent, fixed-complexity
	learners whenever the representer is known or parametrically estimable. The third is
	interpretation: a pairing identity shows that the window effect of any estimator
	of bounded variation equals its jump measure integrated against a locally averaged
	covariate density, characterizing the population quantity that a tree-based
	model's marginal effect measures; and gradient consistency is established for
	neural networks in the Sobolev space $W^{1,p}$ via Rellich--Kondrachov compactness
	\citep{evans2010pde}, and for gradient boosted models with tree base learners and
	Mondrian forests \citep{mourtada2017mondrian} in the space of functions of bounded
	variation \citep{ambrosio2000bv} via BV compactness, characterizing the regime in
	which the uncorrected plug-in is already consistent. For Breiman random forests
	\citep{breiman2001rf, scornet2015forests}, gradient consistency is conjectured
	rather than proved --- the obstruction lies in the data-dependent partition
	boundaries --- but the inference theory does not require it, so what remains open
	for that class is an $L^2$ rate question rather than a function-space one.
	
	The empirical results across three datasets --- UCI Adult Income \citep{kohavi1996,
		UCI}, UCI Credit Card Default \citep{yeh2009, UCI}, and Ames Housing
	\citep{decock2011} --- validate the framework on binary classification and continuous
	regression outcomes. Under the debiased estimator the four model classes agree far
	more closely than they did under the plug-in, and they do so because they are now
	estimating the same quantity rather than each reporting the slope of its own fit.
	The credit default application shows this most sharply: the plug-in put the effect
	of the payment ratio at $-$0.431 and $-$0.440 for the two tree ensembles, $-$0.045
	for logistic regression and 0.0009 --- effectively nothing --- for the neural
	network, a set of estimates that cannot all be describing the same world; the
	corrected estimates are $-$0.489, $-$0.584, $-$0.265 and $-$0.194, all negative,
	all significant, and all telling one story. The variation that the plug-in
	displayed was in the main the regularization bias, which differs across learners
	because their regularization differs.

	The standard errors move in the direction the theory predicts as well, and in
	several places that means upward. The refitting bootstrap reported standard errors
	of a few dollars for the neural network on Ames Housing and of 0.0000 to four
	decimal places for the random forest on credit default; the influence-function
	standard errors for the same quantities are three to four orders of magnitude
	larger. Those earlier intervals were narrow around a displaced center, which is the
	failure mode isolated in Simulation~2. The cost of the correction is visible in the
	smallest application: with $n = 1{,}460$ observations, several Ames Housing effects
	that the plug-in reported as precisely determined are no longer distinguishable
	from zero, and the estimator reports that rather than concealing it.

	The comparison with SHAP values reveals a systematic divergence: SHAP values are on
	a different scale, in several cases have the wrong sign relative to the AME, and
	carry no inferential apparatus. The Ames Housing application is particularly
	informative because the dollar-denominated units make the divergence concrete ---
	SHAP attributes to above-ground living area an effect fifty times larger than the
	AME for the linear model, a number that is not interpretable as a marginal effect
	in any standard sense. Partial dependence slopes require a subtler reading than the
	earlier draft gave them. A partial dependence slope is itself a plug-in average
	over a single full-sample fit, so it agrees with the debiased AME where the
	plug-in is nearly unbiased --- on smooth, well-supported, continuous features ---
	and falls short of it, in about three quarters of the resolved cases toward zero,
	wherever regularization flattens the fitted surface. What once read as two
	estimators corroborating each other is better read as a direct measurement of the
	bias between them. An
	additional finding is that Breiman random forests produce less stable plug-in AME
	estimates than the other estimator classes, and the simulation study traces the
	resulting coverage failure to the same regularization bias --- the first-order term
	that the orthogonal correction removes, and the clearest empirical argument for
	reporting the debiased estimator rather than the uncorrected average.
	
	Several important problems remain open. The most pressing concerns Breiman
	forests, and the present framework splits it into two questions of very different
	difficulty. The inferential question --- valid confidence intervals for the window
	effect when the learner is a Breiman forest --- now requires only an $L^2$
	convergence rate for the forest, since Theorem~\ref{thm:main} consumes nothing
	else; consistency is available under structural assumptions
	\citep{scornet2015forests}, and converting it into a rate meeting the
	$n^{-1/4}$ benchmark, or exploiting the product form of
	Assumption~\ref{ass:nuisance}(b) with a fast representer estimate, is a concrete
	and plausibly tractable target. Simulation~3 indicates how much of that question
	is already settled in the case the theory handles cleanly. With the representer
	known, Proposition~\ref{prop:robust} asks nothing of the forest beyond
	convergence to some fixed limit, and the simulation exhibits exactly that
	situation: a depth-constrained forest whose plug-in error grows with the sample
	size, and therefore is not approaching $f$, nevertheless yields intervals whose
	coverage averages $0.948$ across the whole forest design and never leaves the
	band $0.922$ to $0.968$. What
	the open $L^2$ question governs is the case where the representer must itself be
	estimated, and it is there rather than in the known-representer case that a rate
	for Breiman forests would buy something. The interpretive question --- convergence of the
	forest's own gradient information, in the sense of
	Section~\ref{sec:learners} --- remains genuinely hard: the obstruction identified
	in the proof of Theorem~\ref{thm:forests}, that data-dependent splits couple
	partition geometry to response variation in a way that prevents the key
	decomposition, suggests that new analytical tools will be required, possibly a
	coupling argument relating Breiman partitions to Mondrian partitions. Separating
	the two questions is itself progress: the most widely used forest implementation
	can be brought under the inferential guarantees without waiting for the harder
	interpretation theory.
	
	A second open problem is the extension to additional model classes, and the
	framework now presents two templates with distinct proof obligations. For
	inference, the obligation is minimal and uniform across classes: an $L^2(P)$
	convergence rate for the learner --- support vector machines with smooth kernels,
	Bayesian additive regression trees, generalized additive models, and kernel ridge
	regression all have such rates in the literature or within reach --- together
	with a Riesz-regression class rich enough to approximate the bounded representer,
	for which \citet{chernozhukov2022riesz} provide general guarantees. For
	interpretation, the Sobolev and BV architecture developed here remains the
	template: identify the natural function space of the estimator class, establish
	derivative existence and consistency there, and conclude gradient convergence
	via the appropriate compactness theorem. Splines and ridge regression are
	particularly natural candidates for the second template since they live in
	Sobolev spaces by construction, and the pairing identity of
	Theorem~\ref{thm:pairing} extends the first template to any learner of bounded
	variation without further work.
	
	A third open problem is the extension to causal inference. The average marginal effect
	as defined here is a predictive quantity --- the average derivative of the conditional
	expectation function with respect to $x_j$, under the observational distribution of
	the covariates. It is not in general a causal effect. Extending the framework to
	settings where causal identification is possible, for example through instrumental
	variables or regression discontinuity designs, would substantially broaden its
	applicability to policy-relevant questions. The honest forest framework of
	\citet{wager2018forests} and the broader literature on causal machine learning suggest
	a natural path forward, but the connection to the AME estimand studied here has not
	been formalized.
	
	A fourth direction concerns the discrete and mixed-type feature setting. The current
	framework handles binary variables through a finite difference contrast and continuous
	variables through an adaptive centered finite difference, but ordered categorical
	variables, count variables, and interactions among features of mixed type are not
	treated systematically. The marginal effects at the mean versus average marginal
	effects distinction \citep{greene2012, wooldridge2010} also warrants further study in
	the nonparametric setting, where the mean of the covariates may lie in a sparse region
	of the feature space and the marginal effect at the mean may therefore be a poor
	summary of the average effect.
	
	Finally, several practical questions about the debiased procedure itself deserve
	systematic study. The influence-function standard errors of
	Section~2.4 already remove the computational objection to inference --- $K$ model
	fits in place of $B$ refits with hyperparameter selection --- but the finite-sample
	behavior of the procedure depends on choices the asymptotic theory does not
	pin down: the number of folds, the truncation level for the estimated
	representer, the trimming weight near the boundary, and above all the window
	$h$, which is part of the estimand and whose sensitivity should be reported the
	way bandwidth sensitivity is. The multiplier bootstrap of
	Remark~\ref{rem:bootstrap} makes simultaneous confidence bands across features
	essentially free, and characterizing its finite-sample accuracy relative to the
	plug-in intervals, along with honest procedures when the density bounds of
	Assumption~\ref{ass:sampling} are weak near the boundary, is a practical program
	we leave to future work.
	
	\appendix
	
	% ============================================================
	\section{Proof of Lemma~\ref{lem:riesz} (Riesz Representation)}
	\label{app:lem_riesz}
	% ============================================================
	
	\begin{proof}
		Recall $q = wp$ extended by zero outside $\Omega$, so that $q$ is defined on
		all of $\R^d$, $\norm{q}_\infty \leq p_{\max}$, and $q(x) \neq 0$ implies
		$x \in \operatorname{supp}(w) \subseteq \Omega_{j,h}$. Fix $g \in L^1(P)$ and
		consider the positive-shift half of the finite difference:
		\begin{equation}
			\E\left[w(X)\, g(X + h e_j)\right]
			= \int_{\R^d} q(x)\, g(x + h e_j)\, dx
			= \int_{\R^d} q(u - h e_j)\, g(u)\, du
		\end{equation}
		by the substitution $u = x + h e_j$. The manipulation is legitimate because
		on the set where the integrand is nonzero we have $x \in \Omega_{j,h}$, hence
		$u = x + h e_j \in \Omega$, so $g$ is evaluated only on $\Omega$; and the
		integral is absolutely convergent since:
		\begin{equation}
			\int q(u - h e_j)\, \abs{g(u)}\, du
			\leq p_{\max} \int_\Omega \abs{g(u)}\, du
			\leq \frac{p_{\max}}{p_{\min}} \, \E\abs{g(X)} < \infty
		\end{equation}
		The negative-shift half is identical with $u = x - h e_j$. Subtracting and
		dividing by $2h$:
		\begin{equation}
			\E\left[w(X)\, \Dh g(X)\right]
			= \int_\Omega \frac{q(u - h e_j) - q(u + h e_j)}{2h}\, g(u)\, du
			= \E\left[\alpha_h(X)\, g(X)\right]
		\end{equation}
		where the restriction of the domain of integration to $\Omega$ in the second
		step uses that $g$ is supported there, and the final step multiplies and
		divides by $p(u) \geq p_{\min} > 0$. The uniform bound follows from
		$\abs{q(u - h e_j) - q(u + h e_j)} \leq 2 p_{\max}$ and $p(u) \geq p_{\min}$.
		Applying the representation to $g \equiv 1$ gives
		$\E[\alpha_h(X)] = \E[w(X) \Dh 1] = 0$ since the finite difference of a
		constant vanishes.
	\end{proof}
	
	% ============================================================
	\section{Proof of Proposition~\ref{prop:hlimit} (Recovery of the Average
		Derivative)}
	\label{app:prop_hlimit}
	% ============================================================
	
	\begin{proof}
		Let $f \in W^{1,1}(\Omega)$ and fix $h \leq h_0$, so that
		$\operatorname{supp}(w) \subseteq \Omega_{j,h_0} \subseteq \Omega_{j,h}$ and
		every segment $\{x + t e_j : \abs{t} \leq h\}$ with
		$x \in \operatorname{supp}(w)$ lies in $\Omega$. Sobolev functions are
		absolutely continuous on almost every line parallel to a coordinate axis, with
		the classical derivative along the line agreeing almost everywhere with the
		weak derivative \citep[Section~4.9]{evans2015measure}. Hence for almost every
		$x \in \operatorname{supp}(w)$:
		\begin{equation}
			f(x + h e_j) - f(x - h e_j)
			= \int_{-h}^{h} \partial_j f(x + t e_j)\, dt
		\end{equation}
		and therefore:
		\begin{equation}
			\theta_h(f) - \E\left[w(X)\, \partial_j f(X)\right]
			= \int_{\R^d} q(x)\, \frac{1}{2h} \int_{-h}^{h}
			\left[\partial_j f(x + t e_j) - \partial_j f(x)\right] dt\, dx
		\end{equation}
		Bounding $q$ by $p_{\max}$, applying Tonelli, and extending
		$\partial_j f$ by zero outside $\Omega$:
		\begin{equation}
			\abs{\theta_h(f) - \E\left[w\, \partial_j f\right]}
			\leq p_{\max} \sup_{\abs{t} \leq h}
			\int_{\R^d} \abs{\partial_j f(x + t e_j) - \partial_j f(x)}
			\, \mathbf{1}\{x \in \operatorname{supp}(w)\}\, dx
		\end{equation}
		where the extension introduces no error because $x + t e_j \in \Omega$ for
		$x \in \operatorname{supp}(w)$ and $\abs{t} \leq h \leq h_0$. Continuity of
		translation in $L^1(\R^d)$ applied to the extension of $\partial_j f$ sends
		the right side to zero as $h \downarrow 0$. For the second claim,
		$\abs{w_m \partial_j f} \leq \abs{\partial_j f}$, which is $P$-integrable
		because $\partial_j f \in L^1(\Omega)$ and $p \leq p_{\max}$; dominated
		convergence gives
		$\E[w_m \partial_j f] \to \E[\partial_j f] = \mathrm{AME}_j$.
	\end{proof}
	
	% ============================================================
	\section{Proof of Lemma~\ref{lem:orth} (Exact Product Bias)}
	\label{app:lem_orth}
	% ============================================================
	
	\begin{proof}
		Conditioning on $X$ and using $\E[Y \mid X] = f(X)$:
		\begin{equation}
			\E\left[\tilde{\alpha}(X)\,(Y - \tilde{f}(X))\right]
			= \E\left[\tilde{\alpha}(X)\,(f(X) - \tilde{f}(X))\right]
		\end{equation}
		By linearity of $\Dh$ and Lemma~\ref{lem:riesz} applied to
		$g = \tilde{f} - f \in L^1(P)$:
		\begin{equation}
			\E\left[w(X)\, \Dh \tilde{f}(X)\right] - \theta_h(f)
			= \E\left[w(X)\, \Dh (\tilde{f} - f)(X)\right]
			= \E\left[\alpha_h(X)\,(\tilde{f}(X) - f(X))\right]
		\end{equation}
		Summing the two displays gives:
		\begin{equation}
			\E\left[\psi(X, Y; \tilde{f}, \tilde{\alpha}, \theta_h(f))\right]
			= \E\left[(\alpha_h(X) - \tilde{\alpha}(X))(\tilde{f}(X) - f(X))\right]
		\end{equation}
		and the bound is Cauchy--Schwarz.
	\end{proof}
	
	% ============================================================
	\section{Proof of Theorem~\ref{thm:main} (Asymptotic Linearity and Normality)}
	\label{app:thm_main}
	% ============================================================
	
	\subsection{Step 0: The Trimmed Difference Operator is Bounded on $L^2(P)$}
	
	For any $g \in L^2(\Omega)$, using $(a - b)^2 \leq 2a^2 + 2b^2$ and the
	change of variables from the proof of Lemma~\ref{lem:riesz}:
	\begin{align}
		\E\left[\bigl(w(X)\, \Dh g(X)\bigr)^2\right]
		&\leq \frac{1}{2h^2}\left(
		\E\left[w(X)^2 g(X + h e_j)^2\right]
		+ \E\left[w(X)^2 g(X - h e_j)^2\right]\right) \notag \\
		&\leq \frac{1}{h^2} \cdot \frac{p_{\max}}{p_{\min}}\, \E\left[g(X)^2\right]
		\label{eq:opbound_proof}
	\end{align}
	which establishes \eqref{eq:opbound} with constant
	$\kappa = h^{-1}(p_{\max}/p_{\min})^{1/2}$. This is the only property of the
	finite difference used below; no derivative of $g$ appears anywhere.
	
	\subsection{Step 1: Decomposition by Folds}
	
	Write $\theta = \theta_h(f)$,
	$\psi_i = \psi(\mathbf{x}_i, y_i; f, \alpha_h, \theta)$, and for
	$i \in I_k$ let
	$\hat{\psi}_i = \psi(\mathbf{x}_i, y_i; \hat{f}^{(-k)},
	\hat{\alpha}^{(-k)}, \theta)$. Then:
	\begin{equation}
		\sqrt{n}\,(\hat{\theta}_h - \theta)
		= \frac{1}{\sqrt{n}} \sum_{i=1}^{n} \psi_i
		+ \sum_{k=1}^{K} \sqrt{\frac{n_k}{n}}\,\bigl(\sqrt{n_k}\, A_k
		+ \sqrt{n_k}\, B_k\bigr)
	\end{equation}
	where, with $\E_k$ denoting expectation over a fresh draw $(X, Y)$ holding the
	fold-$k$ nuisances fixed:
	\begin{equation}
		A_k = \E_k\left[\hat{\psi} - \psi\right],
		\qquad
		B_k = \frac{1}{n_k} \sum_{i \in I_k}
		\left(\hat{\psi}_i - \psi_i - A_k\right)
	\end{equation}
	Since $K$ is fixed and $n_k / n \to 1/K$, it suffices that
	$\sqrt{n_k} A_k = o_p(1)$ and $\sqrt{n_k} B_k = o_p(1)$ for each $k$.
	
	\subsection{Step 2: The Bias Term}
	
	By Lemma~\ref{lem:orth} with
	$(\tilde{f}, \tilde{\alpha}) = (\hat{f}^{(-k)}, \hat{\alpha}^{(-k)})$:
	\begin{equation}
		\abs{\sqrt{n_k}\, A_k}
		\leq \sqrt{n_k}\, \Ptwo{\hat{\alpha}^{(-k)} - \alpha_h}
		\cdot \Ptwo{\hat{f}^{(-k)} - f} = o_p(1)
	\end{equation}
	by Assumption~\ref{ass:nuisance}(b). This is the only place the product rate is
	used, and it is used without any expansion: the bound is exact.
	
	\subsection{Step 3: The Empirical Process Term}
	
	Conditional on the out-of-fold data, the summands of $B_k$ are i.i.d.\ and
	mean zero, so:
	\begin{equation}
		\E\left[(\sqrt{n_k}\, B_k)^2 \mid \text{out-of-fold}\right]
		\leq \E_k\left[(\hat{\psi} - \psi)^2\right]
	\end{equation}
	Decompose, adding and subtracting $\hat{\alpha}^{(-k)}(x)(y - f(x))$:
	\begin{equation}
		\hat{\psi} - \psi
		= w\, \Dh(\hat{f}^{(-k)} - f)
		+ (\hat{\alpha}^{(-k)} - \alpha_h)(y - f)
		- \hat{\alpha}^{(-k)}\,(\hat{f}^{(-k)} - f)
	\end{equation}
	and bound each piece in $L^2(P)$: the first by
	$\kappa \Ptwo{\hat{f}^{(-k)} - f}$ via \eqref{eq:opbound_proof}; the second by
	$\bar{\sigma}\, \Ptwo{\hat{\alpha}^{(-k)} - \alpha_h}$ after conditioning on
	$X$; the third by $\bar{A}\, \Ptwo{\hat{f}^{(-k)} - f}$ using
	Assumption~\ref{ass:nuisance}(c). All three are $o_p(1)$ by
	Assumption~\ref{ass:nuisance}(a), so
	$\E_k[(\hat{\psi} - \psi)^2] = o_p(1)$ and conditional Chebyshev gives
	$\sqrt{n_k}\, B_k = o_p(1)$.
	
	\subsection{Step 4: Central Limit Theorem and Variance Estimation}
	
	The $\psi_i$ are i.i.d., mean zero by Lemma~\ref{lem:orth} at
	$(\tilde{f}, \tilde{\alpha}) = (f, \alpha_h)$, with variance:
	\begin{equation}
		V_h = \E[\psi_i^2]
		\leq 2\,\E\left[(w \Dh f)^2\right] + 2\,\norm{\alpha_h}_\infty^2\,
		\E[\sigma^2(X)] < \infty
	\end{equation}
	by \eqref{eq:opbound_proof} and Lemma~\ref{lem:riesz}, so the classical CLT
	applies. For the variance estimator, write
	$\hat{\psi}_i(\hat{\theta}_h) = \hat{\psi}_i - (\hat{\theta}_h - \theta)$; then:
	\begin{equation}
		\hat{V}_h = \frac{1}{n} \sum_i \hat{\psi}_i^2
		- 2(\hat{\theta}_h - \theta)\,\frac{1}{n}\sum_i \hat{\psi}_i
		+ (\hat{\theta}_h - \theta)^2
	\end{equation}
	The last two terms vanish in probability since
	$\hat{\theta}_h - \theta = o_p(1)$ and
	$n^{-1}\sum_i \hat{\psi}_i = O_p(n^{-1/2})$ by Steps~1--3. For the first term:
	\begin{equation}
		\abs{\frac{1}{n}\sum_i \hat{\psi}_i^2 - \frac{1}{n}\sum_i \psi_i^2}
		\leq \frac{1}{n}\sum_i (\hat{\psi}_i - \psi_i)^2
		+ 2\left(\frac{1}{n}\sum_i (\hat{\psi}_i - \psi_i)^2\right)^{1/2}
		\left(\frac{1}{n}\sum_i \psi_i^2\right)^{1/2}
	\end{equation}
	by Cauchy--Schwarz; $n^{-1}\sum_i (\hat{\psi}_i - \psi_i)^2 = o_p(1)$ by the
	fold-wise conditional Markov inequality and the Step~3 bound, and
	$n^{-1}\sum_i \psi_i^2 \xrightarrowp V_h$ by the law of large numbers. Hence
	$\hat{V}_h \xrightarrowp V_h$.
	\qed
	
	% ============================================================
	\section{Proof of Proposition~\ref{prop:robust} (Double Robustness)}
	\label{app:prop_robust}
	% ============================================================
	
	The argument retraces the proof of Theorem~\ref{thm:main} with the pseudo-true
	influence function in place of the efficient one; we record the steps that change.
	
	\subsection{Step 1: The Pseudo-True Influence Function}
	
	Define $\psi_{g,i} = w(\mathbf{x}_i)\, \Dh g(\mathbf{x}_i)
	+ \alpha_h(\mathbf{x}_i)\,(y_i - g(\mathbf{x}_i)) - \theta_h(f)$. By
	Lemma~\ref{lem:orth} with $(\tilde{f}, \tilde{\alpha}) = (g, \alpha_h)$:
	\begin{equation}
		\E[\psi_{g,i}]
		= \E\bigl[(\alpha_h - \alpha_h)(g - f)\bigr] = 0
	\end{equation}
	for every $g \in L^2(\Omega)$, which also establishes
	$\E[m_g(X)] = \theta_h(f)$ after conditioning on $X$. Conditioning likewise gives
	$\E[\psi_{g,i} \mid X] = m_g(X) - \theta_h(f)$ and
	$\mathrm{Var}(\psi_{g,i} \mid X) = \alpha_h(X)^2 \sigma^2(X)$, so
	$\mathrm{Var}(\psi_{g,i}) = V_h(g)$ by the law of total variance; finiteness
	follows from \eqref{eq:opbound_proof}, $g \in L^2(\Omega)$, the bound on
	$\alpha_h$ from Lemma~\ref{lem:riesz}, and $\E[Y^2] < \infty$.
	
	\subsection{Step 2: The Bias Term}
	
	With $A_k = \E_k[\hat{\psi} - \psi_g] = \E_k[\hat{\psi}]$, Lemma~\ref{lem:orth}
	gives $A_k = \E[(\alpha_h - \hat{\alpha}^{(-k)})(\hat{f}^{(-k)} - f)]$, so:
	\begin{equation}
		\abs{\sqrt{n_k}\, A_k}
		\leq \sqrt{n_k}\, \Ptwo{\hat{\alpha}^{(-k)} - \alpha_h}
		\left( \Ptwo{\hat{f}^{(-k)} - g} + \Ptwo{g - f} \right)
		= \sqrt{n_k}\, o_p(n^{-1/2}) \cdot O_p(1) = o_p(1)
	\end{equation}
	by conditions (a) and (b); the second factor is $O_p(1)$ rather than $o_p(1)$
	because $g$ need not equal $f$, and it is exactly here that condition (b) must be
	strengthened to $o_p(n^{-1/2})$ relative to Assumption~\ref{ass:nuisance}(b).
	When $\hat{\alpha}^{(-k)} = \alpha_h$, $A_k = 0$ identically at every $n$;
	averaging the fold-wise identity
	$\E[\hat{\psi}_i + \theta_h(f) \mid \text{out-of-fold data}] = \theta_h(f)$,
	which holds because the observations in fold $k$ are independent of
	$\hat{f}^{(-k)}$, gives the exact unbiasedness claim.
	
	\subsection{Step 3: The Empirical Process Term}
	
	Decompose, for $i \in I_k$:
	\begin{equation}
		\hat{\psi}_i - \psi_{g,i}
		= w\, \Dh(\hat{f}^{(-k)} - g)
		+ (\hat{\alpha}^{(-k)} - \alpha_h)(y_i - \hat{f}^{(-k)})
		+ \alpha_h\,(g - \hat{f}^{(-k)})
	\end{equation}
	The first and third terms have conditional $L^2(P)$ norms bounded by
	$\kappa\, \Ptwo{\hat{f}^{(-k)} - g}$ and
	$\norm{\alpha_h}_\infty \Ptwo{\hat{f}^{(-k)} - g}$ respectively, both $o_p(1)$
	by condition (a). For the middle term, write
	$y - \hat{f}^{(-k)} = (y - f) + (f - g) + (g - \hat{f}^{(-k)})$ and bound the
	three resulting pieces: the first by
	$\bar{\sigma}^2 \Ptwo{\hat{\alpha}^{(-k)} - \alpha_h}^2 = o_p(1)$ after
	conditioning on $X$; the third by
	$(\bar{A} + \norm{\alpha_h}_\infty)^2 \Ptwo{\hat{f}^{(-k)} - g}^2 = o_p(1)$; and
	the second piece, $\E[(\hat{\alpha}^{(-k)} - \alpha_h)^2 (f - g)^2]$, tends to
	zero by dominated convergence along subsequences: the factor
	$(\hat{\alpha}^{(-k)} - \alpha_h)^2$ is bounded by
	$(\bar{A} + \norm{\alpha_h}_\infty)^2$ and converges to zero in $P$-measure by
	condition (b), while $(f - g)^2 \in L^1(P)$ dominates. Conditional Chebyshev then
	gives $\sqrt{n_k}\, B_k = o_p(1)$ exactly as in Step~3 of the proof of
	Theorem~\ref{thm:main}.
	
	\subsection{Step 4: Conclusion}
	
	The $\psi_{g,i}$ are i.i.d., mean zero, with finite positive variance $V_h(g)$,
	so the classical CLT applies to their normalized sum, and Steps~2--3 identify
	$\sqrt{n}(\hat{\theta}_h - \theta_h(f))$ with that sum up to $o_p(1)$. The
	consistency of $\hat{V}_h$ for $V_h(g)$ repeats Step~4 of the proof of
	Theorem~\ref{thm:main} verbatim with $\psi_{g}$ in place of $\psi$, using the
	Step~3 bound just established. Finally, when $g \neq f$ the influence function
	$\psi_g$ differs from the efficient influence function $\psi$ of
	Theorem~\ref{thm:efficiency}; since in the nonparametric model every regular
	asymptotically linear estimator has influence function equal to the efficient
	one, the estimator is not regular in this case, and no conflict with the
	convolution theorem arises.
	\qed
	
	% ============================================================
	\section{Proof of Theorem~\ref{thm:efficiency} (Semiparametric Efficiency)}
	\label{app:thm_eff}
	% ============================================================
	
	\begin{proof}
		Consider one-dimensional submodels
		$dP_t = (1 + t s)\, dP$ indexed by bounded functions
		$s \in L^2_0(P) = \{s : \E[s(X, Y)] = 0\}$, which are well defined for
		$\abs{t} < 1/\norm{s}_\infty$ and differentiable in quadratic mean at $t = 0$
		with score $s$; bounded mean-zero functions are dense in $L^2_0(P)$, so the
		tangent space of the nonparametric model is all of $L^2_0(P)$
		\citep[Chapter~25]{vaart1998}. Write
		$\bar{s}_X(x) = \E[s(X, Y) \mid X = x]$ for the induced score of the
		covariate marginal and
		$f_t(x) = \E_t[Y \mid X = x]$ for the regression function along the path.
		Differentiating the Bayes formula for $f_t$ at $t = 0$ gives, for
		$P$-almost every $x$:
		\begin{equation}
			\dot{f}(x) := \frac{\partial}{\partial t} f_t(x)\Big|_{t=0}
			= \E\left[(Y - f(X))\, s(X, Y) \mid X = x\right]
		\end{equation}
		where differentiation under the integral is justified by
		$\E[Y^2] < \infty$ and the boundedness of $s$. The functional along the path
		is $\theta(t) = \E_t[w(X)\, \Dh f_t(X)]$, and the product rule gives:
		\begin{equation}
			\frac{d\theta}{dt}\Big|_{t=0}
			= \E\left[w\, \Dh f(X)\cdot \bar{s}_X(X)\right]
			+ \E\left[w\, \Dh \dot{f}(X)\right]
		\end{equation}
		For the first term, note $\E[\bar{s}_X] = 0$, so the term equals
		$\E[(w \Dh f(X) - \theta_h(f))\, \bar{s}_X(X)]
		= \E[(w \Dh f(X) - \theta_h(f))\, s(X, Y)]$, the last equality because
		$w \Dh f(X) - \theta_h(f)$ is $X$-measurable. For the second term,
		$\dot{f} \in L^1(P)$ (indeed $L^2$) since $s$ is bounded, so
		Lemma~\ref{lem:riesz} applies:
		\begin{equation}
			\E\left[w\, \Dh \dot{f}(X)\right]
			= \E\left[\alpha_h(X)\, \dot{f}(X)\right]
			= \E\left[\alpha_h(X)\,(Y - f(X))\, s(X, Y)\right]
		\end{equation}
		by the tower property. Summing:
		\begin{equation}
			\frac{d\theta}{dt}\Big|_{t=0}
			= \E\left[\psi(X, Y; f, \alpha_h, \theta_h(f))\; s(X, Y)\right]
		\end{equation}
		for every score $s$, so $\theta_h$ is pathwise differentiable with influence
		function $\psi$; since $\psi \in L^2_0(P)$ lies in the tangent space, it is
		the efficient influence function and the bound is
		$\E[\psi^2]$ \citep[Theorem~25.20]{vaart1998}. The cross term in
		$\E[\psi^2]$ vanishes by conditioning on $X$, yielding
		\eqref{eq:effbound}. Attainment follows from Theorem~\ref{thm:main}, whose
		expansion exhibits $\hat{\theta}_h$ as asymptotically linear with influence
		function $\psi$, hence regular and efficient
		\citep[Lemma~25.23]{vaart1998}.
	\end{proof}
	
	% ============================================================
	\section{Proof of Lemma~\ref{lem:sobolev_existence}
		(Derivative Existence, Sobolev)}
	\label{app:lem_sobolev_existence}
	% ============================================================
	
	\subsection{What Needs to Be Shown}
	
	That $\fn \in W^{2,p}(\Omega)$ for $p > d$ implies $\nabla \fn \in
	W^{1,p}(\Omega)^d$ with the norm bound:
	\begin{equation}
		\norm{\nabla \fn}_{L^p}
		\leq \norm{\fn}_{W^{1,p}}
		\leq \norm{\fn}_{W^{2,p}}
	\end{equation}
	
	\begin{remark}
		The condition $p > d$ is required by the proof. It is the condition under
		which Theorem~6.5 of \citet{evans2015measure} guarantees pointwise
		differentiability of $W^{1,p}_{loc}$ functions. This is slightly stronger
		than the condition $p > 1$ stated in the main text and is an explicit
		assumption of the lemma.
	\end{remark}
	
	\subsection{Step 1: Existence of the Gradient as an $L^p$ Function}
	
	Since $\fn \in W^{2,p}(\Omega)$ and $p > d$, we have in particular that
	$\fn \in W^{1,p}(\Omega)$. Extending $\fn$ to a $W^{1,p}_{loc}(\R^d)$
	function by a standard extension operator \citep[Section~4.4]{evans2015measure},
	we may apply Theorem~6.5 of \citet{evans2015measure}:
	
	\begin{quote}
		\textit{Theorem~6.5 \citep{evans2015measure}.}
		Assume $f \in W^{1,p}_{loc}(\R^n)$ for some $n < p \leq \infty$. Then $f$
		is differentiable $\mathcal{L}^n$-a.e., and its classical derivative equals
		its weak derivative $\mathcal{L}^n$-a.e.
	\end{quote}
	
	Applying this theorem to $\fn$ with $n = d$ and $p > d$, we conclude that
	$\fn$ is classically differentiable at $\mathcal{L}^d$-almost every $x \in
	\Omega$, and:
	\begin{equation}
		\nabla \fn(x) = D\fn(x) \quad \mathcal{L}^d\text{-a.e.}
	\end{equation}
	where $D\fn$ denotes the weak gradient of $\fn$. Since $\fn \in W^{1,p}(\Omega)$,
	the weak gradient $D\fn \in L^p(\Omega)^d$ by definition of the Sobolev space.
	Therefore $\nabla \fn \in L^p(\Omega)^d$.
	
	Applying the same argument to each first-order weak derivative $\partial_j \fn
	\in W^{1,p}(\Omega)$ — which belongs to $W^{1,p}(\Omega)$ because $\fn \in
	W^{2,p}(\Omega)$ — we conclude that each $\partial_j \fn$ is classically
	differentiable $\mathcal{L}^d$-a.e.\ with classical second derivatives in
	$L^p(\Omega)$. Therefore $\nabla \fn \in W^{1,p}(\Omega)^d$.
	
	\subsection{Step 2: The Norm Bound}
	
	We establish the chain of inequalities:
	\begin{equation}
		\norm{\nabla \fn}_{L^p}
		\leq \norm{\fn}_{W^{1,p}}
		\leq \norm{\fn}_{W^{2,p}}
	\end{equation}
	
	\subsubsection*{First Inequality}
	
	By definition of the $W^{1,p}(\Omega)$ norm:
	\begin{equation}
		\norm{\fn}_{W^{1,p}}
		= \norm{\fn}_{L^p} + \sum_{j=1}^d \norm{\partial_j \fn}_{L^p}
		= \norm{\fn}_{L^p} + \norm{\nabla \fn}_{L^p}
	\end{equation}
	Since $\norm{\fn}_{L^p} \geq 0$:
	\begin{equation}
		\norm{\fn}_{W^{1,p}}
		= \norm{\fn}_{L^p} + \norm{\nabla \fn}_{L^p}
		\geq \norm{\nabla \fn}_{L^p}
	\end{equation}
	
	\subsubsection*{Second Inequality}
	
	By definition of the $W^{2,p}(\Omega)$ norm using multi-index notation:
	\begin{equation}
		\norm{\fn}_{W^{2,p}}
		= \sum_{|\alpha| \leq 2} \norm{\partial^\alpha \fn}_{L^p}
	\end{equation}
	The multi-indices with $|\alpha| \leq 1$ form a strict subset of those with
	$|\alpha| \leq 2$, and all terms are nonneg­ative, so:
	\begin{equation}
		\norm{\fn}_{W^{2,p}}
		= \sum_{|\alpha| \leq 2} \norm{\partial^\alpha \fn}_{L^p}
		\geq \sum_{|\alpha| \leq 1} \norm{\partial^\alpha \fn}_{L^p}
		= \norm{\fn}_{W^{1,p}}
	\end{equation}
	
	\subsection{Conclusion}
	
	For $p > d$, membership of $\fn$ in $W^{2,p}(\Omega)$ implies via
	Theorem~6.5 of \citet{evans2015measure} that $\nabla \fn$ exists
	$\mathcal{L}^d$-almost everywhere as a classical gradient, coincides with
	the weak gradient, and belongs to $W^{1,p}(\Omega)^d \subset L^p(\Omega)^d$.
	The gradient operator $\nabla: W^{2,p}(\Omega) \to W^{1,p}(\Omega)^d$ is
	bounded with:
	\begin{equation}
		\norm{\nabla \fn}_{L^p}
		\leq \norm{\fn}_{W^{1,p}}
		\leq \norm{\fn}_{W^{2,p}}
	\end{equation}
	\qed
	
	% ============================================================
	\section{Proof of Lemma~\ref{lem:sobolev_consistency}
		(Derivative Consistency, Sobolev)}
	\label{app:lem_sobolev_consistency}
	% ============================================================
	
	\subsection{What Needs to Be Shown}
	
	That well-behavedness in the Sobolev sense implies
	$\norm{\nabla \fn - \nabla f}_{L^p} \xrightarrowp 0$.
	
	\subsection{Proof Strategy}
	
	The proof uses a subsequence argument. We show that every subsequence of
	$\{\fn\}$ has a further subsequence along which
	$\norm{\nabla \fn - \nabla f}_{L^p} \to 0$ almost surely. By the standard
	metric space characterization of convergence in probability --- a sequence
	converges in probability if and only if every subsequence has an almost
	surely convergent further subsequence --- this implies
	$\norm{\nabla \fn - \nabla f}_{L^p} \xrightarrowp 0$.
	
	Suppose for contradiction that $\norm{\nabla \fn - \nabla f}_{L^p}$ does
	not converge to zero in probability. Then there exist $\epsilon > 0$ and
	$\delta > 0$ and a subsequence $\{n_k\}$ such that:
	\begin{equation}
		\mathbb{P}\!\left(\norm{\nabla \hat{f}_{n_k} - \nabla f}_{L^p}
		> \epsilon\right) > \delta
		\quad \forall k
	\end{equation}
	We derive a contradiction by extracting a further subsequence along which
	the gradient converges to zero almost surely.
	
	\subsection{Step 1: Extract a Strongly Convergent Subsequence}
	
	Take any subsequence $\{\hat{f}_{n_k}\}$. By the uniform $\Wtwop$
	boundedness condition:
	\begin{equation}
		\sup_k \norm{\hat{f}_{n_k}}_{W^{2,p}} < \infty \quad \text{almost surely}
	\end{equation}
	The Rellich--Kondrachov compact embedding theorem states that for bounded
	Lipschitz $\Omega$, the embedding $W^{2,p}(\Omega) \hookrightarrow\hookrightarrow
	W^{1,p}(\Omega)$ is compact \citep[Theorem~5,~Section~5.7]{evans2010pde}.
	Compactness of the embedding means that every sequence bounded in $\Wtwop$
	has a subsequence converging strongly in $\Wonep$. Therefore there exists a
	further subsequence $\{\hat{f}_{n_{k_j}}\}$ and a function $g \in \Wonep$
	such that:
	\begin{equation}
		\norm{\hat{f}_{n_{k_j}} - g}_{W^{1,p}} \to 0
		\quad \text{almost surely as } j \to \infty
	\end{equation}
	
	\subsection{Step 2: Identify the Limit}
	
	Strong $\Wonep$ convergence implies strong $L^p$ convergence, since by
	definition of the $W^{1,p}$ norm:
	\begin{equation}
		\norm{\hat{f}_{n_{k_j}} - g}_{L^p}
		\leq \norm{\hat{f}_{n_{k_j}} - g}_{W^{1,p}}
		\to 0 \quad \text{almost surely}
	\end{equation}
	Therefore $\hat{f}_{n_{k_j}} \to g$ in $L^p$ almost surely. By the $L^p$
	consistency assumption, $\norm{\fn - f}_{L^p} \xrightarrowp 0$, so in
	particular along the subsequence $\hat{f}_{n_{k_j}} \to f$ in $L^p$ in
	probability, and hence almost surely along a further subsequence if
	necessary. By uniqueness of limits in $L^p(\Omega)$:
	\begin{equation}
		g = f \quad \mathcal{L}^d\text{-a.e.\ on } \Omega
	\end{equation}
	Therefore the strongly convergent subsequence satisfies:
	\begin{equation}
		\norm{\hat{f}_{n_{k_j}} - f}_{W^{1,p}} \to 0
		\quad \text{almost surely}
	\end{equation}
	
	\subsection{Step 3: Extract Gradient Convergence}
	
	Since $\hat{f}_{n_{k_j}} \to f$ strongly in $\Wonep$, expanding the
	$W^{1,p}$ norm:
	\begin{equation}
		\norm{\hat{f}_{n_{k_j}} - f}_{W^{1,p}}
		= \norm{\hat{f}_{n_{k_j}} - f}_{L^p}
		+ \norm{\nabla \hat{f}_{n_{k_j}} - \nabla f}_{L^p}
		\to 0
		\quad \text{almost surely}
	\end{equation}
	Since both terms are nonneg­ative and their sum goes to zero, each term goes
	to zero individually. In particular:
	\begin{equation}
		\norm{\nabla \hat{f}_{n_{k_j}} - \nabla f}_{L^p} \to 0
		\quad \text{almost surely as } j \to \infty
	\end{equation}
	
	\subsection{Conclusion}
	
	We have shown that from any subsequence $\{\hat{f}_{n_k}\}$ we can extract
	a further subsequence $\{\hat{f}_{n_{k_j}}\}$ along which
	$\norm{\nabla \hat{f}_{n_{k_j}} - \nabla f}_{L^p} \to 0$ almost surely.
	This contradicts the assumption that
	$\mathbb{P}(\norm{\nabla \hat{f}_{n_k} - \nabla f}_{L^p} > \epsilon) > \delta$
	for all $k$. Therefore the assumption was false and:
	\begin{equation}
		\norm{\nabla \fn - \nabla f}_{L^p} \xrightarrowp 0
	\end{equation}
	which is the conclusion of Lemma~\ref{lem:sobolev_consistency}. \qed
	
	\begin{remark}
		The key step is the application of Rellich--Kondrachov in Step~1. The
		theorem requires $\Omega$ to be bounded and Lipschitz --- both of which are
		assumed throughout --- and $p > 1$, which is assumed in the lemma statement.
		The compact embedding $W^{2,p} \hookrightarrow\hookrightarrow W^{1,p}$ is
		strictly stronger than the continuous embedding $W^{2,p} \hookrightarrow
		W^{1,p}$: the latter follows from the norm bound in
		Lemma~\ref{lem:sobolev_existence}, while the former provides the
		compactness needed to extract convergent subsequences. It is this
		compactness, not merely the norm bound, that allows the interchange of
		limit and gradient in the proof.
	\end{remark}
	
	% ============================================================
	\section{Proof of Lemma~\ref{lem:bv_existence}
		(Derivative Existence, BV)}
	\label{app:lem_bv_existence}
	% ============================================================
	
	\subsection{What Needs to Be Shown}
	
	Two things: that $BV$ membership implies $D\fn$ exists as a Radon measure
	with finite total variation, and the explicit formula for piecewise constant
	functions.
	
	\subsection{Step 1: Existence of $D\fn$ as a Radon Measure}
	
	We establish existence in three steps, using the Riesz representation
	theorem and the Radon-Nikodym theorem.
	
	\subsubsection*{Substep 1a: The Distributional Gradient as a Linear Functional}
	
	For any $\fn \in \BV$, the distributional gradient is defined as the linear
	functional $\Lambda: C_c^1(\Omega)^d \to \R$ given by:
	\begin{equation}
		\Lambda(\phi) = -\int_\Omega \fn \operatorname{div} \phi \, dx
		\quad \forall \phi \in C_c^1(\Omega)^d
	\end{equation}
	By the definition of $\BV$, this functional has finite total variation:
	\begin{equation}
		|\Lambda|(\Omega)
		= \sup\left\{
		\Lambda(\phi) : \phi \in C_c^1(\Omega)^d,\, \norm{\phi}_{L^\infty} \leq 1
		\right\}
		< \infty
	\end{equation}
	This is precisely the condition that $\fn \in \BV$ imposes --- the
	distributional gradient is a bounded linear functional on $C_c^1(\Omega)^d$
	with respect to the supremum norm.
	
	\subsubsection*{Substep 1b: Riesz Representation}
	
	Since $\Lambda$ is a bounded linear functional on $C_c^1(\Omega)^d$ with
	finite total variation, the Riesz representation theorem
	\citep[Theorem~1.38]{evans2015measure} guarantees the existence of a
	vector-valued Radon measure $D\fn \in \mathcal{M}(\Omega)^d$ such that:
	\begin{equation}
		\Lambda(\phi) = \int_\Omega \phi \cdot d(D\fn)
		\quad \forall \phi \in C_c^1(\Omega)^d
	\end{equation}
	That is:
	\begin{equation}
		-\int_\Omega \fn \operatorname{div} \phi \, dx
		= \int_\Omega \phi \cdot d(D\fn)
		\quad \forall \phi \in C_c^1(\Omega)^d
	\end{equation}
	This is the distributional gradient equation, and $D\fn$ is the
	distributional gradient of $\fn$ represented as a Radon measure.
	
	\subsubsection*{Substep 1c: Finiteness of Total Variation}
	
	By the Radon-Nikodym theorem for Radon measures
	\citep[Theorem~1.29]{evans2015measure}, the derivative $D_{\mathcal{L}^d}(D\fn)$
	exists and is finite $\mathcal{L}^d$-almost everywhere on $\Omega$. Since
	$\Omega$ is bounded and the total variation of $\Lambda$ is finite by the
	$BV$ assumption, we conclude:
	\begin{equation}
		|D\fn|(\Omega) = |\Lambda|(\Omega) < \infty
	\end{equation}
	Therefore $D\fn \in \mathcal{M}(\Omega)^d$ with $|D\fn|(\Omega) < \infty$,
	as claimed.
	
	\subsection{Step 2: Explicit Formula for Piecewise Constant Functions}
	
	Let $\fn = \sum_j \bar{y}_j \mathbf{1}_{R_j}$ where $\{R_j\}$ is a finite
	partition of $\Omega$ into sets with Lipschitz boundaries, and $\bar{y}_j$
	denotes the response mean on leaf $R_j$. We derive the explicit formula for
	$D\fn$ by applying the Gauss--Green theorem to each leaf.
	
	\subsubsection*{Substep 2a: Applying the Definition}
	
	For any test function $\phi \in C_c^1(\Omega)^d$, the distributional
	gradient equation from Substep~1b gives:
	\begin{equation}
		\int_\Omega \phi \cdot d(D\fn)
		= -\int_\Omega \fn \operatorname{div} \phi \, dx
		= -\sum_j \bar{y}_j \int_{R_j} \operatorname{div} \phi \, dx
	\end{equation}
	where the second equality uses the piecewise constant structure of $\fn$ and
	linearity of the integral.
	
	\subsubsection*{Substep 2b: Gauss--Green on Each Leaf}
	
	Since each $R_j$ has Lipschitz boundary, the Gauss--Green theorem for $BV$
	functions \citep[Theorem~5.16]{evans2015measure} applies to each leaf:
	\begin{equation}
		\int_{R_j} \operatorname{div} \phi \, dx
		= \int_{\partial R_j} \phi \cdot \nu_j \, d\mathcal{H}^{d-1}
	\end{equation}
	where $\nu_j$ denotes the outward unit normal to $\partial R_j$ and
	$\mathcal{H}^{d-1}$ is the $(d-1)$-dimensional Hausdorff measure on the
	boundary.
	
	\subsubsection*{Substep 2c: Cancellation on Interior Boundaries}
	
	Substituting the Gauss--Green formula into the expression from Substep~2a:
	\begin{equation}
		\int_\Omega \phi \cdot d(D\fn)
		= -\sum_j \bar{y}_j \int_{\partial R_j} \phi \cdot \nu_j \, d\mathcal{H}^{d-1}
	\end{equation}
	Every interior boundary $\partial R_j \cap \partial R_k$ between adjacent
	leaves $j$ and $k$ appears twice in this sum --- once with outward normal
	$\nu_j$ for leaf $R_j$ and once with outward normal $\nu_k = -\nu_j$ for
	leaf $R_k$. Collecting these paired contributions and writing $\nu_{jk} =
	\nu_j$ for the unit normal pointing from $R_j$ into $R_k$:
	\begin{align}
		-\sum_j \bar{y}_j \int_{\partial R_j} \phi \cdot \nu_j \, d\mathcal{H}^{d-1}
		&= -\sum_{j \sim k} \bar{y}_j \int_{\partial R_j \cap \partial R_k}
		\phi \cdot \nu_{jk} \, d\mathcal{H}^{d-1} \notag \\
		&\quad + \sum_{j \sim k} \bar{y}_k \int_{\partial R_j \cap \partial R_k}
		\phi \cdot \nu_{jk} \, d\mathcal{H}^{d-1} \notag \\
		&= \sum_{j \sim k} (\bar{y}_k - \bar{y}_j)
		\int_{\partial R_j \cap \partial R_k} \phi \cdot \nu_{jk}
		\, d\mathcal{H}^{d-1}
	\end{align}
	Boundary terms on $\partial\Omega$ vanish because $\phi \in C_c^1(\Omega)^d$
	has compact support in $\Omega$.
	
	\subsubsection*{Substep 2d: Identifying the Measure}
	
	Since the equation:
	\begin{equation}
		\int_\Omega \phi \cdot d(D\fn)
		= \sum_{j \sim k} (\bar{y}_k - \bar{y}_j)
		\int_{\partial R_j \cap \partial R_k} \phi \cdot \nu_{jk}
		\, d\mathcal{H}^{d-1}
	\end{equation}
	holds for all $\phi \in C_c^1(\Omega)^d$, the uniqueness part of the Riesz
	representation theorem \citep[Theorem~1.38]{evans2015measure} identifies
	$D\fn$ uniquely as the measure:
	\begin{equation}
		D\fn = \sum_{j \sim k} (\bar{y}_k - \bar{y}_j) \cdot \nu_{jk} \cdot
		\mathcal{H}^{d-1}\lfloor_{\partial R_j \cap \partial R_k}
	\end{equation}
	where $\mathcal{H}^{d-1}\lfloor_{\partial R_j \cap \partial R_k}$ denotes
	the $(d-1)$-dimensional Hausdorff measure restricted to the shared boundary
	between leaves $j$ and $k$. This is the formula stated in the lemma.
	\qed
	
	% ============================================================
	\section{Proof of Lemma~\ref{lem:bv_consistency}
		(Derivative Consistency, BV)}
	\label{app:lem_bv_consistency}
	% ============================================================
	
	\subsection{What Needs to Be Shown}
	
	That well-behavedness in the BV sense implies $D\fn \xrightarrow{w*} Df$
	in $\cM$, and the strict convergence upgrade under the additional condition
	$\E[|D\fn|(\Omega)] \to |Df|(\Omega)$.
	
	\subsection{Proof Strategy}
	
	The proof mirrors the subsequence argument of
	Lemma~\ref{lem:sobolev_consistency}, replacing the Rellich--Kondrachov
	compact embedding with the BV compactness theorem and replacing strong $L^p$
	convergence with weak-* convergence of measures. We show that every
	subsequence of $\{D\fn\}$ has a further subsequence converging weak-* to
	$Df$. By the metric space characterization of convergence in probability,
	this implies $D\fn \xrightarrow{w*} Df$ in $\cM$ in probability.
	
	Suppose for contradiction that $D\fn$ does not converge weak-* to $Df$ in
	probability. Then there exist $\epsilon > 0$, $\delta > 0$, a test function
	$\phi \in C_0(\Omega)^d$ with $\norm{\phi}_\infty \leq 1$, and a subsequence
	$\{n_k\}$ such that:
	\begin{equation}
		\mathbb{P}\!\left(\left|
		\int_\Omega \phi \cdot d(D\hat{f}_{n_k})
		- \int_\Omega \phi \cdot d(Df)
		\right| > \epsilon\right) > \delta
		\quad \forall k
	\end{equation}
	We derive a contradiction by extracting a further subsequence along which
	weak-* convergence holds almost surely.
	
	\subsection{Part I: Weak-* Convergence}
	
	\subsubsection*{Step 1: Bound the Full BV Norm}
	
	To apply the BV compactness theorem we require a uniform bound on the full
	$BV$ norm:
	\begin{equation}
		\norm{\fn}_{BV} = \norm{\fn}_{L^1} + |D\fn|(\Omega)
	\end{equation}
	For the total variation term, the well-behaved condition gives:
	\begin{equation}
		\sup_n \E\!\left[|D\fn|(\Omega)\right] < \infty
	\end{equation}
	so $|D\fn|(\Omega)$ is bounded in expectation and hence bounded almost
	surely along subsequences by Markov's inequality.
	
	For the $L^1$ term, the consistency assumption
	$\norm{\fn - f}_{L^1} \xrightarrowp 0$ implies that $\norm{\fn}_{L^1}$ is
	eventually bounded in probability. Since $f \in \Wonep \subset \Lone$, we
	have $\norm{f}_{L^1} < \infty$, and:
	\begin{equation}
		\norm{\fn}_{L^1}
		\leq \norm{\fn - f}_{L^1} + \norm{f}_{L^1}
	\end{equation}
	which is bounded in probability. Combining the two terms:
	\begin{equation}
		\sup_n \norm{\fn}_{BV} < \infty \quad \text{almost surely}
	\end{equation}
	
	\subsubsection*{Step 2: Extract a Weakly-* Convergent Subsequence}
	
	Take any subsequence $\{\hat{f}_{n_k}\}$. By Step~1 it is bounded in
	$\BV$ almost surely. The BV compactness theorem
	\citep[Theorem~5.5]{evans2015measure} states that every sequence bounded in
	$BV(\Omega)$ for bounded Lipschitz $\Omega$ has a subsequence
	$\{\hat{f}_{n_{k_j}}\}$ and $g \in \BV$ such that:
	\begin{enumerate}
		\item $\hat{f}_{n_{k_j}} \to g$ strongly in $\Lone$
		\item $D\hat{f}_{n_{k_j}} \xrightarrow{w*} Dg$ in $\cM$
	\end{enumerate}
	both holding almost surely. This is the BV analog of the Rellich--Kondrachov
	compact embedding: the BV compactness theorem gives simultaneously the $L^1$
	convergence of the functions and the weak-* convergence of their gradient
	measures, in a single extraction.
	
	\subsubsection*{Step 3: Identify the Limit}
	
	From item~(1) of Step~2, $\hat{f}_{n_{k_j}} \to g$ strongly in $\Lone$
	almost surely. By the $L^1$ consistency assumption
	$\norm{\fn - f}_{L^1} \xrightarrowp 0$, along the subsequence
	$\hat{f}_{n_{k_j}} \to f$ in $L^1$ in probability, hence almost surely
	along a further subsequence. By uniqueness of $L^1$ limits:
	\begin{equation}
		g = f \quad \mathcal{L}^d\text{-a.e.\ on } \Omega
	\end{equation}
	Therefore $Dg = Df$ as Radon measures, and from item~(2) of Step~2:
	\begin{equation}
		D\hat{f}_{n_{k_j}} \xrightarrow{w*} Df \quad \text{in } \cM
		\quad \text{almost surely}
	\end{equation}
	In particular, for the test function $\phi$ from the contradiction
	assumption:
	\begin{equation}
		\int_\Omega \phi \cdot d(D\hat{f}_{n_{k_j}})
		\to \int_\Omega \phi \cdot d(Df)
		\quad \text{almost surely}
	\end{equation}
	
	\subsubsection*{Step 4: Derive the Contradiction}
	
	The almost sure convergence in Step~3 implies:
	\begin{equation}
		\mathbb{P}\!\left(\left|
		\int_\Omega \phi \cdot d(D\hat{f}_{n_{k_j}})
		- \int_\Omega \phi \cdot d(Df)
		\right| > \epsilon\right) \to 0
	\end{equation}
	This contradicts the assumption that the probability exceeds $\delta > 0$
	for all $k$. Therefore the assumption was false.
	
	\subsubsection*{Conclusion of Part I}
	
	Since every subsequence $\{D\hat{f}_{n_k}\}$ has a further subsequence
	converging weak-* to $Df$ almost surely, the full sequence satisfies:
	\begin{equation}
		D\fn \xrightarrow{w*} Df \quad \text{in } \cM
		\quad \text{in probability}
	\end{equation}
	
	\subsection{Part II: Strict Convergence Upgrade}
	
	We now prove the additional claim: if $\E[|D\fn|(\Omega)] \to |Df|(\Omega)$
	then $D\fn \to Df$ strictly in $\cM$, meaning weak-* convergence holds and
	$|D\fn|(\Omega) \to |Df|(\Omega)$.
	
	\subsubsection*{Step 1: Lower Semicontinuity Gives One Direction}
	
	The total variation functional $\mu \mapsto |\mu|(\Omega)$ is lower
	semicontinuous with respect to weak-* convergence of Radon measures
	\citep[Theorem~5.2]{evans2015measure}. Since $D\fn \xrightarrow{w*} Df$
	in $\cM$ by Part~I:
	\begin{equation}
		|Df|(\Omega) \leq \liminf_{n \to \infty} |D\fn|(\Omega)
	\end{equation}
	This gives the lower bound on the total variation of the limit.
	
	\subsubsection*{Step 2: The Additional Condition Gives the Other Direction}
	
	By the additional assumption $\E[|D\fn|(\Omega)] \to |Df|(\Omega)$, we have:
	\begin{equation}
		\limsup_{n \to \infty} |D\fn|(\Omega) \leq |Df|(\Omega)
	\end{equation}
	in expectation, hence in probability along subsequences.
	
	\subsubsection*{Step 3: Pin the Total Variation}
	
	Step~1 gives $|Df|(\Omega) \leq \liminf_{n \to \infty} |D\fn|(\Omega)$ by
	lower semicontinuity of total variation under weak-* convergence. Step~2
	gives $\limsup_{n \to \infty} |D\fn|(\Omega) \leq |Df|(\Omega)$ from the
	additional assumption $\E[|D\fn|(\Omega)] \to |Df|(\Omega)$. Since
	$\liminf \leq \limsup$ always, combining gives:
	\begin{equation}
		|Df|(\Omega)
		\leq \liminf_{n \to \infty} |D\fn|(\Omega)
		\leq \limsup_{n \to \infty} |D\fn|(\Omega)
		\leq |Df|(\Omega)
	\end{equation}
	All inequalities are therefore equalities and $|D\fn|(\Omega) \to |Df|(\Omega)$
	in probability. Combined with the weak-* convergence $D\fn \xrightarrow{w*}
	Df$ established in Part~I, this is precisely the definition of strict
	convergence in $\cM$ \citep[Definition~3.13]{ambrosio2000bv}. \qed
	
	\begin{remark}
		The parallel between Lemma~\ref{lem:sobolev_consistency} and
		Lemma~\ref{lem:bv_consistency} is exact at the structural level. In both
		cases: uniform boundedness in a strong space gives compactness via a
		compact embedding or compactness theorem; the consistency condition
		identifies the limit; and the subsequence argument promotes deterministic
		convergence to convergence in probability. The difference is entirely in
		the function spaces and the topology of convergence --- strong $L^p$ in the
		Sobolev case, weak-* of measures in the BV case --- reflecting the
		fundamental difference in the inductive biases of smooth and piecewise
		constant estimators.
	\end{remark}
	
	\begin{remark}
		The strict convergence upgrade requires the additional assumption
		$\E[|D\fn|(\Omega)] \to |Df|(\Omega)$, which is not implied by weak-*
		convergence alone. This is because weak-* convergence of measures does not
		in general imply convergence of total masses --- mass can escape to the
		boundary or concentrate in ways invisible to continuous test functions. The
		lower semicontinuity of total variation \citep[Theorem~5.2]{evans2015measure}
		gives the liminf inequality for free; pinning the limsup requires the
		additional assumption. In applications where only weak-* convergence is
		needed --- as in the bootstrap corollaries --- the additional assumption
		is not required.
	\end{remark}
	
	% ============================================================
	\section{Proof of Theorem~\ref{thm:pairing} (Pairing Identity)}
	\label{app:thm_pairing}
	% ============================================================
	
	\begin{proof}
		Write points as $x = (x_{-j}, x_j)$ and for fixed $x_{-j}$ let
		$\Omega^{x_{-j}} = \{ \tau \in \R : (x_{-j}, \tau) \in \Omega \}$ denote the
		one-dimensional section of the domain and
		$g^{x_{-j}}(\tau) = g(x_{-j}, \tau)$ the section of $g$. By the
		one-dimensional section theory of BV functions
		\citep[Section~3.11]{ambrosio2000bv}, for $\mathcal{H}^{d-1}$-almost every
		$x_{-j}$ the section $g^{x_{-j}}$ of the precise representative is a good
		representative of a function of bounded variation on $\Omega^{x_{-j}}$, its
		distributional derivative $D g^{x_{-j}}$ is the corresponding section of
		$D_j g$, and the total variations disintegrate:
		\begin{equation}
			\int_\Omega \phi\; d D_j g
			= \int \left( \int_{\Omega^{x_{-j}}} \phi(x_{-j}, \tau)\;
			d D g^{x_{-j}}(\tau) \right) d\mathcal{H}^{d-1}(x_{-j})
			\label{eq:disintegration}
		\end{equation}
		for every bounded Borel $\phi$. For a good representative of a BV function of
		one variable, the fundamental theorem of calculus holds in the form
		$g^{x_{-j}}(b) - g^{x_{-j}}(a) = D g^{x_{-j}}\bigl((a, b]\bigr)$ for all
		$a < b$ in $\Omega^{x_{-j}}$ that are not atoms of $D g^{x_{-j}}$; since
		each section has at most countably many atoms, for
		Lebesgue-almost every $x_j$ the pair $(x_j - h, x_j + h)$ avoids them. Hence
		for almost every $x \in \operatorname{supp}(w)$ (where the whole segment lies
		in $\Omega$ by the trimming condition):
		\begin{equation}
			g(x + h e_j) - g(x - h e_j)
			= D g^{x_{-j}}\bigl((x_j - h,\; x_j + h]\bigr)
		\end{equation}
		Multiplying by $q(x)/2h$, integrating, and applying
		Tonelli on each section (justified since
		$\abs{D g^{x_{-j}}}$ is finite for a.e.\ section and $q$ is bounded):
		\begin{align}
			\E\left[w(X)\, \Dh g(X)\right]
			&= \int d\mathcal{H}^{d-1}(x_{-j}) \int dx_j\;
			\frac{q(x_{-j}, x_j)}{2h}
			\int \mathbf{1}\{\tau \in (x_j - h, x_j + h]\}\;
			d D g^{x_{-j}}(\tau) \notag \\
			&= \int d\mathcal{H}^{d-1}(x_{-j}) \int d D g^{x_{-j}}(\tau)\;
			\frac{1}{2h} \int \mathbf{1}\{x_j \in [\tau - h, \tau + h)\}\,
			q(x_{-j}, x_j)\; dx_j \notag \\
			&= \int d\mathcal{H}^{d-1}(x_{-j}) \int (K_h \star q)(x_{-j}, \tau)\;
			d D g^{x_{-j}}(\tau)
			= \int_\Omega (K_h \star q)\; d D_j g
		\end{align}
		using $\tau \in (x_j - h, x_j + h] \Leftrightarrow
		x_j \in [\tau - h, \tau + h)$ in the second line, the definition of
		$K_h \star q$ (half-open versus closed windows agree up to Lebesgue-null
		sets) in the third, and \eqref{eq:disintegration} in the last step. The
		piecewise constant case follows by substituting the explicit jump formula of
		Lemma~\ref{lem:bv_existence} for $D_j g$.
		
		For the corollary: if $q$ extends to a uniformly continuous function of
		compact support on $\R^d$, then $K_h \star q \to q$ uniformly as
		$h \downarrow 0$, so:
		\begin{equation}
			\abs{\int (K_h \star q)\; d D_j g - \int q\; d D_j g}
			\leq \norm{K_h \star q - q}_\infty \cdot \abs{D_j g}(\Omega) \to 0
		\end{equation}
		since $\abs{D_j g}(\Omega) < \infty$. When $g \in W^{1,1}(\Omega)$, the
		measure $D_j g$ is absolutely continuous with density $\partial_j g$, so
		$\int q\, d D_j g = \int w\, \partial_j g\, p\, dx
		= \E[w\, \partial_j g(X)]$.
	\end{proof}
	
	% ============================================================
	\section{Proof of Theorem~\ref{thm:neural_nets} (Neural Networks)}
	\label{app:neural_nets}
	% ============================================================
	
	\subsection{Setup and Notation}
	
	Let $\fn: \Omega \to \R$ be a feedforward neural network with $L$ layers,
	weight matrices $W_\ell^{(n)} \in \R^{d_\ell \times d_{\ell-1}}$, bias vectors
	$b_\ell^{(n)} \in \R^{d_\ell}$, and activation function $\sigma \in C^2(\R)$.
	The network computes:
	\begin{equation}
		\fn(x) = W_L^{(n)} z_{L-1} + b_L^{(n)}
	\end{equation}
	where the intermediate representations are defined recursively by:
	\begin{equation}
		z_\ell = \sigma\!\left(W_\ell^{(n)} z_{\ell-1} + b_\ell^{(n)}\right),
		\quad z_0 = x
	\end{equation}
	with $\sigma$ applied elementwise. Denote the pre-activation at layer $\ell$
	by $a_\ell = W_\ell^{(n)} z_{\ell-1} + b_\ell^{(n)}$ and the diagonal
	matrix of activation derivatives by $\Sigma'_\ell = \operatorname{diag}
	(\sigma'(a_\ell^{(1)}), \ldots, \sigma'(a_\ell^{(d_\ell)}))$ and similarly
	$\Sigma''_\ell = \operatorname{diag}(\sigma''(a_\ell^{(1)}), \ldots,
	\sigma''(a_\ell^{(d_\ell)}))$.
	
	Define the product of weight matrices from layer $k$ to layer $\ell$ as:
	\begin{equation}
		M_{\ell:k} = W_\ell^{(n)} \Sigma'_{\ell-1} W_{\ell-1}^{(n)}
		\Sigma'_{\ell-2} \cdots W_{k+1}^{(n)} \Sigma'_k
		\quad \text{for } \ell > k
	\end{equation}
	with the convention $M_{\ell:\ell} = \Sigma'_\ell$. This notation captures
	the Jacobian of the composition of layers from $k$ to $\ell$.
	
	The proof proceeds in two steps. Step~1 establishes uniform $\Wtwop$
	boundedness. Step~2 applies Lemma~\ref{lem:sobolev_consistency} to conclude.
	
	\subsection{Step 1: Uniform $W^{2,p}$ Boundedness}
	
	\textbf{Claim:} $\sup_n \norm{\fn}_{\Wtwop} < \infty$.
	
	\subsubsection*{Substep 1a: Bound on $\fn$ in $L^p$}
	
	Since $\norm{\fn - f}_{L^p} \xrightarrowp 0$ and $f \in \Wonep \subset \Lp$,
	the triangle inequality gives:
	\begin{equation}
		\norm{\fn}_{L^p} \leq \norm{\fn - f}_{L^p} + \norm{f}_{L^p}
	\end{equation}
	The first term converges to zero in probability and the second is finite
	since $f \in \Lp$. Therefore $\norm{\fn}_{L^p}$ is eventually bounded in
	probability and this bound can be made uniform in $n$ by passing to
	a subsequence or imposing a mild boundedness condition on the network output.
	
	\subsubsection*{Substep 1b: First Derivative Bound}
	
	Differentiating $\fn$ with respect to the input $x$ via the chain rule
	through all $L$ layers, the gradient at $x \in \Omega$ is:
	\begin{equation}
		\nabla \fn(x) = \left(W_L^{(n)} M_{L-1:1}\right)^\top
	\end{equation}
	where $M_{L-1:1} = W_{L-1}^{(n)} \Sigma'_{L-2} W_{L-2}^{(n)} \cdots
	\Sigma'_1 W_1^{(n)}$ is the product of weight-activation Jacobians from
	layer 1 to layer $L-1$. More explicitly, for the $j$-th input coordinate:
	\begin{equation}
		\frac{\partial \fn}{\partial x_j}(x)
		= \left(W_L^{(n)}\right)_{:} \cdot
		\prod_{\ell=L-1}^{1} \left(W_\ell^{(n)} \Sigma'_{\ell-1}\right) \cdot e_j
	\end{equation}
	where $e_j$ is the $j$-th standard basis vector.
	
	Taking the operator norm of this product and applying submultiplicativity:
	\begin{equation}
		\abs{\frac{\partial \fn}{\partial x_j}(x)}
		\leq \norm{W_L^{(n)}} \cdot \prod_{\ell=1}^{L-1}
		\norm{W_\ell^{(n)}} \cdot \norm{\Sigma'_\ell}
		\leq \prod_{\ell=1}^{L} \norm{W_\ell^{(n)}} \cdot \norm{\sigma'}_\infty^{L-1}
	\end{equation}
	using that $\norm{\Sigma'_\ell} = \max_i |\sigma'(a_\ell^{(i)})| \leq
	\norm{\sigma'}_\infty < \infty$ since $\sigma \in C^2$. This bound holds
	pointwise for every $x \in \Omega$, so integrating over $\Omega$:
	\begin{equation}
		\norm{\nabla \fn}_{L^p}
		\leq \norm{\nabla \fn}_{L^\infty} \cdot |\Omega|^{1/p}
		\leq \prod_\ell \norm{W_\ell^{(n)}} \cdot \norm{\sigma'}_\infty^{L-1}
		\cdot |\Omega|^{1/p}
	\end{equation}
	
	\subsubsection*{Substep 1c: Second Derivative Bound}
	
	Differentiating the gradient formula once more with respect to $x_i$ using
	the product rule and chain rule, the $(i,j)$ entry of the Hessian
	$\nabla^2 \fn(x)$ is:
	\begin{equation}
		\frac{\partial^2 \fn}{\partial x_i \partial x_j}(x)
		= \sum_{\ell=1}^{L-1}
		\left(W_L^{(n)} M_{L-1:\ell+1}\right)_{:}
		\cdot \Sigma''_\ell \cdot
		\left(M_{\ell:1} e_j\right) \cdot
		\left(M_{\ell:1} e_i\right)^\top \cdot
		\left(W_L^{(n)} M_{L-1:\ell+1}\right)_{:}^\top
	\end{equation}
	where the sum runs over all layers $\ell$ at which the second derivative of
	$\sigma$ is applied. The structure is: the forward pass from the input to
	layer $\ell$ contributes $M_{\ell:1} e_j$ for the $j$-th coordinate
	direction, the activation second derivative at layer $\ell$ contributes
	$\Sigma''_\ell$, and the backward pass from layer $\ell$ to the output
	contributes $W_L^{(n)} M_{L-1:\ell+1}$.
	
	Taking absolute values and applying submultiplicativity of operator norms:
	\begin{align}
		\abs{\frac{\partial^2 \fn}{\partial x_i \partial x_j}(x)}
		&\leq \sum_{\ell=1}^{L-1}
		\norm{W_L^{(n)} M_{L-1:\ell+1}}^2 \cdot
		\norm{\Sigma''_\ell} \cdot
		\norm{M_{\ell:1}}^2 \notag \\
		&\leq \sum_{\ell=1}^{L-1}
		\left(\prod_{k=\ell+1}^{L} \norm{W_k^{(n)}}\right)^2 \cdot
		\norm{\sigma''}_\infty \cdot
		\left(\prod_{k=1}^{\ell} \norm{W_k^{(n)}}\right)^2 \cdot
		\norm{\sigma'}_\infty^{2(L-2)} \notag \\
		&\leq (L-1) \cdot \left(\prod_{\ell=1}^{L} \norm{W_\ell^{(n)}}\right)^2
		\cdot \norm{\sigma''}_\infty \cdot \norm{\sigma'}_\infty^{2(L-2)}
	\end{align}
	using $\norm{\Sigma''_\ell} \leq \norm{\sigma''}_\infty < \infty$ since
	$\sigma \in C^2$, and collecting all weight norm factors into a single
	product. This bound holds pointwise for every $x \in \Omega$, so:
	\begin{equation}
		\norm{\nabla^2 \fn}_{L^p}
		\leq \norm{\nabla^2 \fn}_{L^\infty} \cdot |\Omega|^{1/p}
		\leq (L-1) \cdot \left(\prod_\ell \norm{W_\ell^{(n)}}\right)^2
		\cdot \norm{\sigma''}_\infty \cdot \norm{\sigma'}_\infty^{2(L-2)}
		\cdot |\Omega|^{1/p}
	\end{equation}
	
	\subsubsection*{Combining Substeps 1a--1c}
	
	Summing the three bounds:
	\begin{equation}
		\norm{\fn}_{W^{2,p}}
		= \norm{\fn}_{L^p} + \norm{\nabla \fn}_{L^p} + \norm{\nabla^2 \fn}_{L^p}
		\leq C\!\left(\prod_\ell \norm{W_\ell^{(n)}},\,
		\norm{\sigma'}_\infty,\, \norm{\sigma''}_\infty,\, L,\, |\Omega|\right)
	\end{equation}
	where $C$ is a finite constant depending on the indicated quantities but not
	on $n$ or $x$. By the assumption $\sup_n \prod_\ell \norm{W_\ell^{(n)}}
	< \infty$ --- which formalizes the effect of $L^2$ regularization during
	training \citep{bartlett2017spectrally} --- and the finiteness of
	$\norm{\sigma'}_\infty$ and $\norm{\sigma''}_\infty$ from $\sigma \in C^2$,
	the right-hand side is uniformly bounded in $n$. Therefore:
	\begin{equation}
		\sup_n \norm{\fn}_{W^{2,p}} < \infty
	\end{equation}
	This completes Step~1.
	
	\subsection{Step 2: $L^p$ Consistency}
	
	The $L^p$ consistency condition $\norm{\fn - f}_{L^p} \xrightarrowp 0$ is
	the standard universal approximation result for smooth neural networks.
	\citet{hornik1990universal} establish that multilayer feedforward networks
	with smooth activation functions are dense in $W^{1,p}(\Omega)$ with respect
	to the $W^{1,p}$ norm, implying in particular density in $L^p(\Omega)$.
	Combined with standard empirical risk minimization results establishing that
	the training procedure converges to the best approximation in the function
	class, this gives $\norm{\fn - f}_{L^p} \xrightarrowp 0$.
	
	\subsection{Step 3: Applying Lemma~\ref{lem:sobolev_consistency}}
	
	By Step~1, $\sup_n \norm{\fn}_{W^{2,p}} < \infty$. By Step~2,
	$\norm{\fn - f}_{L^p} \xrightarrowp 0$. These are precisely the two
	conditions of the well-behaved definition in the Sobolev sense
	(Definition~\ref{def:wellbehaved_sobolev}). Lemma~\ref{lem:sobolev_consistency}
	then gives:
	\begin{equation}
		\norm{\nabla \fn - \nabla f}_{L^p} \xrightarrowp 0
	\end{equation}
	which is the conclusion of Theorem~\ref{thm:neural_nets}. 
	\qed
	
	\begin{remark}
		The weight norm assumption $\sup_n \prod_\ell \norm{W_\ell^{(n)}} < \infty$
		is the key regularity condition on the estimator. It formalizes the effect
		of $L^2$ weight decay regularization during training, which penalizes large
		weights and keeps the product of operator norms bounded. Under standard
		conditions on the regularization parameter and optimization procedure,
		this bound holds with high probability over the training randomness
		\citep{bartlett2017spectrally}. The condition is stated as an assumption
		of the theorem rather than derived from first principles because its
		verification depends on details of the training procedure that are outside
		the scope of the present analysis.
	\end{remark}
	
	\begin{remark}
		Theorem~\ref{thm:neural_nets} is stated for smooth activations $\sigma \in
		C^2$. For ReLU networks $\sigma \notin C^2$ and the Hessian calculation in
		Substep~1c does not apply directly. The first derivative exists almost
		everywhere with $\sigma' = \mathbf{1}_{(0,\infty)}$ bounded, so Substep~1b
		goes through unchanged. The distributional second derivative is $\sigma'' =
		\delta_0$, a Dirac mass at zero, placing $\fn$ in a mixed $W^{1,p} \cap BV$
		space rather than $\Wtwop$. A modified version of
		Lemma~\ref{lem:sobolev_consistency} applicable in this mixed space, or a
		restriction to input distributions placing zero mass on the kink set
		$\{x : a_\ell(x) = 0\}$, would be required to handle this case. We leave
		this extension for future work.
	\end{remark}
	
	% ============================================================
	\section{Proof of Theorem~\ref{thm:boosting} (Gradient Boosting)}
	\label{app:boosting}
	% ============================================================
	
	\subsection{Setup and Notation}
	
	By Proposition~1 of \citet{buhlmann2003boosting}, the $L_2$Boost estimator after
	$m$ iterations can be represented as:
	\begin{equation}
		\hat{F}_m = (I - (I-S)^{m+1})Y
	\end{equation}
	where $S: \mathbb{R}^n \to \mathbb{R}^n$ is the base learner operator. The residual
	at step $m$ is:
	\begin{equation}
		u_m = Y - \hat{F}_m = (I-S)^m Y
	\end{equation}
	so that each base learner fit is $\hat{f}_m = Su_m$ and the ensemble updates as
	$\hat{F}_m = \hat{F}_{m-1} + \hat{f}_m$.
	
	By Proposition~2 of \citet{buhlmann2003boosting}, for a symmetric linear learner $S$
	with eigenvalues $\{\lambda_k; k = 1,\ldots,n\}$, the eigenvalues of the boosting
	operator $\mathcal{B}_m = I - (I-S)^{m+1}$ are $\{1-(1-\lambda_k)^{m+1}; k =
	1,\ldots,n\}$.
	
	Under the eigenvalue condition $0 < \lambda_k \leq 1$ for all $k$
	\citep[Theorem~1]{buhlmann2003boosting}, let:
	\begin{equation}
		\rho = \max_k (1-\lambda_k) \in [0,1)
	\end{equation}
	Since $\lambda_k > 0$ for all $k$, we have $\rho < 1$. The bias decays
	exponentially fast: by Theorem~1(a) of \citet{buhlmann2003boosting}:
	\begin{equation}
		\norm{u_m}_{L^2}^2 = \norm{(I-S)^m Y}_{L^2}^2
		\leq \rho^{2m} \norm{Y}_{L^2}^2
	\end{equation}
	In particular:
	\begin{equation}
		\norm{u_m}_{L^2} \leq C\rho^m, \quad C = \norm{Y}_{L^2}
		\label{eq:geometric_decay}
	\end{equation}
	This geometric decay is the key engine for both cases of the proof.
	
	The two cases of the theorem are handled separately.
	
	\subsection{Case (i): Smooth Base Learners}
	
	\subsubsection*{Step 1: $W^{2,p}$ Bound for a Single Base Learner}
	
	Each $\hat{f}_m = Su_m \in \Wtwop$ by assumption on the base learner class. For
	shallow neural net base learners this follows from Appendix~\ref{app:neural_nets}.
	For spline base learners it follows from standard spline approximation theory
	\citep{wahba1990splines}. We require a uniform bound:
	\begin{equation}
		\sup_m \norm{\hat{f}_m}_{W^{2,p}} \leq C' \norm{u_m}_{L^2} < \infty
	\end{equation}
	where the bound on $\norm{\hat{f}_m}_{W^{2,p}}$ in terms of $\norm{u_m}_{L^2}$
	follows from the boundedness of $S$ as an operator from $L^2$ to $W^{2,p}$.
	
	\subsubsection*{Step 2: Ensemble Bound via Geometric Decay}
	
	By linearity of the $W^{2,p}$ norm and the triangle inequality:
	\begin{equation}
		\norm{\hat{F}_m}_{W^{2,p}}
		\leq \sum_{j=1}^m \norm{\hat{f}_j}_{W^{2,p}}
		\leq C' \sum_{j=1}^m \norm{u_j}_{L^2}
		\leq C'C \sum_{j=1}^m \rho^j
		\leq \frac{C'C\rho}{1-\rho}
		< \infty
	\end{equation}
	using the geometric decay \eqref{eq:geometric_decay} and the fact that
	$\sum_{j=1}^\infty \rho^j = \rho/(1-\rho) < \infty$ for $\rho < 1$. The bound is
	uniform in $m$.
	
	\subsubsection*{Step 3: Apply Lemma~\ref{lem:sobolev_consistency}}
	
	$L^p$ consistency is assumed. Uniform $\Wtwop$ boundedness follows from Step~2.
	Lemma~\ref{lem:sobolev_consistency} gives
	$\norm{\nabla \hat{F}_m - \nabla f}_{L^p} \xrightarrowp 0$. \qed
	
	\subsection{Case (ii): Tree Stump Base Learners}
	
	Each base learner $\hat{f}_m = Su_m$ is a depth-1 tree (stump) of the form:
	\begin{equation}
		\hat{f}_m(x) = c_m^L \,\mathbf{1}[x_j \leq t_m] + c_m^R \,\mathbf{1}[x_j > t_m]
	\end{equation}
	for some split variable $j$, threshold $t_m$, and leaf constants:
	\begin{equation}
		c_m^L = \frac{1}{|I_L|}\sum_{i \in I_L} u_m(x_i), \quad
		c_m^R = \frac{1}{|I_R|}\sum_{i \in I_R} u_m(x_i)
	\end{equation}
	where $I_L = \{i : x_{ij} \leq t_m\}$ and $I_R = \{i : x_{ij} > t_m\}$ are the
	left and right leaf index sets.
	
	\subsubsection*{Step 1: BV Norm of a Single Stump}
	
	By Lemma~\ref{lem:bv_existence} the distributional gradient of $\hat{f}_m$ is:
	\begin{equation}
		D\hat{f}_m = (c_m^R - c_m^L) \cdot \nu_m \cdot
		\mathcal{H}^{d-1}\lfloor_{\{x_j = t_m\} \cap \Omega}
	\end{equation}
	so the total variation is:
	\begin{equation}
		|D\hat{f}_m|(\Omega) =
		|c_m^R - c_m^L| \cdot
		\mathcal{H}^{d-1}\!\left(\{x_j = t_m\} \cap \Omega\right)
	\end{equation}
	The geometric term satisfies
	$\mathcal{H}^{d-1}(\{x_j = t_m\} \cap \Omega) \leq \operatorname{diam}(\Omega)^{d-1}$,
	a constant depending only on $\Omega$. For the response term, by Cauchy--Schwarz
	applied to the leaf means:
	\begin{equation}
		|c_m^R - c_m^L|
		\leq \frac{1}{|I_R|}\sum_{i \in I_R} |u_m(x_i)|
		+ \frac{1}{|I_L|}\sum_{i \in I_L} |u_m(x_i)|
		\leq 2\norm{u_m}_{L^2}
	\end{equation}
	Therefore:
	\begin{equation}
		|D\hat{f}_m|(\Omega)
		\leq 2\operatorname{diam}(\Omega)^{d-1} \cdot \norm{u_m}_{L^2}
		\leq 2\operatorname{diam}(\Omega)^{d-1} \cdot C\rho^m
	\end{equation}
	using the geometric decay \eqref{eq:geometric_decay}.
	
	\subsubsection*{Step 2: Ensemble BV Bound via Geometric Series}
	
	By subadditivity of total variation:
	\begin{equation}
		|D\hat{F}_m|(\Omega)
		\leq \sum_{j=1}^m |D\hat{f}_j|(\Omega)
		\leq 2\operatorname{diam}(\Omega)^{d-1} \cdot C \sum_{j=1}^m \rho^j
		\leq \frac{2\operatorname{diam}(\Omega)^{d-1} \cdot C\rho}{1-\rho}
		< \infty
	\end{equation}
	The geometric series $\sum_{j=1}^\infty \rho^j = \rho/(1-\rho)$ converges because
	$\rho < 1$, which follows from the eigenvalue condition $\lambda_k > 0$ for all $k$.
	Taking expectations:
	\begin{equation}
		\mathbb{E}[|D\hat{F}_m|(\Omega)]
		\leq \frac{2\operatorname{diam}(\Omega)^{d-1} \cdot \mathbb{E}[C]\rho}{1-\rho}
		< \infty
	\end{equation}
	uniformly in $m$, where $\mathbb{E}[C] = \mathbb{E}[\norm{Y}_{L^2}] < \infty$
	by assumption.
	
	\subsubsection*{Step 3: Apply Lemma~\ref{lem:bv_consistency}}
	
	$L^1$ consistency is assumed. Uniform BV boundedness follows from Step~2.
	Lemma~\ref{lem:bv_consistency} gives $D\hat{F}_m \xrightarrow{w*} Df$ in $\cM$.
	\qed
	
	\begin{remark}
		The key step in both cases is the geometric decay $\norm{u_m}_{L^2} \leq C\rho^m$,
		which follows from the eigenvalue condition $0 < \lambda_k \leq 1$ via Proposition~1
		and Theorem~1 of \citet{buhlmann2003boosting}. This decay is much faster than any
		polynomial rate and implies that the geometric series $\sum_j \rho^j$ converges
		without any condition on step sizes. The eigenvalue condition is the operative
		regularity assumption for both cases of the theorem, replacing the step size
		conditions common in the boosting literature.
	\end{remark}
	
	\begin{remark}
		When $\lambda_k \in \{0,1\}$ for all $k$ --- as occurs for linear projection
		operators --- Corollary~1 of \citet{buhlmann2003boosting} shows that
		$\mathcal{B}_m \equiv S$ for all $m$, meaning boosting does not improve on the
		base learner. In this degenerate case $\rho = \max_k(1-\lambda_k)$ takes the value
		1 for any $k$ with $\lambda_k = 0$, and the geometric series does not converge.
		The strict positivity condition $\lambda_k > 0$ excludes this case and corresponds
		to the requirement that the base learner has positive but incomplete explanatory
		power --- the standard ``weak learner'' assumption in the boosting literature.
	\end{remark}
	
	% ============================================================
	\section{Proof of Theorem~\ref{thm:forests} (Mondrian Forests)}
	\label{app:forests}
	% ============================================================
	
	\subsection{Setup}
	
	Let $\fn = \frac{1}{B}\sum_{b=1}^B T_b$ where $T_b$ are individual Mondrian
	trees. Each $T_b$ is piecewise constant on a random partition $\mathcal{P}_b$
	of $\Omega = [0,1]^d$ generated by the Mondrian process $MP(\lambda_n, [0,1]^d)$
	independently of the data, with budget $\lambda_n \asymp n^{1/(d+2)}$
	\citep{mourtada2017mondrian}. The partition $\mathcal{P}_b$ consists of
	$K_{\lambda_n} + 1$ hyperrectangular cells, where $K_{\lambda_n}$ is the
	number of splits \citep{roy2009mondrian}.
	
	\subsection{Step 1: Each Tree is in BV}
	
	Each $T_b$ is piecewise constant on a finite partition and therefore in
	$\BV$ by Lemma~\ref{lem:bv_existence}. The forest average $\fn =
	\frac{1}{B}\sum_b T_b$ is a finite convex combination of $BV$ functions,
	hence $\fn \in \BV$ with:
	\begin{equation}
		|D\fn|(\Omega) \leq \frac{1}{B} \sum_{b=1}^B |DT_b|(\Omega)
	\end{equation}
	by convexity of the total variation functional
	\citep[Proposition~3.6]{ambrosio2000bv}.
	
	\subsection{Step 2: Uniform BV Boundedness via the Second Moment Argument}
	
	\textbf{Claim:} $\sup_n \E[|D\fn|(\Omega)] < \infty$.
	
	The triangle inequality applied to the forest average is too crude because
	individual tree BV norms grow with $n$. Instead we work with the second
	moment and apply Cauchy--Schwarz:
	\begin{equation}
		\E\!\left[|D\fn|(\Omega)\right]
		\leq \sqrt{\E\!\left[|D\fn|(\Omega)^2\right]}
	\end{equation}
	It therefore suffices to show $\E[|D\fn|(\Omega)^2] < \infty$ uniformly
	in $n$.
	
	\subsubsection*{Substep 2a: Expand the Second Moment}
	
	Writing $D\fn = \frac{1}{B}\sum_b DT_b$ and expanding:
	\begin{equation}
		\E\!\left[|D\fn|(\Omega)^2\right]
		= \frac{1}{B^2} \E\!\left[\left|\sum_b DT_b\right|(\Omega)^2\right]
		= \frac{1}{B^2} \sum_{b,b'} \E\!\left[
		\int \phi \cdot d(DT_b) \cdot \int \phi \cdot d(DT_{b'})
		\right]
	\end{equation}
	for a suitable test function $\phi$ with $\norm{\phi}_\infty \leq 1$.
	The double sum splits into diagonal terms $b = b'$ and cross terms
	$b \neq b'$.
	
	\subsubsection*{Substep 2b: Cross Terms Vanish}
	
	For $b \neq b'$ the partitions $\mathcal{P}_b$ and $\mathcal{P}_{b'}$ are
	independent of each other and of the data. By independence:
	\begin{equation}
		\E\!\left[\int \phi \cdot d(DT_b) \cdot \int \phi \cdot d(DT_{b'})\right]
		= \E\!\left[\int \phi \cdot d(DT_b)\right] \cdot
		\E\!\left[\int \phi \cdot d(DT_{b'})\right]
	\end{equation}
	We claim each factor is zero. By Lemma~\ref{lem:bv_existence} and the
	explicit formula for $DT_b$:
	\begin{equation}
		\E\!\left[\int \phi \cdot d(DT_b)\right]
		= \E\!\left[\sum_{j \sim k}
		(\bar{y}_k^b - \bar{y}_j^b)
		\int_{\partial R_j^b \cap \partial R_k^b} \phi \cdot \nu_{jk}
		\, d\mathcal{H}^{d-1}\right]
	\end{equation}
	Since $\mathcal{P}_b$ is independent of the data, we condition on the
	partition and take the expectation over the data first:
	\begin{equation}
		= \mathbb{E}_{\mathcal{P}_b}\!\left[\sum_{j \sim k}
		\mathbb{E}_{\mathrm{data}}\!\left[\bar{y}_k^b - \bar{y}_j^b\right]
		\int_{\partial R_j^b \cap \partial R_k^b} \phi \cdot \nu_{jk}
		\, d\mathcal{H}^{d-1}\right]
	\end{equation}
	The leaf means satisfy $\mathbb{E}_{\mathrm{data}}[\bar{y}_j^b] \to f(\bar{x}_j^b)$
	as $n \to \infty$, where $\bar{x}_j^b$ is the center of leaf $R_j^b$. For
	adjacent leaves $j$ and $k$, by the Lipschitz continuity of $f$:
	\begin{equation}
		\mathbb{E}_{\mathrm{data}}\!\left[\bar{y}_k^b - \bar{y}_j^b\right]
		\approx f(\bar{x}_k^b) - f(\bar{x}_j^b) = O(D_{\lambda_n})
	\end{equation}
	where $D_{\lambda_n}$ is the cell diameter. As $\lambda_n \to \infty$ the
	cell diameter $D_{\lambda_n} \to 0$ and each term vanishes. More precisely,
	since $f \in \Wonep$ with $p > d$, the Sobolev embedding
	$W^{1,p} \hookrightarrow C^{0,\alpha}$ for $\alpha = 1 - d/p$ gives:
	\begin{equation}
		\left|\mathbb{E}_{\mathrm{data}}\!\left[\bar{y}_k^b - \bar{y}_j^b\right]\right|
		\leq C \cdot \operatorname{diam}(R_j^b \cup R_k^b)^\alpha
		\leq C \cdot D_{\lambda_n}^\alpha \to 0
	\end{equation}
	Therefore $\mathbb{E}[\int \phi \cdot d(DT_b)] \to 0$ and the cross terms
	vanish asymptotically. For $n$ large enough the cross terms are negligible
	and the second moment is dominated by the diagonal terms.
	
	\subsubsection*{Substep 2c: Bound the Diagonal Terms}
	
	For the diagonal terms $b = b'$:
	\begin{equation}
		\E\!\left[|DT_b|(\Omega)^2\right]
		= \E\!\left[\left(\sum_{j \sim k} |\bar{y}_k^b - \bar{y}_j^b|
		\cdot \mathcal{H}^{d-1}(\partial R_j^b \cap \partial R_k^b)\right)^2\right]
	\end{equation}
	Expanding the square and separating geometric and response terms using the
	independence of $\mathcal{P}_b$ from the data:
	\begin{equation}
		\E\!\left[|DT_b|(\Omega)^2\right]
		\leq \mathbb{E}_{\mathcal{P}_b}\!\left[K_{\lambda_n}^2\right]
		\cdot \sup_{j \sim k} \mathbb{E}_{\mathrm{data}}\!\left[
		|\bar{y}_k^b - \bar{y}_j^b|^2\right]
	\end{equation}
	where $K_{\lambda_n}$ is the number of adjacent pairs, which equals the
	number of splits.
	
	\paragraph{Geometric factor.}
	The Mondrian process generates cuts according to a recursive tree structure.
	By the conditional independence of subtrees \citep{roy2009mondrian}, the
	number of splits satisfies the distributional recursion:
	\begin{equation}
		K_{\lambda_n} = 1 + K_{\lambda_n, L} + K_{\lambda_n, R}
	\end{equation}
	where $K_{\lambda_n, L}$ and $K_{\lambda_n, R}$ are the split counts in the
	left and right subtrees, conditionally independent given the first cut.
	From \citet{mourtada2017mondrian}:
	\begin{equation}
		\mathbb{E}[K_{\lambda_n}] \leq (e(\lambda_n + 1))^d = O(n^{d/(d+2)})
	\end{equation}
	Taking the second moment via the recursion and using conditional independence:
	\begin{equation}
		\mathbb{E}[K_{\lambda_n}^2]
		= \mathbb{E}[K_{\lambda_n}] + \mathbb{E}[K_{\lambda_n}]^2
		= O(n^{d/(d+2)}) + O(n^{2d/(d+2)})
		= O(n^{2d/(d+2)})
	\end{equation}
	where the first equality uses the Poisson-like variance structure of the
	Mondrian process implied by the conditional independence recursion.
	
	\paragraph{Response factor.}
	For adjacent leaves $j$ and $k$, by Cauchy--Schwarz on the sample means:
	\begin{equation}
		\mathbb{E}_{\mathrm{data}}\!\left[|\bar{y}_k^b - \bar{y}_j^b|^2\right]
		\leq 4\,\mathbb{E}_{\mathrm{data}}\!\left[\|r^{(b)}\|_{L^2}^2\right]
	\end{equation}
	where $r^{(b)}$ denotes the residual within the relevant cells. Using the
	cell diameter bound of \citet{roy2009mondrian}:
	\begin{equation}
		\mathbb{E}[D_{\lambda_n}(x)^2] \leq \frac{4d}{\lambda_n^2}
		= O(n^{-2/(d+2)})
	\end{equation}
	and the Hölder continuity of $f \in W^{1,p}$ with $p > d$ giving
	$\alpha = 1 - d/p$:
	\begin{equation}
		\mathbb{E}_{\mathrm{data}}\!\left[|\bar{y}_k^b - \bar{y}_j^b|^2\right]
		\leq C^2 \left(\mathbb{E}[D_{\lambda_n}(x)^2]\right)^\alpha
		\leq C^2 \left(\frac{4d}{\lambda_n^2}\right)^\alpha
		= O(n^{-2\alpha/(d+2)})
	\end{equation}
	
	\paragraph{Combining.}
	\begin{equation}
		\E\!\left[|DT_b|(\Omega)^2\right]
		\leq O(n^{2d/(d+2)}) \cdot O(n^{-2\alpha/(d+2)})
		= O(n^{2(d-\alpha)/(d+2)})
	\end{equation}
	Therefore:
	\begin{equation}
		\E\!\left[|D\fn|(\Omega)^2\right]
		= \frac{1}{B^2} \sum_b \E\!\left[|DT_b|(\Omega)^2\right]
		= \frac{1}{B} O(n^{2(d-\alpha)/(d+2)})
	\end{equation}
	For this to be bounded uniformly in $n$ we require:
	\begin{equation}
		B \gg n^{2(d-\alpha)/(d+2)}
	\end{equation}
	This is a condition on the number of trees relative to the dimension and
	smoothness of $f$. Under this condition:
	\begin{equation}
		\E\!\left[|D\fn|(\Omega)\right]
		\leq \sqrt{\E\!\left[|D\fn|(\Omega)^2\right]}
		= \sqrt{\frac{1}{B} O(n^{2(d-\alpha)/(d+2)})}
		< \infty
	\end{equation}
	uniformly in $n$, completing the uniform BV boundedness claim.
	
	\begin{remark}
		The condition $B \gg n^{2(d-\alpha)/(d+2)}$ requires the number of trees
		to grow with $n$ at a rate depending on dimension $d$ and Hölder exponent
		$\alpha = 1 - d/p$. For $\alpha$ close to $d$ --- achieved when $p \gg d$,
		meaning $f$ is highly smooth --- the condition is mild and a slowly growing
		number of trees suffices. For $\alpha$ small --- meaning $f$ has limited
		smoothness or $d$ is large --- more trees are required. In practice random
		forests with $B = 500$ or $B = 1000$ trees satisfy this condition for
		moderate $d$ and $n$.
	\end{remark}
	
	\subsection{Step 3: Apply Lemma~\ref{lem:bv_consistency}}
	
	The $L^1$ consistency condition follows from the $L^2$ consistency result of
	\citet{mourtada2017mondrian} under the budget $\lambda_n \asymp n^{1/(d+2)}$
	by the Cauchy--Schwarz inequality on the bounded domain $\Omega$:
	\begin{equation}
		\norm{\fn - f}_{L^1}
		\leq |\Omega|^{1/2} \norm{\fn - f}_{L^2}
		\xrightarrowp 0
	\end{equation}
	Uniform BV boundedness follows from Step~2 under the condition
	$B \gg n^{2(d-\alpha)/(d+2)}$. Lemma~\ref{lem:bv_consistency} then gives
	$D\fn \xrightarrow{w*} Df$ in $\cM$. \qed
	
	\begin{remark}
		Theorem~\ref{thm:forests} is stated for Mondrian forests where the
		partition is generated independently of the data. The independence is
		used critically in the decomposition of Substep~2a and in the cross
		term cancellation of Substep~2b. For Breiman forests with data-dependent
		splits the partition geometry and response variation are coupled ---
		splits are chosen to maximize response differences across boundaries,
		which is precisely the worst case for the BV bound, and the cross term
		cancellation argument fails because the jump directions are not
		independent across trees. Extending the result to Breiman forests
		requires additional control over this coupling and is left as future
		work.
	\end{remark}
	
	% ============================================================
	\bibliographystyle{apalike}  % or plainnat, or whatever style you prefer
	\bibliography{references}
	
\end{document}
